\chapter{Pharmacology}

\section{Autonomic \& Nervous System Pharmacology}
\medterm{Antiemetic}-- A medication that prevents or treats nausea and vomiting. Classes target different pathways and receptors: 5-HT3 receptor antagonists (ondansetron, granisetron) — block serotonin receptors in the chemoreceptor trigger zone and GI tract, first-line for chemotherapy-induced nausea and vomiting; NK1 receptor antagonists (aprepitant, fosaprepitant) — block substance P/neurokinin-1 receptors in the vomiting center, used for highly emetogenic chemotherapy; dopamine D2 antagonists (metoclopramide, prochlorperazine, droperidol) — block dopamine receptors in the CTZ, used for migraine and postoperative nausea; antihistamines (dimenhydrinate, meclizine) — block H1 receptors, effective for motion sickness; anticholinergics (scopolamine) — block muscarinic receptors, used for motion sickness; and corticosteroids (dexamethasone) — mechanism unclear, used for chemotherapy-induced nausea and postoperative nausea.

\medterm{Antiemetics}-- Medications that prevent or treat nausea and vomiting. Classes and mechanisms: 5-HT$_3$ antagonists (ondansetron — block serotonin receptors in CTZ and gut, effective for chemotherapy and post-operative nausea), dopamine antagonists (metoclopramide, prochlorperazine — block D$_2$ receptors in CTZ, prokinetic effect), antihistamines (cyclizine, promethazine — block H$_1$ receptors, effective for motion sickness), anticholinergics (hyoscine — motion sickness), neurokinin-1 antagonists (aprepitant — chemotherapy), and cannabinoids (nabilone). The cause of vomiting dictates drug choice.

\medterm{Antihistamines}-- Medications that block histamine receptors, primarily H1 and H2 receptors, used to treat allergic conditions and other histamine-mediated disorders. H1 antihistamines are classified as first-generation (sedating, with significant anticholinergic effects) and second-generation (non-sedating, with minimal anticholinergic activity). First-generation H1 antihistamines include diphenhydramine, chlorpheniramine, hydroxyzine, and doxylamine; they cross the blood-brain barrier, causing sedation, cognitive impairment, and anticholinergic side effects including dry mouth, urinary retention, and constipation; second-generation H1 antihistamines include loratadine, cetirizine, and fexofenadine, which are preferred for long-term use due to their favourable safety profile.

\medterm{Antiparkinson}-- A class of medications used to manage Parkinson disease symptoms by restoring dopaminergic activity or modulating other neurotransmitter systems. Dopaminergic agents: levodopa (combined with carbidopa to prevent peripheral decarboxylation) — the most effective symptomatic treatment, converted to dopamine in the brain; dopamine agonists (pramipexole, ropinirole, rotigotine patch) — directly stimulate dopamine receptors; MAO-B inhibitors (selegiline, rasagiline) — inhibit dopamine breakdown; COMT inhibitors (entacapone, tolcapone) — inhibit peripheral levodopa metabolism, extending its half-life. Anticholinergic agents (benztropine, trihexyphenidyl) — reduce tremor by blocking muscarinic receptors. Amantadine — increases dopamine release and reduces dyskinesias.

\medterm{Antispasmodic}-- A medication that reduces smooth muscle spasm, primarily in the gastrointestinal tract. Anticholinergic agents: hyoscyamine, dicyclomine, atropine — block muscarinic M3 receptors on GI smooth muscle, reducing contractions and peristalsis, used for irritable bowel syndrome and functional dyspepsia. Direct smooth muscle relaxants: mebeverine, papaverine, peppermint oil — act directly on smooth muscle cells without anticholinergic effects, providing relief from GI spasms. Urinary antispasmodics: oxybutynin, tolterodine, solifenacin — block M3 receptors on detrusor muscle in the bladder wall, reducing overactive bladder symptoms.

\medterm{Antithyroid}-- A class of medications that reduce thyroid hormone synthesis, used to treat hyperthyroidism (Graves disease, toxic nodular goiter). Thioamide derivatives: methimazole (propylthiouracil as second-line) — inhibit thyroid peroxidase, blocking iodination of tyrosine residues and coupling of iodotyrosines, reducing production of T3 and T4. PTU also inhibits peripheral conversion of T4 to T3. Methimazole is first-line due to better tolerability and once-daily dosing. Adverse effects: agranulocytosis (0.3-0.5\%, the most serious, presenting with fever and sore throat), hepatotoxicity (rare with methimazole, more common with PTU), rash, arthralgia, and cholestatic jaundice. Treatment monitoring: thyroid function tests every 4-6 weeks during dose titration. Beta-blockers (propranolol, atenolol) are used adjunctively to control adrenergic symptoms (tachycardia, tremor) until euthyroid.

\medterm{Anxiolytic}-- A medication that reduces anxiety without producing excessive sedation. First-line for generalized anxiety disorder: SSRIs (escitalopram, paroxetine, sertraline) and SNRIs (venlafaxine, duloxetine) — increase serotonin/norepinephrine activity, require 2-4 weeks for therapeutic effect, first-line for chronic anxiety disorders. Buspirone — a 5-HT1A partial agonist, effective for GAD, no sedation, no dependence, requires 2-4 weeks for effect. Benzodiazepines (alprazolam, lorazepam, diazepam, clonazepam) — enhance GABA-A receptor activity, rapid onset (minutes to hours), effective for acute anxiety and panic attacks, limited to short-term use (2-4 weeks) due to tolerance, dependence, and withdrawal risks. Beta-blockers (propranolol) — reduce somatic symptoms of performance anxiety (tachycardia, tremor) through peripheral beta-blockade.

\medterm{Atypical Antipsychotics}-- A class of modern psychiatric medications used to treat schizophrenia and bipolar disorder. They work by balancing dopamine and serotonin levels in the brain, which helps reduce symptoms like hallucinations and delusions with fewer movement-related side effects than older antipsychotics.

\medterm{Cardioselective Beta-Blockers}-- Cardioselective beta-blockers selectively block beta-1 adrenergic receptors, which are predominant in the heart, while having minimal effects on beta-2 receptors in the bronchial smooth muscle and peripheral vasculature at therapeutic doses. Common cardioselective beta-blockers include metoprolol, atenolol, bisoprolol, and nebivolol. At low to moderate doses, these agents produce cardiac effects including decreased heart rate, reduced myocardial contractility, and decreased cardiac output without significant bronchospasm, making them preferred in patients with asthma or COPD, though cardioselectivity is dose-dependent and may be lost at higher doses.

\medterm{Cholinesterase Inhibitor}-- A class of medications that inhibit acetylcholinesterase, the enzyme that breaks down acetylcholine in the synaptic cleft, thereby increasing cholinergic neurotransmission in the central nervous system and peripheral tissues. Centrally acting cholinesterase inhibitors: donepezil, rivastigmine, and galantamine — used for symptomatic treatment of Alzheimer disease and other dementias by enhancing central cholinergic activity, providing modest cognitive and functional benefit (typically 6-12 months of stabilization before decline resumes). Rivastigmine is also approved for Parkinson disease dementia. Peripherally acting: neostigmine, pyridostigmine — used for myasthenia gravis to increase neuromuscular junction acetylcholine levels; and physostigmine — used as an antidote for anticholinergic poisoning. Adverse effects: GI symptoms (nausea, vomiting, diarrhea, most common), bradycardia, increased gastric acid secretion, muscle cramps, and insomnia.

\medterm{Decongestants}-- Medications that reduce nasal congestion by constricting blood vessels in the nasal mucosa, decreasing swelling and mucus production. Systemic decongestants: pseudoephedrine and phenylephrine (oral), which stimulate alpha-adrenergic receptors causing vasoconstriction. Topical decongestants: oxymetazoline, xylometazoline, naphazoline, phenylephrine (nasal sprays)--more potent and rapid-acting but risk of rhinitis medicamentosa (rebound congestion) with use beyond 3--5 days. Indications: nasal congestion from allergic rhinitis, common cold, sinusitis. Contraindications: uncontrolled hypertension, severe coronary artery disease, narrow-angle glaucoma, and MAOI use.

\medterm{Dopamine Agonist}-- A class of medications that stimulate dopamine receptors (D1-D5), mimicking the effects of endogenous dopamine. Used primarily in Parkinson disease: pramipexole and ropinirole (non-ergot dopamine agonists) are effective for early and advanced PD, delaying levodopa initiation and reducing motor fluctuations; rotigotine is available as a transdermal patch providing continuous delivery. Apomorphine is a potent, rapid-acting injectable dopamine agonist used for acute rescue in refractory motor fluctuations and off episodes. Adverse effects: impulse control disorders (pathological gambling, hypersexuality, compulsive shopping, binge eating — unique to dopamine agonists, occurring in 15-25\%), nausea, orthostatic hypotension, hallucinations, and daytime somnolence (sudden sleep attacks). Ergot-derived dopamine agonists (bromocriptine, cabergoline) are rarely used due to risk of pulmonary and retroperitoneal fibrosis.

\medterm{Drug Antagonism}-- An interaction where one drug blocks or decreases the effectiveness of another, such as when an antidote blocks a poison.

\medterm{Epinephrin (Adrenaline)}-- Alternate spelling of epinephrine (adrenaline). See Epinephrine (Adrenaline).

\medterm{Epinephrine (Adrenaline)}-- A catecholamine hormone and neurotransmitter produced by the adrenal medulla and sympathetic nerve terminals, acting on alpha-1, alpha-2, beta-1, beta-2, and beta-3 adrenergic receptors. Effects: alpha-1: vasoconstriction (increases blood pressure), mydriasis; beta-1: increased heart rate and contractility; beta-2: bronchodilation, vasodilation in skeletal muscle, glycogenolysis. Clinical uses: first-line treatment for anaphylaxis (0.3--0.5 mg IM, repeat every 5--15 minutes), cardiac arrest (1 mg IV every 3--5 minutes during ACLS), severe asthma exacerbation (subcutaneous or IV), septic shock (IV infusion), and local vasoconstriction (combined with local anesthetics to prolong duration and reduce bleeding). Side effects: hypertension, tachycardia, arrhythmias, tremor, anxiety, hyperglycemia. Auto-injectors (EpiPen) are prescribed for patients with severe allergies.

\medterm{Ganglionic Blocker}-- A class of drugs that block nicotinic acetylcholine receptors at autonomic ganglia (both sympathetic and parasympathetic), inhibiting neurotransmission in both divisions of the autonomic nervous system. Due to their non-selective blockade of all autonomic ganglia, they produce widespread and unpredictable effects including severe orthostatic hypotension, blurred vision (cycloplegia), dry mouth, constipation, urinary retention, and impotence. Prototype: hexamethonium. Clinical use: historically used for hypertensive emergencies before the development of safer agents. Currently used rarely, primarily for inducing controlled hypotension during surgery (trimethaphan). Ganglionic blockers are of more pharmacological than clinical importance, serving as tools for understanding autonomic physiology.

\medterm{Inverse Agonist}-- A type of drug that binds to a cell receptor and actively produces the opposite biological effect of an agonist, effectively turning down the receptor's baseline activity.

\medterm{MAOIs}-- Monoamine oxidase inhibitors are a class of antidepressant drugs that inhibit the enzyme monoamine oxidase, which breaks down serotonin, norepinephrine, and dopamine in the brain. By inhibiting this enzyme, MAOIs increase the levels of these neurotransmitters in the synaptic cleft. There are two forms of MAO: MAO-A preferentially metabolizes serotonin and norepinephrine, while MAO-B preferentially metabolizes dopamine and tyramine. Non-selective MAOIs including phenelzine and tranylcypromine inhibit both MAO-A and MAO-B, while selective MAO-B inhibitors such as selegiline are used in Parkinson disease. MAOIs have significant potential for drug interactions and require dietary restriction of tyramine-containing foods due to the risk of hypertensive crisis.

\medterm{Neuroleptic Malignant Syndrome}-- A rare but life-threatening idiosyncratic reaction to dopamine receptor-blocking agents, typically antipsychotics, characterized by hyperthermia, severe muscle rigidity, altered mental status, and autonomic dysfunction. The pathogenesis involves sudden blockade of dopamine D2 receptors in the hypothalamus causing thermoregulatory dysfunction, in the nigrostriatal pathway causing extrapyramidal symptoms, and in the mesocortical pathway causing altered mental statstatus; treatment involves immediate discontinuation of the causative agent, supportive care, and pharmacological therapy including dantrolene (a direct-acting muscle relaxant that inhibits calcium release from the sarcoplasmic reticulum), bromocriptine (a dopamine agonist), and benzodiazepines for agitation and muscle rigidity, with careful monitoring of vital signs in an intensive care setting.

\medterm{Nonselective Beta-Blockers}-- Nonselective beta-adrenergic blocking agents block both beta-1 and beta-2 adrenergic receptors, reducing heart rate, myocardial contractility, and cardiac output through beta-1 blockade, while also blocking beta-2 receptors in the bronchial smooth muscle and peripheral vasculature. Common nonselective beta-blockers include propranolol, nadolol, timolol, and pindolol. Propranolol is the prototype agent, with lipophilic properties allowing good CNS penetration, making it useful for performance anxxiety and essential tremor, though this property also increases the risk of CNS side effects including fatigue, depression, and sleep disturbances.

\medterm{Over-the-Counter (OTC) Medications}-- Nonprescription drugs approved as safe and effective for self-management of common symptoms, available at pharmacies and retail stores without a prescription. Major categories: analgesics/antipyretics (acetaminophen—max 3000--4000 mg/day, hepatotoxicity risk with overdose; NSAIDs—ibuprofen, naproxen—GI bleeding and cardiovascular risk with chronic use), cough and cold preparations (decongestants—pseudoephedrine, phenylephrine; antihistamines—diphenhydramine, cetirizine; dextromethorphan for cough), gastrointestinal (PPIs/omeprazole, H2 blockers/famotidine for GERD; antacids; loperamide for diarrhea; psyllium, polyethylene glycol for constipation; bismuth subsalicylate for dyspepsia), allergy (oral/lora-tadine, cetirizine; intranasal corticosteroids, cromolyn), dermatologic (hydrocortisone, antifungals—clotrimazole, terbinafine), and sleep aids (diphenhydramine, doxylamine, melatonin). Risks: drug--drug interactions (NSAIDs with anticoagulants, antacids with antibiotics), unintentional overdose (especially combination products with acetaminophen), and delayed diagnosis of serious conditions. Pharmacist counseling improves safe OTC selection.

\medterm{Parasympatholytic}-- Also called anticholinergic, a drug that blocks muscarinic (cholinergic) receptors, inhibiting parasympathetic effects. Agents: atropine (bradycardia, organophosphate poisoning), ipratropium/tiotropium (COPD/asthma—bronchodilation), scopolamine (motion sickness), oxybutynin/tolterodine (overactive bladder), and glycopyrrolate. Effects: mydriasis, cycloplegia, tachycardia, dry mouth, urinary retention, constipation, and bronchodilation. Adverse effects: confusion (elderly), heat intolerance, and glaucoma precipitation. Reversed by physostigmine.

\medterm{Parasympathomimetic}-- A drug that mimics or enhances parasympathetic (cholinergic) activity. Direct agonists: muscarinic (pilocarpine—glaucoma, dry mouth) and nicotinic agonists. Indirect (acetylcholinesterase inhibitors): neostigmine (myasthenia gravis, postoperative ileus/urinary retention), physostigmine (anticholinergic poisoning), donepezil/rivastigmine (Alzheimer disease). Effects: miosis, bradycardia, bronchoconstriction, salivation, lacrimation, GI stimulation, and bladder contraction. Adverse effects: bradycardia, bronchospasm, diarrhea, and sweating; reversed by atropine.

\medterm{Placebo Effect}-- A psychological and physiological phenomenon where a patient experiences a real improvement in health or symptoms after receiving an inactive treatment (like a sugar pill) due to their expectations of healing.

\medterm{Receptor Downregulation}-- A decrease in the number or sensitivity of receptors in response to chronic receptor stimulation by agonists. When receptors are chronically activated, the cells compensate by internalizing and degrading receptors, reducing receptor density on the cell surface, and decreasing receptor responsiveness. This process contributes to drug tolerance, where increasing doses are needed to achieve the same therapeutic effect. Chronic use of beta-agonists such as albuterol for asasthma leads to reduced bronchodilator response, and tolerance to opioid analgesics develops through downregulation of mu-opioid receptors in the central nervous system.

\medterm{Receptor Upregulation}-- An increase in the number or sensitivity of receptors in response to chronic receptor blockade or decreased neurotransmitter activity. When receptors are chronically blocked by antagonists, the cells compensate by synthesizing more receptors to maintain homeostatic signaling, leading to increased sensitivity to agonists. This phenomenon contributes to drug tolerance, where increasing doses are needed to achieve the same effect, and to withdrawal syndromes when the blockidrug is discontinued; well-recognised examples include the upregulation of beta-adrenergic receptors following chronic beta-blocker therapy, causing rebound tachycardia and hypertension upon abrupt withdrawal, and the upregulation of histamine H2 receptors after chronic proton pump inhibitor use, leading to rebound acid hypersecretion.

\medterm{Salbutamol (Albuterol)}-- A selective beta-2 adrenergic agonist bronchodilator used for acute relief and prevention of bronchospasm in asthma and COPD. Relaxes airway smooth muscle; short-acting (SABA), onset 5-15 minutes, duration 4-6 hours via inhaler or nebulizer. Adverse effects: tremor, tachycardia, palpitations, hypokalemia, and headache. First-line rescue therapy; overuse signals poor control. Levalbuterol is the R-enantiomer. Salmeterol and formoterol are long-acting (LABAs, not for acute relief).

\medterm{Sedatives}-- A class of drugs that produce a calming or drowsy effect by depressing the central nervous system. Sedatives reduce anxiety, promote sleep, and induce relaxation. Classes: (1) Benzodiazepines—enhance GABA$_A$ receptor activity (diazepam, lorazepam, midazolam, temazepam). Used for anxiety, insomnia, seizures, muscle relaxation, procedural sedation, and alcohol withdrawal. (2) Z-drugs—non-benzodiazepine hypnotics (zolpidem, zopiclone, zaleplon)—selective for GABA$_A$ $\alpha$1 subunit, primarily for insomnia. (3) Barbiturates—enhance GABA$_A$ activity but with a narrower therapeutic index (phenobarbital, thiopental)—now rarely used except for refractory epilepsy and anesthetic induction. (4) Antihistamines—H1 antagonists with sedative properties (diphenhydramine, hydroxyzine, promethazine)—mild sedation. (5) Antipsychotics—low-dose quetiapine, olanzapine—used as sedatives off-label. (6) Melatonin receptor agonists (ramelteon, prolonged-release melatonin)—regulate circadian rhythm. (7) Orexin receptor antagonists (suvorexant, daridorexant)—promote sleep by blocking wake-promoting orexin signalling. (8) $\alpha$2-adrenergic agonists (clonidine, dexmedetomidine)—sedation without respiratory depression. Adverse effects: drowsiness, daytime sedation, cognitive impairment, impaired coordination, falls (especially elderly), tolerance, dependence, and withdrawal. Respiratory depression (especially with alcohol, opioids). Long-term use of benzodiazepines is discouraged (dependence, tolerance, cognitive decline, increased dementia risk—controversial).

\medterm{Serotonin Syndrome}-- A potentially life-threatening condition caused by excessive serotonergic activity in the central and peripheral nervous systems, typically resulting from drug interactions or overdose involving serotonergic medications. The classic triad includes neuromuscular hyperactivity with tremor, clonus, rigidity, and hyperreflexia; autonomic dysfunction with hyperthermia, diaphoresis, tachycardia, and hypertension; and altered mental status with agitation, confusion, and delirium. Severe cases may progress to hyperpyrexia, rhabdomyolysis, multi-organ failure, and death. Treatment involves immediate discontinuation of all serotonergic agents, supportive care with intravenous fluids and cooling, and administration of cyproheptadine, a serotonin antagonist, for moderate to severe cases.

\medterm{Skeletal Muscle Relaxant}-- A medication that reduces muscle spasm, spasticity, and associated pain. Centrally acting muscle relaxants: baclofen — GABA-B agonist that reduces spinal reflex activity, first-line for spasticity in cerebral palsy, multiple sclerosis, and spinal cord injury; tizanidine — alpha-2 adrenergic agonist that inhibits spinal motor neuron activity; diazepam — enhances GABA-A activity, effective for acute muscle spasm; cyclobenzaprine — structurally related to tricyclic antidepressants, reduces brainstem descending noradrenergic activity; methocarbamol and carisoprodol — centrally acting, mechanism less defined. Peripherally acting: dantrolene — inhibits calcium release from the sarcoplasmic reticulum by blocking ryanodine receptors, used for malignant hyperthermia and chronic spasticity. Botulinum toxin — blocks acetylcholine release at the neuromuscular junction, used for focal spasticity and dystonia.

\medterm{SNRIs}-- Serotonin-norepinephrine reuptake inhibitors are a class of antidepressant drugs that inhibit the reuptake of both serotonin and norepinephrine from the synaptic cleft, increasing both serotonergic and noradrenergic neurotransmission. Common SNRIs include venlafaxine, duloxetine, desvenlafaxine, and levomilnacipran. They are used for major depressive disorder, generalized anxiety disorder, panic disorder, social anxiety disorder, and chronic pain conditions. Duloxetine is also approved for diabeabetic peripheral neuropathy, fibromyalgia, and chronic musculoskeletal pain. Common adverse effects include nausea, dry mouth, insomnia, sweating, and sexual dysfunction; discontinuation syndrome with dizziness and paraesthesias may occur with abrupt cessation, requiring gradual tapering.

\medterm{Tachyphylaxis}-- A rapid, acute decrease in response to a drug after repeated administration over a short period, distinct from tolerance which develops more gradually. Tachyphylaxis results from rapid depletion of neurotransmitter stores, receptor desensitization, or receptor downregulation. A classic example is the rapid development of tachyphylaxis to ephedrine, where repeated doses lead to depletion of norepinephrine stores in nerve terminals, reducing the sympathomimetic effect. Nicotine alalso demonstrates tachyphylaxis, and this phenomenon contributes to the need for dose escalation during the first few days of therapy. Tachyphylaxis is distinct from tolerance, which develops more slowly, and from desensitisation, which results from reduced receptor responsiveness.

\medterm{TCAs}-- Tricyclic antidepressants are a class of drugs that inhibit the reuptake of serotonin and norepinephrine, similar to SNRIs, but also block histamine H1 receptors, muscarinic acetylcholine receptors, and alpha-1 adrenergic receptors, producing numerous side effects. Common TCAs include amitriptyline, nortriptyline, imipramine, and clomipramine. They are effective antidepressants but are now used less frequently for depression due to their side effect profile and toxicity in overdose. TCAs are used for a range of other conditions including neuropathic pain (amitriptyline, nortriptyline), migraine prophylaxis (amitriptyline), chronic pain syndromes, enuresis, and insomnia (doxepin, trimipramine). Their anticholinergic effects cause dry mouth, blurred vision, constipation, urinary retention, and cognitive impairment in elderly patients; cardiotoxicity in overdose manifests as QT prolongation, conduction abnormalities, and potentially fatal arrhythmias, limiting their use in patients with cardiac disease or at risk of suicide.

\medterm{Tocolytics}-- Medications used to suppress preterm labor (contractions before 37 weeks gestation), to delay delivery for corticosteroid administration (fetal lung maturity) and in utero transport to a tertiary care center. First-line: calcium channel blockers (nifedipine) — inhibit smooth muscle contraction by blocking voltage-gated calcium channels, oral administration, fewer side effects than other tocolytics. Beta-2 adrenergic agonists (terbutaline, ritodrine) — activate beta-2 receptors on uterine smooth muscle, causing relaxation via increased intracellular cAMP, side effects include maternal tachycardia, pulmonary edema, and hypokalemia. NSAIDs (indomethacin) — inhibit prostaglandin synthesis by COX inhibition, effective for short-term tocolysis (<48 hours), side effects include oligohydramnios and premature closure of the fetal ductus arteriosus with prolonged use. Magnesium sulfate — blocks calcium entry into smooth muscle cells, used for neuroprotection in preterm infants (<32 weeks) and seizure prophylaxis in preeclampsia, but has marginal tocolytic efficacy.

\medterm{Typical Antipsychotics}-- Typical antipsychotics, also known as first-generation antipsychotics, primarily block dopamine D2 receptors in the mesolimbic pathway, reducing positive symptoms of schizophrenia including hallucinations, delusions, and disorganized thinking. They are classified by potency: high-potency agents including haloperidol and fluphenazine require lower doses and have more extrapyramidal side effects; low-potency agents including chlorpromazine and thioridazine require higher doses and have more anticholinergic side effects including sedation, dry mouth, constipation, and postural hypotension; all typical antipsychotics carry a risk of tardive dyskinesia, a potentially irreversible movement disorder characterized by involuntary, repetitive movements of the face, mouth, and extremities, with higher potency agents posing a greater risk.

\medterm{Uterotonic}-- A medication that stimulates uterine contraction, used for induction/augmentation of labor, prevention and treatment of postpartum hemorrhage, and medical termination of pregnancy. Oxytocin — a synthetic analog of the endogenous hormone, stimulates oxytocin receptors on uterine smooth muscle, producing rhythmic contractions, administered IV for labor induction/augmentation; also used IV or IM for active management of the third stage of labor to prevent postpartum hemorrhage. Ergot alkaloids (ergometrine, methylergonovine) — cause sustained uterine tetanic contraction by acting on alpha-adrenergic and serotonin receptors, used IM/IV for postpartum hemorrhage, contraindicated in hypertension and preeclampsia. Prostaglandins: carboprost (PGF2-alpha) — IM for refractory postpartum hemorrhage; misoprostol (PGE1 analog) — oral, sublingual, rectal, or vaginal for cervical ripening, labor induction, and postpartum hemorrhage; dinoprostone (PGE2) — vaginal insert for cervical ripening.

\section{Cardiovascular \& Renal Pharmacology}
\medterm{5-Alpha Reductase Inhibitor}-- A class of medications that inhibit the enzyme 5-alpha-reductase, which converts testosterone to dihydrotestosterone (DHT), the more potent androgen responsible for prostate growth and male pattern baldness. Two isoforms: Type I (predominant in skin and scalp) and Type II (predominant in the prostate). Finasteride inhibits primarily Type II (blocks 70\% of DHT); dutasteride inhibits both Type I and Type II (blocks >90\% of DHT). Indications: benign prostatic hyperplasia — reduces prostate volume, improves urinary flow, reduces risk of acute urinary retention and need for surgery (requires 6-12 months for maximal effect); male androgenetic alopecia — slows hair loss and promotes regrowth. Adverse effects: decreased libido, erectile dysfunction (2-5\%), ejaculatory dysfunction, and gynecomastia. Serum PSA levels decrease by approximately 50\% — adjusted PSA values are needed for prostate cancer screening.

\medterm{Absorption}-- The process by which a drug enters the bloodstream from its site of administration, such as the digestive tract, skin, or muscle.

\medterm{ACE Inhibitors}-- Angiotensin-converting enzyme inhibitors are a class of medications that block the conversion of angiotensin I to angiotensin II by inhibiting the angiotensin-converting enzyme (ACE, also known as kininase II). This results in decreased aldosterone secretion, reduced sodium and water retention, decreased vasoconstriction, and reduced sympathetic nervous system activity. ACE inhibitors also inhibit the breakdown of bradykinin, a vasodilator, which contributes to their antihypertensive effect but also contributes to the characteristic adverse effect of dry cough, which occurs in approximately 10-20\% of patients, and the less common but more serious angioedema.

\medterm{Acetaminophen Mechanism}-- Acetaminophen, also known as paracetamol, is an analgesic and antipyretic drug with weak anti-inflammatory activity, differing from NSAIDs in its mechanism and side effect profile. Its precise mechanism of action remains incompletely understood, but it is thought to inhibit a COX variant in the central nervous system, possibly COX-3, and may also activate descending serotonergic pain inhibition pathways and interact with the endocannabinoid system through its metabolite AM404. Acetaminophen inhibits prostaglandin synthesis in the central nervous system, reducing pain perception and fever, and also activates descending serotonergic pathways that modulate pain signalling at the spinal cord level.

\medterm{Acetaminophen (Paracetamol)}-- An analgesic and antipyretic drug that inhibits central COX enzymes and may have serotonergic and endocannabinoid mechanisms. It is first-line for mild pain and fever with fewer side effects than NSAIDs. Maximum daily dose is 4 grams (reduced to 2-3 grams in liver disease or alcohol use). Acetaminophen overdose causes hepatotoxicity due to the toxic metabolite NAPQI. Antidote is N-acetylcysteine, which replenishes glutathione stores.

\medterm{Acetazolamide}-- A carbonic anhydrase inhibitor diuretic. Blocks carbonic anhydrase in the proximal tubule, reducing bicarbonate reabsorption, producing a mild diuresis and metabolic acidosis. Uses: acute mountain sickness prophylaxis/treatment, glaucoma (reduces aqueous humor production), idiopathic intracranial hypertension, metabolic alkalosis, and as an adjunct in some seizure disorders and urinary alkalinization. Given orally or IV. Side effects: paresthesias, hypokalemia, metabolic acidosis, renal stones, sulfonamide allergy cross-reactivity. Contraindicated in hepatic cirrhosis (hepatic encephalopathy risk) and severe renal impairment.

\medterm{Acinetobacter baumannii resistance}-- Acinetobacter baumannii is a gram-negative coccobacillus that has emerged as a significant cause of healthcare-associated infections, particularly in intensive care units. The organism has intrinsic and acquired resistance mechanisms, including beta-lactamases, efflux pumps, and porin mutations. Multidrug-resistant and carbapenem-resistant strains are increasingly common. Treatment options are limited and may include polymyxins, tigecycline, and minocycline. Infection prevention strategies are critical for preventing nosocomial transmission, particularly as Acinetobacter can survive on environmental surfaces for extended periods.

\medterm{Adverse Drug Reaction}-- An adverse drug reaction is any noxious or unintended response to a drug that occurs at doses used for prophylaxis, diagnosis, or therapy. ADRs are classified as type A reactions, which are dose-dependent and predictable extensions of the drug's pharmacological effects, such as bleeding with warfarin or hypoglycemia with insulin; and type B reactions, which are idiosyncratic, unpredictable, and not dose-related, such as drug allergies and Stevens-Johnson syndrome. ADRs are further classified as type C, dose- and time-dependent reactions such as cumulative toxicity; type D, delayed effects such as carcinogenicity and teratogenicity; and type E, withdrawal reactions upon discontinuation.

\medterm{Adverse Drug Reactions}-- Any harmful or unintended response to a drug at normal doses. ADRs are classified as Type A (augmented, dose-dependent, predictable: side effects, toxicity) or Type B (bizarre, dose-independent, unpredictable: allergy, idiosyncrasy). ADRs are a major cause of morbidity and mortality, accounting for 5-7\% of hospital admissions. Risk factors include polypharmacy, renal/hepatic impairment, age extremes, and genetic predisposition. Prevention involves careful prescribing and monitoring.

\medterm{Agonist-Antagonist}-- A drug that has mixed pharmacological activity, acting as an agonist at one receptor type and an antagonist at another, or as a partial agonist at a single receptor. Mixed opioid agonist-antagonists: pentazocine — kappa agonist and mu partial antagonist, producing analgesia with less respiratory depression and lower abuse potential than pure mu agonists; nalbuphine — kappa agonist and mu antagonist; butorphanol — kappa agonist and mu antagonist. Partial agonists: buprenorphine — partial mu agonist and kappa antagonist, used for opioid use disorder due to ceiling effect on respiratory depression and lower abuse liability. The clinical advantage of agonist-antagonists is the reduced risk of respiratory depression and dependence compared to full mu agonists.

\medterm{Aminoglycoside nephrotoxicity}-- Kidney damage caused by aminoglycoside antibiotics, characterized by acute tubular necrosis. Nephrotoxicity is caused by accumulation of aminoglycosides in the proximal tubular cells. The risk is increased with prolonged therapy, high doses, concurrent nephrotoxic drugs, and renal impairment. Gentamicin is considered the most nephrotoxic aminoglycoside. The nephrotoxicity is usually reversible with early detection and discontinuation. Regular monitoring of serum creatinine and drug levels is essential for early detection of nephrotoxicity, and dose adjustment based on renal function is required to minimise the risk of acute kidney injury.

\medterm{Aminoglycoside ototoxicity}-- Damage to the inner ear caused by aminoglycoside antibiotics, resulting in hearing loss and vestibular dysfunction. Ototoxicity is caused by accumulation of aminoglycosides in the inner ear, damaging hair cells. The risk is increased with prolonged therapy, high doses, concurrent ototoxic drugs, and renal impairment. Amikacin is considered the most ototoxic aminoglycoside. Symptoms include tinnitus, hearing loss, and vertigo. Hearing loss may be irreversible. Baseline and periodic audiometry should be performed at baseline and repeated regularly during treatment, particularly in patients receiving high cumulative doses or those with pre-existing hearing impairment.

\medterm{Aminoglycosides}-- A class of bactericidal antibiotics that bind irreversibly to the 30S ribosomal subunit, causing misreading of messenger RNA and inhibiting protein synthesis. They include gentamicin, tobramycin, amikacin, and streptomycin. Aminoglycosides are particularly effective against aerobic gram-negative organisms including Pseudomonas aeruginosa and have synergistic activity with beta-lactams against enterococci and staphylococci. They are poorly absorbed from the gastrointestinal tract and must be administered parenterally for systemic infections, with careful monitoring of serum concentrations to balance efficacy against nephrotoxicity and ototoxicity.

\medterm{Analgesic}-- A drug that relieves pain without producing loss of consciousness. Classes: non-opioids (acetaminophen, NSAIDs—act peripherally and on central prostaglandin synthesis), opioids (morphine, oxycodone, fentanyl—act on central and spinal opioid receptors), and adjuvants (anticonvulsants such as gabapentin/pregabalin, tricyclic antidepressants, and local anesthetics) used for neuropathic pain. Selection depends on pain type (nociceptive vs neuropathic), severity (WHO analgesic ladder), and comorbidities. Adverse effects: GI bleeding (NSAIDs), hepatotoxicity (acetaminophen), respiratory depression and dependence (opioids).

\medterm{Angiotensin Receptor Blockers}-- A class of cardiovascular drugs that selectively block the angiotensin II type 1 receptor, preventing the binding of angiotensin II and its effects including vasoconstriction, aldosterone secretion, and sympathetic activation. Unlike ACE inhibitors, ARBs do not affect bradykinin metabolism, resulting in lower rates of cough and angioedema. Common ARBs include losartan, valsartan, irbesartan, and olmesartan. ARBs have similar hemodynamic effects to ACE inhibitorsbut are generally better tolerated due to a lower incidence of cough and angioedema.

\medterm{Antacids}-- Medications that neutralise gastric acid, providing rapid symptomatic relief in dyspepsia, gastro-esophageal reflux, and peptic ulcer disease. Commonly contain aluminium hydroxide, magnesium hydroxide, calcium carbonate, or sodium bicarbonate. Aluminium-based antacids can cause constipation; magnesium-based antacids can cause diarrhea; combination preparations balance these effects. Antacids also promote gastric mucosal defence by stimulating bicarbonate and prostaglandin secretion. Drug interactions: reduce absorption of tetracyclines, iron, digoxin, and bisphosphonates (separate dosing by 2 hours). Not first-line for chronic GERD (PPIs preferred).

\medterm{Antianginal}-- A class of medications used to prevent or relieve angina pectoris (chest pain from myocardial ischemia). The three main classes of antianginal drugs target different mechanisms: nitrates (nitroglycerin, isosorbide mononitrate/dinitrate) — venodilators that reduce preload and myocardial oxygen demand; beta-blockers — reduce heart rate, contractility, and oxygen demand; and calcium channel blockers (verapamil, diltiazem, nifedipine) — vasodilators that reduce afterload and coronary spasm. Ranolazine is a newer antianginal agent that inhibits the late sodium current in cardiac myocytes, reducing intracellular calcium overload and improving diastolic function without affecting heart rate or blood pressure.

\medterm{Antiarrhythmic}-- A drug used to treat or prevent cardiac arrhythmias, classified by the Vaughan Williams system: Class I (sodium channel blockers: quinidine, lidocaine, flecainide), Class II (beta-blockers), Class III (potassium channel blockers prolonging repolarization: amiodarone, sotalol), Class IV (calcium channel blockers: verapamil, diltiazem), and Class V (other: digoxin, adenosine). Selection depends on the arrhythmia (AF, SVT, VT) and underlying heart disease. Antiarrhythmics can themselves be proarrhythmic; amiodarone requires monitoring for pulmonary, thyroid, hepatic, and ocular toxicity.

\medterm{Antibiotic cycling}-- The scheduled rotation of antimicrobial agents within a healthcare facility to reduce selective pressure for resistance. The strategy involves periodically changing the preferred antibiotics for specific infections. For example, cycling between carbapenems and fluoroquinolones for gram-negative infections. The evidence for antibiotic cycling is mixed, with some studies showing reductions in resistance and others showing no benefit. Cycling may be combined with other stewardship strategies includincluding education, surveillance, and the development of evidence-based prescribing guidelines tailored to local resistance patterns.

\medterm{Antibiotic de-escalation}-- The process of narrowing antimicrobial spectrum based on microbiological results and clinical response. De-escalation aims to reduce selective pressure for resistance while maintaining therapeutic efficacy. The process involves identifying the causative organism and susceptibility pattern, then switching from broad-spectrum to narrow-spectrum agents. De-escalation should be performed as soon as possible, typically within 48-72 hours when culture results are available. This strategy reduces the rrisk of antibiotic resistance, selective pressure on the microbiome, superinfection risk, and adverse drug effects, while also lowering overall treatment costs.

\medterm{Antibiotic mixing}-- A stewardship strategy where different patients with the same indication receive different antibiotics concurrently, with the goal of reducing resistance emergence. For example, using different agents for ventilator-associated pneumonia in different patients within the same unit. The evidence for antibiotic mixing is limited, but theoretical models suggest it may reduce resistance emergence. This strategy may be combined with other stewardshipstrategies including prospective audit and feedback, formulary restriction, and computerised decision support to achieve optimal antimicrobial use.

\medterm{Antibiotics}-- Medications that kill or inhibit the growth of bacteria, used to treat bacterial infections such as pneumonia, strep throat, and urinary tract infections. They are ineffective against viral infections like colds or influenza. Overuse and misuse have driven antibiotic resistance, a growing public health crisis that makes infections harder to treat.

\medterm{Anticonvulsants}-- Medications used to prevent or control seizures in epilepsy and other seizure disorders. Mechanisms: enhance GABA activity (benzodiazepines, barbiturates, vigabatrin, tiagabine), block voltage-gated sodium channels (phenytoin, carbamazepine, lamotrigine), block calcium channels (gabapentin, pregabalin, ethosuximide), or inhibit glutamate release (levetiracetam). Choice depends on seizure type: sodium valproate for generalised, carbamazepine/lamotrigine for focal. Side effects: sedation, dizziness, ataxia, rash, hepatotoxicity, teratogenicity (valproate). Therapeutic drug monitoring required for phenytoin and carbamazepine.

\medterm{Antidiabetic}-- A class of medications used to manage diabetes mellitus by lowering blood glucose levels through various mechanisms. Major classes include: biguanides (metformin) — reduce hepatic glucose production; sulfonylureas (glipizide, glyburide) — stimulate insulin secretion; meglitinides (repaglinide) — stimulate rapid, short-acting insulin secretion; thiazolidinediones (pioglitazone) — improve insulin sensitivity; DPP-4 inhibitors (sitagliptin) — increase incretin levels; GLP-1 receptor agonists (semaglutide, liraglutide) — stimulate insulin secretion, slow gastric emptying; SGLT2 inhibitors (empagliflozin, dapagliflozin) — increase urinary glucose excretion; insulin preparations — replace endogenous insulin; and alpha-glucosidase inhibitors (acarbose) — delay carbohydrate absorption.

\medterm{Antidiarrheal}-- A drug that reduces the frequency or volume of diarrhea. Loperamide, a peripheral mu-opioid receptor agonist, slows intestinal motility and reduces fluid loss without CNS effects; diphenoxylate/atropine combinations and bismuth subsalicylate are alternatives. Antidiarrheals are avoided in invasive/inflammatory diarrhea (dysentery, C. difficile infection) because they may prolong infection or precipitate toxic megacolon. Oral rehydration is first-line for dehydration. Octreotide is used for secretory diarrhea; cholestyramine for bile acid diarrhea.

\medterm{Antiepileptic Drugs}-- Medications used to prevent and control seizures by modulating neuronal excitability through various mechanisms. They can be classified by their primary mechanism of action. Sodium channel blockers including phenytoin, carbamazepine, and lamotrigine stabilize the inactivated state of voltage-gated sodium channels, reducing high-frequency repetitive firing. GABA-enhancing drugs including valproic acid, benzodiazepines, and vigabatrin increase inhibitory GABAergic transmissmission; glutamate receptor antagonists such as topiramate and felbamate reduce excitatory neurotransmission; and drugs with multiple mechanisms including valproic acid and levetiracetam provide broad-spectrum seizure control through combined action on GABA, sodium channels, and calcium channels.

\medterm{Antifungal}-- A drug used to treat fungal infections. Classes: polyenes (amphotericin B—broad, fungicidal, for severe/systemic mycoses), azoles (fluconazole, itraconazole, voriconazole, posaconazole—inhibit ergosterol synthesis), echinocandins (caspofungin, micafungin, anidulafungin—inhibit cell wall glucan synthesis), and allylamines (terbinafine—for dermatophytes). Griseofulvin, flucytosine, and nystatin (topical/enteral) are additional agents. Selection depends on the organism, site, and host immunity. Antifungal resistance is emerging (Candida auris).

\medterm{Antifungal prophylaxis}-- The use of antifungal agents to prevent invasive fungal infections in high-risk patients. See microbiology.

\medterm{Antigout}-- A class of medications used to treat gout by reducing uric acid levels or controlling acute inflammation. Acute gout treatment: NSAIDs (indomethacin, naproxen) — first-line for pain and inflammation; colchicine — inhibits microtubule polymerization, reducing neutrophil chemotaxis and inflammation, most effective when started within 24 hours of attack; corticosteroids (oral, intra-articular, or IV) for refractory cases. Long-term urate-lowering therapy: xanthine oxidase inhibitors (allopurinol, febuxostat) — reduce uric acid production by inhibiting xanthine oxidase; uricosuric agents (probenecid, lesinurad) — increase renal uric acid excretion by inhibiting URAT1 transporter; and recombinant uricase (pegloticase) — enzymatically degrades uric acid, reserved for severe tophaceous gout.

\medterm{Antihistamine}-- A drug that blocks histamine receptors. H1 antihistamines treat allergic rhinitis, urticaria, prpruritus, and (older sedating agents) insomnia and motion sickness. First-generation (diphenhydramine, chlorpheniramine): CNS-penetrant—sedating, anticholinergic (dry mouth, urinary retention, confusion in elderly). Second-generation (loratadine, cetirizine, fexofenadine): non-sedating, peripherally selective. H2 antihistamines (famotidine, ranitidine) reduce gastric acid (see H2 Blockers). Used in anaphylaxis as adjuncts (after epinephrine).

\medterm{Antihypertensive}-- A drug that lowers blood pressure, used for hypertension management. Major classes: ACE inhibitors and ARBs (preferred first-line, renoprotective), calcium channel blockers, thiazide diuretics, beta-blockers (with CAD/heart failure), alpha-blockers, centrally acting agents (clonidine), and direct vasodilators (hydralazine, minoxidil). Combination therapy is often needed; lifestyle modification is foundational. Goals vary by risk (typically <130/80 mmHg). Adverse effects are class-specific; ACE inhibitor cough, hyperkalemia, and angioedema require vigilance. Hypertensive urgency/emergency requires urgent IV therapy.

\medterm{Anti-inflammatory}-- An agent that reduces inflammation. Nonsteroidal anti-inflammatory drugs (NSAIDs) inhibit cyclooxygenase, reducing prostaglandin synthesis; corticosteroids have broader anti-inflammatory actions via glucocorticoid receptor-mediated gene regulation; colchicine, methotrexate, sulfasalazine, and biologic agents (TNF inhibitors, IL-6 inhibitors) are used for inflammatory and autoimmune disease. Anti-inflammatories are used for pain, fever, arthritis, inflammatory bowel disease, asthma, and allergic conditions. Long-term use carries risks: GI ulceration, renal impairment, cardiovascular events (NSAIDs); immunosuppression, hyperglycemia, osteoporosis (corticosteroids).

\medterm{Antimicrobial desensitization}-- A procedure that involves administering gradually increasing doses of an antimicrobial to induce tolerance in patients with allergy. Desensitization is used when the causative agent is the only effective treatment and the patient has a history of hypersensitivity. The procedure is performed under close medical supervision with monitoring for anaphylaxis. It is most commonly performed for beta-lactam antibiotics in patients with severe allergy. Desensitization involves starting with very low doses, typically micrograms, and increasing gradually over several hours until the full therapeutic dose is tolerated, after which desensitisation must be maintained by continuous drug exposure.

\medterm{Antimicrobial drug interactions}-- Interactions between antimicrobial agents and other medications. See microbiology.

\medterm{Antimicrobial hypersensitivity reactions}-- Hypersensitivity reactions to antimicrobial agents range from mild rash to life-threatening anaphylaxis. Beta-lactam antibiotics are the most common cause of drug allergy. Cross-reactivity between penicillins and cephalosporins is approximately 2\%. Sulfonamide antibiotics can cause severe cutaneous reactions. Fluoroquinolones can cause photosensitivity and rare anaphylaxis. Vancomycin can cause red man syndrome, which is an infusion-related reaction rather than true allergy. Proper allergy assessment is essential, including skin testing when indicated, and documentation of the reaction type and severity to guide future prescribing decisions.

\medterm{Antimicrobial resistance mechanisms}-- Mechanisms by which bacteria develop resistance to antimicrobial agents. Major mechanisms include enzymatic inactivation (beta-lactamases, aminoglycoside-modifying enzymes), target modification (PBP alterations, ribosomal methylation), efflux pumps, and decreased permeability (porin mutations). Resistance can be intrinsic, acquired through mutation, or transferred horizontally via plasmids, transposons, and integrons. Understanding resistance mechanisms is essential for developing new antibioticantibiotics and for guiding empirical therapy based on local resistance surveillance data.

\medterm{Antimicrobial Stewardship}-- See infectious\_disease.

\medterm{Antineoplastic}-- A drug used to treat cancer (chemotherapy). Categories: alkylating agents (cyclophosphamide, cisplatin), antimetabolites (methotrexate, 5-FU), antitumor antibiotics (doxorubicin), plant alkaloids (vincristine, paclitaxel), topoisomerase inhibitors (irinotecan, etoposide), hormonal agents, targeted therapies (TKIs, monoclonal antibodies), and immunotherapies (checkpoint inhibitors). Most act on rapidly dividing cells, causing myelosuppression, mucositis, alopecia, nausea, and gonadal toxicity. Regimens combine drugs with different mechanisms and toxicities to maximize efficacy and minimize resistance.

\medterm{Antioxidants}-- Substances that prevent or slow oxidative damage to cells caused by free radicals (reactive oxygen and nitrogen species). Endogenous antioxidants: glutathione, superoxide dismutase, catalase. Dietary antioxidants: vitamin C (ascorbic acid — water-soluble, scavenges aqueous-phase radicals), vitamin E (tocopherol — lipid-soluble, protects cell membranes), beta-carotene (provitamin A), selenium (cofactor for glutathione peroxidase), and flavonoids. Despite theoretical benefits, large clinical trials have not consistently demonstrated reduced cardiovascular events or cancer with antioxidant supplementation. A balanced diet rich in fruits and vegetables provides adequate antioxidants.

\medterm{Antiparasitic}-- A drug used to treat infections caused by parasites (protozoa and helminths). Antiprotozoals: metronidazole (Giardia, amebiasis, trichomoniasis), tinidazole, chloroquine and artemisinin combinations (malaria), pyrimethamine (toxoplasmosis), pentamidine, and nifurtimox (Chagas). Anthelmintics: albendazole and mebendazole (roundworm, hookworm, pinworm, tapeworm), ivermectin (strongyloidiasis, onchocerciasis, scabies), praziquantel (schistosomiasis, flukes), and diethylcarbamazine (filariasis). Many require weight-based dosing and repeated courses; adverse effects depend on the class.

\medterm{Antiparkinsonian}-- A drug used to treat Parkinson disease and drug-induced parkinsonism. Classes: levodopa/carbidopa (most effective), dopamine agonists (pramipexole, ropinirole, rotigotine), MAO-B inhibitors (selegiline, rasagiline), COMT inhibitors (entacapone), anticholinergics (benztropine, trihexyphenidyl—for drug-induced parkinsonism and tremor), and amantadine (early/adjunct, dyskinesia). Selection depends on age, severity, and predominant symptoms; levodopa remains the gold standard.

\medterm{Antipruritic}-- A medication that relieves itching (prpruritus). First-line for most causes: topical emollients and moisturizers — restore skin barrier function; topical corticosteroids — reduce inflammation in atopic dermatitis, contact dermatitis, and psoriasis; topical antihistamines (diphenhydramine, doxepin) — block histamine H1 receptors in the skin; topical calcineurin inhibitors (tacrolimus, pimecrolimus) — for atopic dermatitis without steroid side effects. Systemic antiprurities: oral antihistamines (cetirizine, loratadine, hydroxyzine) — especially for urticaria; antidepressants (doxepin, mirtazapine) — for chronic prpruritus; opioid antagonists (naltrexone) — for cholestatic prpruritus; and aprepitant — an NK1 receptor antagonist for refractory prpruritus.

\medterm{Antipsychotic}-- A drug used to treat psychotic disorders (schizophrenia, schizoaffective disorder, psychosis in bipolar/manic states and delirium). Typical (first-generation): haloperidol, chlorpromazine—D2 blockade, high EPS/tardive dyskinesia risk, hyperprolactinemia. Atypical (second-generation): clozapine, risperidone, olanzapine, quetiapine—D2/5-HT2A antagonism, fewer EPS but more metabolic side effects (weight gain, diabetes). Long-acting injectable formulations improve adherence. Adverse effects include neuroleptic malignant syndrome and QT prolongation.

\medterm{Antipyretic}-- A drug that reduces fever by lowering the hypothalamic set point, typically by inhibiting prostaglandin synthesis in the preoptic hypothalamus. Acetaminophen and NSAIDs (ibuprofen, aspirin) are the principal antipyretics. Fever is itself a beneficial host response, so antipyretics are reserved for patient comfort, particularly in children at risk of dehydration. Aspirin is contraindicated in children (Reye syndrome). Antipyretics are used in combination or rotation for control. Overdose of acetaminophen causes hepatotoxicity.

\medterm{Antipyretics}-- Medications that reduce fever by inhibiting cyclo-oxygenase (COX) enzymes, reducing hypothalamic prostaglandin E$_2$ synthesis and resetting the thermoregulatory set point. Common antipyretics: paracetamol (acetaminophen — preferred first-line, particularly in children and patients with contraindications to NSAIDs, but hepatotoxic in overdose), ibuprofen (NSAID, also anti-inflammatory, avoid in dehydration, renal impairment, and peptic ulcer disease), and aspirin (avoid in children due to Reye syndrome risk). Antipyretics should be used for symptomatic relief, not routinely for all fevers, as fever is a physiological defence mechanism.

\medterm{Antiseptic}-- See infectious\_disease.

\medterm{Antitubercular Agents}-- Antimicrobial drugs used to treat Mycobacterium tuberculosis infections. The first-line agents include isoniazid, which inhibits mycolic acid synthesis; rifampicin, which inhibits RNA polymerase; pyrazinamide, which disrupts membrane energetics; ethambutol, which inhibits arabinosyl transferase; and streptomycin, which inhibits protein synthesis. The standard regimen for drug-susceptible tuberculosis is a four-drug combination of isoniazid, rifampicin, pyrazinamide, andethambutol for two months followed by a continuation phase of isoniazid and rifampicin for four months; treatment adherence is critical and directly observed therapy is recommended to prevent treatment failure and acquired drug resistance.

\medterm{Antitussive}-- A medication that suppresses the cough reflex, used for dry, non-productive cough. Centrally acting antitussives: dextromethorphan — acts on the cough center in the medulla oblongata, likely through sigma-1 receptor agonism, NMDA antagonism, and serotonin reuptake inhibition; codeine — a prodrug converted to morphine, acts on mu-opioid receptors in the cough center, controlled substance due to abuse potential, not recommended in children. Peripherally acting antitussives: benzonatate — a local anesthetic that reduces cough reflex by numbing stretch receptors in the respiratory tract; inhaled sodium cromoglycate and nedocromil — stabilize mast cells and reduce cough in asthma. Demulcents (honey, glycerin, lozenges) soothe irritated pharyngeal mucosa. Expectorants (guaifenesin) are not antitussives but increase mucus production to facilitate coughing up phlegm.

\medterm{Antivirals}-- Drugs that inhibit the replication of viruses, reducing the severity and duration of viral infections. They are used for conditions such as influenza, herpes, HIV, hepatitis B and C, and COVID-19. Unlike antibiotics, antivirals target specific viral enzymes or life cycle stages and often work best when started early in the infection.

\medterm{Aromatase Inhibitor}-- A class of medications that inhibit aromatase, the enzyme that converts androgens (androstenedione, testosterone) to estrogens (estrone, estradiol) in peripheral tissues (adipose, breast, muscle, bone). Used primarily in hormone receptor-positive breast cancer in postmenopausal women. Third-generation aromatase inhibitors: anastrozole (nonsteroidal, reversible), letrozole (nonsteroidal, reversible), and exemestane (steroidal, irreversible). In postmenopausal women, AIs are first-line adjuvant therapy and first-line therapy for metastatic breast cancer (superior to tamoxifen). In premenopausal women, AIs are used in combination with ovarian suppression. Adverse effects: arthralgia and myalgia (20-50\%, most common cause of discontinuation), accelerated bone loss (osteoporosis, fracture risk), hot flashes, and dyslipidemia. AIs do not cause endometrial cancer or thromboembolism (unlike tamoxifen).

\medterm{Azole Antifungals}-- A class of antifungal agents that inhibit the enzyme lanosterol 14-alpha-demethylase, a cytochrome P450 enzyme essential for ergosterol synthesis in fungal cell membranes. The depletion of ergosterol and accumulation of toxic methylated sterols disrupt membrane integrity and function. Azoles are divided into imidazoles including ketoconazole and clotrimazole, used primarily for topical infections, and triazoles including fluconazole, itraconazole, voriconazole, and posaconaposaconazole, which provide broader antifungal coverage; posaconazole and voriconazole have activity against Aspergillus species, while fluconazole has limited activity against this genus.

\medterm{Bacteriophage therapy}-- The use of bacterial viruses (bacteriophages) to treat bacterial infections. Bacteriophages are viruses that specifically infect and lyse bacteria. They have narrow specificity for target bacteria and do not affect normal flora. Phage therapy is being studied for multidrug-resistant bacterial infections, particularly those failing conventional antibiotics. Clinical trials have shown promise for chronic pseudomonal infections and prosthetic joint infections. The main challenges include regulatoryhurdles including standardised production, stability testing, pharmacokinetic characterisation, and regulatory approval pathways that vary across jurisdictions.

\medterm{Barbiturate}-- A class of CNS depressants derived from barbituric acid that enhance GABA-A receptor activity by increasing the duration of chloride channel opening (as opposed to benzodiazepines, which increase the frequency). At low doses: anxiolysis and sedation. At moderate doses: hypnosis and anesthesia. At high doses: respiratory depression, coma, and death. Classification by duration of action: ultra-short (thiopental, methohexital — IV induction agents), short/intermediate (pentobarbital, secobarbital — sedation, now rarely used), and long-acting (phenobarbital — anticonvulsant). Indications now limited: phenobarbital for neonatal seizures and status epilepticus when other agents fail; thiopental for rapid sequence induction historically. Adverse effects: narrow therapeutic index, respiratory depression, tolerance, physical dependence, severe withdrawal (anxiety, tremors, seizures, delirium), and enzyme induction (CYP3A4, CYP2C9). Drug interactions: reduces efficacy of oral contraceptives, warfarin, and many other drugs. Barbiturates are Schedule II-IV controlled substances.

\medterm{Benzodiazepines}-- A class of drugs that enhance the effect of the inhibitory neurotransmitter gamma-aminobutyric acid at the GABA-A receptor, increasing the frequency of chloride channel opening and producing anxiolytic, sedative, anticonvulsant, and muscle relaxant effects. They bind to an allosteric site on the GABA-A receptor, distinct from the GABA binding site, and potentiate the effect of GABA when it binds. Common benzodiazepines include diazepam, lorazepam, alprazolam, clonazepam, midamidazolam. They differ in potency, half-life, lipid solubility, and metabolic pathways, influencing their clinical applications for anxiety, insomnia, seizure disorders, and procedural sedation.

\medterm{Beta Blockers}-- A class of medications that block the effects of adrenaline on the body's beta receptors. They slow down the heart rate and reduce blood pressure, making them widely used for treating heart conditions, high blood pressure, and anxiety.

\medterm{Beta-Lactam Antibiotics}-- A large class of antimicrobial agents characterized by a beta-lactam ring in their molecular structure, which is essential for their bactericidal activity. They include penicillins, cephalosporins, carbapenems, and monobactams, all of which share the mechanism of inhibiting bacterial cell wall synthesis by binding to penicillin-binding proteins and preventing cross-linking of peptidoglycan chains. Beta-lactams are bactericidal, killing bacteria by weakening the cell wwall, leading to osmotic lysis and cell death; they are the most widely used class of antibiotics globally, with spectrum determined by their affinity for specific PBPs and their ability to penetrate the outer membrane of gram-negative organisms.

\medterm{Beta-lactamase inhibitor combinations}-- Combinations of beta-lactam antibiotics with beta-lactamase inhibitors that extend activity against beta-lactamase producing organisms. The major combinations include amoxicillin-clavulanate, ampicillin-sulbactam, piperacillin-tazobactam, cefoperazone-sulbactam, and ceftolozane-tazobactam. Clavulanate, sulbactam, and tazobactam are beta-lactamase inhibitors that bind irreversibly to beta-lactamase enzymes, preventing degradation of the partner beta-lactam. Clavulanate has the strongest beta-lactactamase inhibition; tazobactam is the most potent beta-lactamase inhibitor in clinical use, while avibactam is a newer non-beta-lactam inhibitor that effectively restores activity against class A and class C beta-lactamases.

\medterm{Biguanide}-- A class of oral antidiabetic drugs derived from guanidine, with metformin being the only biguanide in current clinical use (phenformin and buformin were withdrawn due to lactic acidosis risk). Metformin is the first-line pharmacotherapy for type 2 diabetes mellitus. Mechanism: activates AMP-activated protein kinase, reducing hepatic gluconeogenesis and glycogenolysis; improves peripheral insulin sensitivity (increases glucose uptake in skeletal muscle); and reduces intestinal glucose absorption. Metformin does not stimulate insulin secretion and therefore does not cause hypoglycemia as monotherapy. Effects: lowers HbA1c by 1-1.5\%, weight neutral or modest weight loss, and improves lipid profile. Adverse effects: GI intolerance (diarrhea, nausea, metallic taste, 20-30\%), vitamin B12 deficiency with chronic use, and lactic acidosis (rare, 3-10 per 100,000 patient-years, primarily in patients with contraindications). Contraindications: eGFR <30 mL/min, acute kidney injury, decompensated liver disease, and acute medical illness with hypoxia (sepsis, MI, shock). Extended-release formulations improve GI tolerability.

\medterm{Bioavailability}-- The fraction of an administered drug that reaches the systemic circulation in its unchanged active form. It is expressed as a percentage and is determined by the route of administration and the extent of first-pass metabolism. For intravenous administration, bioavailability is 100 percent by definition. For oral administration, bioavailability is typically less than 100 percent due to incomplete absorption and pre-systemic metabolism by gut wall enzymes and the liver. The fraction of an oral dose reaching systemic circulation can be increased by altering the formulation, using prodrugs, or co-administering with food or other drugs that affect absorption or first-pass metabolism.

\medterm{Bioequivalence}-- The property of two drug products (brand-name and generic) that have the same rate and extent of absorption when administered at the same molar dose under similar conditions, resulting in comparable bioavailability and therapeutic effect. The FDA requires generic drugs to demonstrate bioequivalence to the reference listed drug through pharmacokinetic studies showing that the 90\% confidence interval of the ratio of generic-to-brand for AUC and Cmax falls within 80-125\%. Bioequivalence ensures that generic drugs can be substituted for brand-name drugs without dose adjustment, providing the same clinical efficacy and safety profile at a lower cost.

\medterm{Biotransformation}-- The process by which drugs and other xenobiotics are chemically modified by metabolic enzymes, primarily in the liver, to facilitate elimination from the body. Biotransformation occurs in two phases. Phase I reactions (functionalization): oxidation, reduction, and hydrolysis catalyzed primarily by cytochrome P450 enzymes (CYP1A2, CYP2C9, CYP2C19, CYP2D6, CYP3A4), introducing or exposing functional groups (e.g., -OH, -NH2, -COOH). Phase II reactions (conjugation): glucuronidation (UGT), sulfation (SULT), acetylation (NAT), methylation, and conjugation with glutathione (GST), attaching endogenous hydrophilic molecules to the drug or Phase I metabolite, increasing water solubility and promoting renal or biliary excretion. Biotransformation generally produces inactive, water-soluble metabolites, but some drugs undergo bioactivation (prodrugs) or produce toxic metabolites that contribute to drug toxicity.

\medterm{Bisacodyl}-- A stimulant laxative (diphenylmethane derivative) that stimulates colonic motility and enhances water/electrolyte secretion, producing a bowel movement within 6-12 hours (oral) or 15-60 minutes (suppository). Used for constipation, bowel preparation, and evacuation before procedures. Adverse effects: abdominal cramps, diarrhea, and electrolyte disturbances with overuse; chronic stimulant use may cause melanosis coli and dependence. Generally safe short-term; avoid in intestinal obstruction, acute abdomen, and with caution in inflammatory bowel disease.

\medterm{Bismuth Subsalicylate}-- An antidiarrheal and anti-ulcer agent (e.g., Pepto-Bismol) with antimicrobial, anti-inflammatory, and coating actions. Used for traveler's diarrhea (prophylaxis and treatment), dyspepsia, Helicobacter pylori eradication (quadruple therapy), and symptoms of indigestion. Adverse effects: black stools and black tongue (harmless, from bismuth sulfide), constipation, and tinnitus (salicylate); avoid in children/adolescents with viral illness (Reye syndrome risk) and with anticoagulants/aspirin.

\medterm{Blood-Brain Barrier}-- The blood-brain barrier is a highly selective semipermeable membrane that separates the circulating blood from the brain extracellular fluid, consisting of endothelial cells connected by tight junctions, astrocyte foot processes, and a basement membrane. The BBB restricts the passage of most molecules from the blood into the brain, protecting the central nervous system from toxins and pathogens while maintaining a stable environment for neural function. Only small, lipophilic molecules and gasescan cross the BBB via passive diffusion or active transport; the BBB is disrupted in pathological conditions such as meningitis, ischemia, and multiple sclerosis, affecting drug distribution to the brain.

\medterm{Bromocriptine}-- An ergot dopamine agonist (D2) used for hyperprolactinemia (prolactinomas, galactorrhea), acromegaly (adjunct), Parkinson disease (less common now), and postpartum breast engorgement (historical). Also used for neuroleptic malignant syndrome. Adverse effects: nausea, orthostatic hypotension, dizziness, headache, cardiac valvulopathy (ergot, high doses), retroperitoneal fibrosis, and hallucinations. Cabergoline (longer-acting, better tolerated) is preferred for hyperprolactinemia.

\medterm{Bumetanide}-- A potent loop diuretic inhibiting the Na-K-2Cl cotransporter in the thick ascending limb of the loop of Henle. Used for edema in heart failure, cirrhosis, and renal disease, and for hypertension in renal impairment. Approximately 40 times more potent than furosemide; more bioavailable orally, so equivalent IV-to-oral dosing is possible. Adverse effects: hypokalemia, hyponatremia, hypomagnesemia, hypocalcemia, metabolic alkalosis, ototoxicity (with aminoglycosides), and dehydration/azotemia.

\medterm{Bupropion}-- An atypical antidepressant that inhibits dopamine and norepinephrine reuptake. Used for major depressive disorder, seasonal affective disorder, and smoking cessation (Zyban). Low sexual side effects and weight neutrality; energizing. Also used off-label for ADHD. Adverse effects: insomnia, dry mouth, agitation, lowered seizure threshold (avoid in epilepsy, eating disorders, alcohol/benzodiazepine withdrawal), and hypertension. Contraindicated with MAOIs and bupropion-containing products. Sustained-release forms allow daily dosing.

\medterm{Caffeine}-- A naturally occurring xanthine alkaloid and central nervous system stimulant found in coffee, tea, chocolate, and many soft drinks and energy beverages. Mechanism: adenosine receptor antagonist (primarily A1 and A2A subtypes), increasing neuronal activity and promoting catecholamine release. Effects: increased alertness, reduced fatigue, improved cognitive performance, diuresis, and bronchodilation. Half-life: 3--7 hours in adults, prolonged in pregnancy, liver disease, and with oral contraceptives. Chronic use produces tolerance; abrupt cessation causes withdrawal syndrome (headache, fatigue, irritability, dysphoria) lasting 2--9 days. Acute toxicity (>500 mg): anxiety, insomnia, tremor, tachycardia; severe toxicity (>5 g): seizures, arrhythmias, death.

\medterm{Carbapenem-resistant organisms}-- Bacteria that are resistant to carbapenem antibiotics, typically through production of carbapenemases. The most important carbapenemases include KPC (Klebsiella pneumoniae carbapenemase), NDM (New Delhi metallo-beta-lactamase), and OXA-48-like enzymes. CRO are a major public health threat due to limited treatment options. Infections are associated with high mortality rates. Treatment options include ceftazidime-avibactam, meropenem-vaborbactam, and colistin-based combinations. Infection prevention and antimicrobial stewardship are critical for controlling the spread of carbapenem-resistant organisms in healthcare settings.

\medterm{Carbidopa}-- A peripheral dopa decarboxylase inhibitor that does not cross the blood-brain barrier; combined with levodopa to prevent its peripheral conversion to dopamine, reducing side effects (nausea, hypotension) and increasing CNS levodopa delivery. Standard combinations: carbidopa/levodopa (Sinemet), including controlled-release and orally disintegrating forms. Higher carbidopa ratios (25/100) reduce 'wearing off'. Carbidopa itself has no antiparkinsonian effect but reduces peripheral dopamine-related adverse effects.

\medterm{Carbolic Acid (Phenol)}-- An aromatic organic compound (C6H5OH) with antiseptic, disinfectant, and caustic properties. Historically significant as the first surgical antiseptic introduced by Joseph Lister in 1867, dramatically reducing post-operative infection and mortality. Mechanism: denatures bacterial proteins and disrupts cell membranes. Phenol is bactericidal, fungicidal, and virucidal but is toxic to human tissues at high concentrations causing chemical burns and systemic toxicity (neurologic depression, cardiovascular collapse). Modern medical uses: phenol-based chemical peels (dermatology), nerve ablation for chronic pain (neurolysis), and as a preservative in some injectable medications. Safer, less toxic antiseptics (chlorhexidine, povidone-iodine) have largely replaced it for routine disinfection.

\medterm{Carvedilol}-- A nonselective beta-blocker with additional alpha-1 adrenergic blockade and antioxidant properties, producing vasodilation with balanced reduction of heart rate, contractility, and afterload. Used for heart failure with reduced ejection fraction (improves survival, NYHA II-IV), hypertension, and post-myocardial infarction ventricular dysfunction. Also treats angina. Metabolized by CYP2D6; must be titrated slowly (e.g., 3.125 mg BID upward). Adverse effects: dizziness, hypotension, bradycardia, fatigue, bronchospasm, and masking of hypoglycemia.

\medterm{Caspofungin}-- An echinocandin antifungal agent used in the treatment of invasive candidiasis, esophageal candidiasis, and invasive aspergillosis (as salvage therapy). Caspofungin inhibits beta-(1,3)-D-glucan synthase, disrupting fungal cell wall synthesis. The drug is available only intravenously and has a half-life of approximately 9-11 hours, allowing for once-daily dosing. Common adverse effects include headache, nausea, and phlebitis at the infusion site. Fever and rash have been reported. Hepatotoxicity has been reported, and the drug has fewer drug interactions than azole antifungals as it is not a significant CYP inhibitor.

\medterm{Ceftazidime-avibactam}-- A combination of the third-generation cephalosporin ceftazidime with the non-beta-lactam beta-lactamase inhibitor avibactam. Avibactam inhibits class A, class C, and some class D beta-lactamases, including KPC carbapenemases. This combination has activity against many carbapenem-resistant Enterobacteriaceae and Pseudomonas aeruginosa. It is indicated for complicated intra-abdominal infections, complicated urinary tract infections, and hospital-acquired pneumonia. The combination is administered administered intravenously and requires dose adjustment based on renal function; the most common adverse effects include nausea, diarrhea, and headache.

\medterm{Cephalosporin}-- A class of beta-lactam antibiotics derived from the fungus Cephalosporium acremonium, with a mechanism identical to penicillins — inhibition of bacterial cell wall synthesis by binding to penicillin-binding proteins and inhibiting transpeptidase. Cephalosporins are classified into five generations based on their antimicrobial spectrum. First generation (cefazolin, cephalexin): good gram-positive coverage (Staph aureus, Strep pyogenes), limited gram-negative coverage. Second generation (cefuroxime, cefoxitin, cefotetan): expanded gram-negative coverage including H. influenzae and anaerobes (cefoxitin). Third generation (ceftriaxone, cefotaxime, ceftazidime): broad gram-negative coverage including Enterobacteriaceae; ceftazidime also covers Pseudomonas. Fourth generation (cefepime): broad coverage including Pseudomonas and gram-positive with stability against AmpC beta-lactamases. Fifth generation (ceftaroline): anti-MRSA activity plus standard cephalosporin coverage. Adverse effects: hypersensitivity reactions (cross-reactivity with penicillins ~10\% based on similar beta-lactam ring structure), GI disturbances, and bleeding (cefotetan, due to N-methylthiotetrazole side chain causing hypoprothrombinemia).

\medterm{Cetirizine}-- A second-generation H1 antihistamine (metabolite of hydroxyzine) used for allergic rhinitis and urticaria. Slightly more sedating than loratadine/fexofenadine (some CNS penetration). Once daily; renally cleared (dose-reduce in renal impairment). Adverse effects: somnolence (dose-related), dry mouth, and fatigue. Available OTC. Levocetirizine is the active enantiomer.

\medterm{Chloramphenicol}-- A broad-spectrum antibiotic with activity against gram-positive and gram-negative organisms, as well as Rickettsia. Chloramphenicol inhibits bacterial protein synthesis by binding to the 50S ribosomal subunit. It is indicated for serious infections when other agents are not appropriate, including bacterial meningitis and typhoid fever. Common adverse effects include bone marrow suppression, gray baby syndrome in neonates, and aplastic anemia. The drug is rarely used due to its toxicity profile. Although effective, its use is limited to life-threatening infections due to the availability of safer alternatives.

\medterm{Citalopram}-- A selective serotonin reuptake inhibitor (SSRI) used for major depressive disorder and (off-label) anxiety disorders. Well tolerated; lowest cost among SSRIs. Adverse effects: nausea, insomnia, sexual dysfunction, and dose-dependent QT prolongation (maximum 40 mg/day; 20 mg in elderly/CYP2C19 poor metabolizers—torsades risk). Serotonin syndrome and discontinuation syndrome; black-box warning for suicidality in young adults. Escitalopram (S-enantiomer) has fewer QT effects.

\medterm{Clearance}-- The volume of plasma from which a drug is completely removed per unit time, typically expressed in milliliters per minute or liters per hour. Total clearance is the sum of all clearance mechanisms, primarily hepatic and renal clearance. Hepatic clearance depends on hepatic blood flow, enzyme activity, and protein binding, with the hepatic extraction ratio determining the contribution of blood flow versus enzyme activity to clearance. High-extraction drugs like lidocaine have clearanclearance that is highly dependent on hepatic blood flow, while low-extraction drugs like warfarin have clearance that is primarily determined by intrinsic enzyme activity and protein binding.

\medterm{Clinical Trial Phases}-- The stages of drug development: Phase I (safety, dosing, pharmacokinetics in 20-100 healthy volunteers); Phase II (efficacy and side effects in 100-300 patients with the target condition); Phase III (large-scale efficacy and safety comparison with standard treatment in 1,000-3,000 patients); Phase IV (post-marketing surveillance for long-term effects and rare adverse events). Each phase has specific endpoints and success criteria before advancing to the next phase.

\medterm{Clostridioides difficile}-- An anaerobic, spore-forming, toxin-producing gram-positive bacillus causing antibiotic-associated diarrhea and pseudomembranous colitis. Risk factors: antibiotics (especially clindamycin, fluoroquinolones, cephalosporins), hospitalization, age >65, PPIs, and immunosuppression. Toxins A and B damage colonic epithelium. Diagnosis: NAAT/PCR for toxin genes or toxin EIA. Treatment: vancomycin or fidaxomicin PO (first-line); metronidazole for mild disease when agents unavailable; fecal microbiota transplantation for recurrent disease. Infection control: contact precautions, handwashing with soap (alcohol gel is less effective).

\medterm{Clostridioides difficile infection}-- An infectious diarrhea caused by the bacterium Clostridioides difficile, typically occurring after disruption of normal gastrointestinal flora by antibiotics. C. difficile produces toxins A and B that cause colonic inflammation and diarrhea. Risk factors include antibiotic exposure, proton pump inhibitor use, and advanced age. Treatment options include vancomycin, fidaxomicin, and metronidazole. Recurrence occurs in approximately 25\% of cases. Fecal microbiota transplantation is highly effectivefor recurrent C. difficile infection, with cure rates exceeding 90\%.

\medterm{Clostridium difficile Infection (CDI)}-- An anaerobic, spore-forming, toxin-producing gram-positive bacillus causing antibiotic-associated diarrhea and pseudomembranous colitis. See Clostridioides difficile.

\medterm{Competitive Antagonist}-- A competitive antagonist is a drug that binds to the same receptor site as an agonist and competes for binding, reducing the effect of the agonist without producing an effect of its own. The antagonism is surmountable by increasing the agonist concentration, which shifts the dose-response curve of the agonist to the right without reducing the maximal response. This means that with enough agonist, the full effect can still be achieved. The degree of rightward shift depends on the concentration anand the concentration of the antagonist; the relationship is quantitated using the Schild regression, which yields a pA2 value representing the concentration of antagonist required to produce a twofold rightward shift of the agonist dose-response curve.

\medterm{Controlled Substance}-- A drug or chemical whose manufacture, possession, and use are regulated by government authorities due to its potential for abuse, dependence, or harm. In the United States, the Controlled Substances Act (CSA) classifies controlled substances into five schedules (I-V) based on their accepted medical use, abuse potential, and dependence liability. Schedule I: high abuse potential, no accepted medical use (heroin, LSD, marijuana federally). Schedule II: high abuse potential, severe dependence liability, accepted medical use with restrictions (morphine, oxycodone, fentanyl, cocaine, amphetamine). Schedule III: moderate to low abuse potential, moderate dependence liability (codeine-containing products, ketamine, anabolic steroids). Schedule IV: low abuse potential, limited dependence liability (benzodiazepines, zolpidem, tramadol). Schedule V: lowest abuse potential, limited dependence liability (codeine-containing antitussives, pregabalin). Healthcare providers must register with the DEA to prescribe controlled substances and comply with state and federal regulations.

\medterm{Corticosteroids Mechanism}-- Corticosteroids exert their anti-inflammatory and immunosuppressive effects primarily by binding to intracellular glucocorticoid receptors, which then translocate to the nucleus and modulate gene transcription. They transactivate anti-inflammatory genes including those encoding lipocortin-1, which inhibits phospholipase A2 and reduces arachidonic acid release, and they transrepress pro-inflammatory genes by inhibiting transcription factors including NF-kappaB and AP-1. This reduces the productioproduction of prostaglandins, leukotrienes, and other inflammatory mediators, which underlies their therapeutic efficacy in inflammatory conditions, autoimmune diseases, and allergic reactions.

\medterm{COX Inhibitor}-- A class of drugs that inhibit cyclooxygenase enzymes (COX-1 and COX-2), which catalyze the conversion of arachidonic acid to prostaglandins and thromboxanes. COX-1 is constitutively expressed in most tissues, producing prostaglandins that protect the gastric mucosa, maintain renal blood flow, and support platelet aggregation. COX-2 is induced during inflammation, producing prostaglandins that mediate pain, fever, and inflammation. Non-selective COX inhibitors (aspirin, ibuprofen, naproxen, indomethacin, diclofenac) inhibit both isoforms, providing analgesic, anti-inflammatory, and antipyretic effects but causing GI toxicity (gastritis, ulcers, bleeding) and platelet inhibition. Selective COX-2 inhibitors (celecoxib) provide anti-inflammatory effects with reduced GI toxicity but do not inhibit platelet aggregation and are associated with increased cardiovascular risk (MI, stroke). The COX hypothesis explains the therapeutic and adverse effects of NSAIDs.

\medterm{Cream}-- A semisolid emulsion (oil in water or water in oil) for topical application, containing a drug in an emulsion base. Creams are lighter, less greasy, and easier to spread than ointments, and are preferred for moist or oozing skin lesions. They penetrate less deeply than ointments and are cosmetically acceptable. Used for dermatologic conditions (corticosteroid creams, antifungals), analgesics, and moisturizers. Formulation (cream, ointment, gel, lotion, solution) affects drug delivery, occlusion, and patient acceptability.

\medterm{CYP2C19 Polymorphism}-- CYP2C19 is a cytochrome P450 enzyme involved in the metabolism of approximately 10 percent of drugs, including the antiplatelet agent clopidogrel, proton pump inhibitors, some antidepressants, and the antiepileptic drug phenytoin. Genetic polymorphisms result in poor metabolizer phenotypes in 2 to 5 percent of Caucasians and 15 to 20 percent of Asian populations. Poor metabolizers have reduced conversion of clopidogrel, a prodrug requiring CYP2C19-mediated activation, to its active metabolite, lleading to reduced antiplatelet efficacy and increased risk of stent thrombosis in poor metabolisers; alternative antiplatelet agents such as ticagrelor or prasugrel, which are not dependent on CYP2C19 activation, are preferred in such patients.

\medterm{CYP2D6 Polymorphism}-- CYP2D6 is a cytochrome P450 enzyme responsible for the metabolism of approximately 25 percent of commonly prescribed drugs, including codeine, tamoxifen, many antidepressants, and antipsychotics. Genetic polymorphisms in the CYP2D6 gene result in highly variable enzyme activity among individuals, classified into four phenotype groups: poor metabolizers with two non-functional alleles and absent enzyme activity; intermediate metabolizers with reduced activity; extensive metabolizers with normal aactivity, and ultrarapid metabolisers with increased enzyme activity due to gene duplication. This phenotypic variability has significant clinical implications for drugs metabolized by CYP2D6: codeine prodrug activation to morphine is reduced in poor metabolisers and enhanced in ultrarapid metabolisers, who may experience opioid toxicity at standard doses; tamoxifen efficacy is reduced in poor metabolisers; and antidepressant dosing may require individualisation based on CYP2D6 phenotype.

\medterm{Cytochrome P450}-- The cytochrome P450 system is a superfamily of heme-containing monooxygenase enzymes primarily located in the endoplasmic reticulum of hepatocytes, responsible for Phase I drug metabolism through oxidation, reduction, and hydrolysis reactions. The most important CYP isoforms in drug metabolism include CYP3A4, which metabolizes approximately 50 percent of all drugs; CYP2D6, responsible for approximately 25 percent of drug metabolism; CYP2C9, CYP2C19, CYP1A2, and CYP2E1. CYP enzymes can be inducedinduced or inhibited by various drugs and xenobiotics, leading to clinically significant drug interactions; enzyme induction occurs over days to weeks through increased gene transcription, while enzyme inhibition occurs rapidly through direct competition for the active site.

\medterm{Dabigatran}-- A direct oral anticoagulant that reversibly inhibits thrombin (factor IIa). Used for stroke prevention in nonvalvular atrial fibrillation and treatment/prophylaxis of venous thromboembolism. Fixed dosing, rapid onset, predictable pharmacokinetics, no routine monitoring. Prodrug requiring acid environment; capsules must not be opened. Adverse effects: bleeding, dyspepsia. Reversal agent: idarucizumab (specific antibody). Avoid in mechanical heart valves (valve thrombosis risk) and severe renal impairment.

\medterm{Daptomycin resistance}-- Resistance to daptomycin can develop through mutations in genes encoding cell membrane phospholipid synthesis. The mechanism involves changes in the bacterial cell membrane that reduce daptomycin binding. Resistance is more common in VRE than MRSA. Resistance can emerge during therapy, particularly in endocarditis and bacteremia. Daptomycin resistance is typically accompanied by vancomycin resistance, a phenomenon known as daptomycin-vancomycin cross-resistance. Combination therapy may be requirred in selected cases and careful monitoring of creatine kinase levels is recommended during daptomycin therapy.

\medterm{Delamanid}-- A nitroimidazole derivative used in the treatment of multidrug-resistant tuberculosis. Delamanid inhibits mycolic acid synthesis in Mycobacterium tuberculosis, disrupting the cell wall. It is indicated for pulmonary multidrug-resistant tuberculosis in combination with other agents. The drug is administered orally twice daily. Common adverse effects include headache, nausea, and QT prolongation. Delamanid is metabolized by plasma esterases and has minimal CYP450 interactions. The drug has demonstonstrated efficacy in MDR-TB treatment cohorts, with improved outcomes compared to background regimens alone.

\medterm{Delavirdine}-- A non-nucleoside reverse transcriptase inhibitor (NNRTI) that was previously used in combination antiretroviral therapy for HIV-1 infection but is now rarely used due to the availability of more effective and better-tolerated agents. Delavirdine binds to the HIV reverse transcriptase enzyme at a site distinct from the active site, causing a conformational change that inhibits enzymatic activity. The drug has relatively low potency compared to newer NNRTIs and requires dosing three times daily. CCommon adverse effects include rash, elevated transaminases, and nausea, and the drug has significant CYP3A4 inhibitory activity leading to numerous drug interactions, limiting its use in contemporary practice.

\medterm{Dextroamphetamine}-- A central nervous system stimulant (d-amphetamine isomer) increasing dopamine and norepinephrine. Used for ADHD and narcolepsy. Available as immediate and extended-release, and in mixed amphetamine salts (Adderall) and lisdexamfetamine (prodrug). Adverse effects: appetite suppression, insomnia, growth delay, tachycardia, hypertension, anxiety, and abuse potential (Schedule II). Cardiovascular screening (ECG for risk factors) before initiation; monitor growth in children.

\medterm{Didanosine}-- A nucleoside reverse transcriptase inhibitor (NRTI) used in the treatment of HIV-1 infection, typically as part of combination antiretroviral therapy. Didanosine is an adenosine analogue that inhibits viral reverse transcriptase, leading to chain termination of viral DNA synthesis. The drug requires acidic pH for stability and must be administered with buffer tablets or antacids to prevent gastric acid degradation. Didanosine has a bioavailability of approximately 30-40\% when taken with buffer, buffer, and is dosed twice daily; common adverse effects include peripheral neuropathy, pancreatitis, and lactic acidosis, with pancreatitis being a potentially fatal complication that requires immediate discontinuation.

\medterm{Diphenhydramine}-- A first-generation H1 antihistamine with potent anticholinergic and sedative effects. Used for allergic reactions, urticaria/prpruritus, motion sickness, insomnia (OTC sleep aid), and as an antitussive; IV form for anaphylaxis (adjunct) and acute dystonic reactions. Adverse effects: sedation, dry mouth, blurred vision, urinary retention, constipation, confusion (elderly—avoid), and anticholinergic toxicity in overdose (treated with physostigmine). Metabolized by CYP2D6.

\medterm{Disulfiram}-- An aversive agent used to treat alcohol use disorder. Inhibits acetaldehyde dehydrogenase, causing accumulation of acetaldehyde when alcohol is consumed—producing flushing, tachycardia, nausea, vomiting, and hypotension ('acetaldehyde syndrome'). Provides a strong deterrent to drinking but requires motivation and abstention. Adverse effects: hepatotoxicity, neuropathy, and psychosis. Contraindicated with alcohol-containing products, metronidazole, and severe cardiac/hepatic disease. Must warn patients to avoid all alcohol, including in foods and medications.

\medterm{Diuretic}-- A drug that increases urine production, used for edema (heart failure, cirrhosis, renal disease), hypertension, and electrolyte disorders. Classes: loop diuretics (furosemide, bumetanide, torsemide—most potent, act on the loop of Henle), thiazides (hydrochlorothiazide, chlorthalidone—distal tubule), potassium-sparing (spironolactone, eplerenone, amiloride—collecting duct), and osmotic (mannitol) and carbonic anhydrase inhibitors (acetazolamide). Adverse effects: electrolyte disturbances (hypokalemia, hyponatremia), dehydration, hypotension, and ototoxicity (loop diuretics). Spironolactone/eplerenone improve outcomes in heart failure.

\medterm{DOACs}-- Direct oral anticoagulants are a class of anticoagulant drugs that directly inhibit specific coagulation factors, providing an alternative to warfarin with more predictable pharmacokinetics and no routine monitoring. The factor Xa inhibitors include rivaroxaban, apixaban, and edoxaban, which directly and reversibly bind to the active site of factor Xa. The direct thrombin inhibitor dabigatran directly binds to and inhibits thrombin. DOACs have rapid onset of action, predictable pharmacokinetics,pharmacokinetics, fixed dosing, fewer drug and food interactions than warfarin, and no requirement for routine coagulation monitoring, though they are more expensive and specific reversal agents are limited.

\medterm{Docusate}-- A stool softener (surfactant laxative) that increases water and fat incorporation into stool by lowering surface tension, facilitating defecation in patients who should avoid straining (post-MI, post-surgery, hemorrhoids) and for opioid-related constipation. Onset 12-72 hours. Limited efficacy as a true laxative; better for softening than for stimulating evacuation. Generally well tolerated; may cause mild cramping. Often combined with stimulant laxatives.

\medterm{Doravirine}-- A non-nucleoside reverse transcriptase inhibitor (NNRTI) used in combination antiretroviral therapy for HIV-1 infection. Doravirine binds to the HIV reverse transcriptase enzyme and inhibits its activity through a unique mechanism that maintains activity against strains resistant to first-generation NNRTIs. The drug has a favorable tolerability profile with minimal impact on lipid levels and body weight compared to other antiretrovirals. Common adverse effects include headache, nausea, and diarrdiarrhoea, and the drug has fewer drug interactions than other NNRTIs as it is not a significant inducer or inhibitor of CYP enzymes.

\medterm{Dosage}-- The amount and frequency of a drug administered to achieve a therapeutic effect while minimizing toxicity. Components: dose (amount, e.g., mg), dosing interval (frequency), route, and duration. Dosing is individualized using weight (mg/kg, especially children), body surface area (chemotherapy), renal and hepatic function (clearance), age, and therapeutic drug monitoring for narrow-therapeutic-index drugs. Loading doses achieve steady state rapidly; maintenance doses sustain it. Overdosing causes toxicity, underdosing causes therapeutic failure and antimicrobial resistance.

\medterm{Dose-Response Curve}-- A graphical representation of the relationship between the dose of a drug and the magnitude of its biological effect. The x-axis represents dose (typically on a logarithmic scale), and the y-axis represents the response (as percentage of maximum effect). The curve is sigmoidal in shape, characterized by: threshold (minimum dose producing a measurable effect), slope (steepness of the curve, indicating how rapidly the effect changes with dose), maximal efficacy (the plateau, Emax), and potency (the dose producing 50\% of maximal effect, ED50). The therapeutic window is the range between the ED50 for the desired effect and the TD50 (toxic dose in 50\% of patients). The therapeutic index (TI = TD50 / ED50) measures relative safety. Dose-response curves are essential for determining appropriate dosing regimens and comparing drug potency and efficacy.

\medterm{Dose-Response Relationship}-- The dose-response relationship describes the quantitative relationship between the dose of a drug and the magnitude of the pharmacological response it produces. This relationship is typically represented as a sigmoidal curve when response is plotted against the log of the dose. The curve is characterized by several key parameters: the maximal effect or efficacy represents the maximum response achievable; the EC50 or ED50 represents the dose producing 50 percent of the maximal response and indicaindicates the potency of the drug. The slope of the curve reflects the mechanism of drug action, with steeper slopes indicating a more direct relationship between dose and response, as seen with drugs that act on a single receptor type.

\medterm{Doxazosin}-- A selective alpha-1 adrenergic blocker causing vascular smooth muscle relaxation and reduced peripheral resistance. Used for hypertension and benign prostatic hyperplasia (relaxes prostatic smooth muscle). Extended-release formulation reduces first-dose hypotension. Adverse effects: orthostatic hypotension, dizziness, syncope (especially first dose), nasal congestion, and fatigue. Alpha-blockers are not first-line for hypertension but useful in men with BPH and as add-on therapy. Potential association with heart failure in some trials—use with caution in patients with heart failure.

\medterm{DPP-4 Inhibitors}-- Dipeptidyl peptidase-4 inhibitors that prevent degradation of endogenous GLP-1 and GIP, enhancing incretin effects on insulin secretion. Examples include sitagliptin, saxagliptin, linagliptin, and alogliptin. They are well-tolerated with low hypoglycemia risk and weight neutrality. They have modest glucose-lowering effects. Saxagliptin was associated with increased heart failure hospitalizations in the SAVOR-TIMI 53 trial. Linagliptin does not require dose adjustment for renal impairment.

\medterm{Drug Allergy}-- An adverse immune-mediated reaction to a medication, representing a hypersensitivity response that is distinct from side effects, toxicities, and idiosyncratic reactions. Drug allergies can be IgE-mediated causing immediate reactions including urticaria, angioedema, and anaphylaxis within minutes to hours; T-cell mediated causing delayed reactions including contact dermatitis, drug rash with eosinophila and systemic symptoms, and Stevens-Johnson syndrome occurring days to weeks after exposure; diagnosis involves skin testing, in vitro assays for specific IgE, and oral challenge in selected cases, with management based on avoidance of the offending drug and use of alternative agents.

\medterm{Drug Classification Systems}-- Drugs are classified using multiple complementary systems. The Anatomical Therapeutic Chemical (ATC) Classification System, maintained by the WHO, groups drugs into five levels: anatomical main group (e.g., cardiovascular system), therapeutic subgroup (e.g., antihypertensives), pharmacological subgroup (e.g., beta-blockers), chemical subgroup (e.g., selective beta-blockers), and chemical substance (e.g., metoprolol). The Therapeutic Classification groups drugs by their clinical use (e.g., antihypertensives, antidepressants, antibiotics). The Mechanism of Action Classification groups drugs by their molecular target (e.g., ACE inhibitors, calcium channel blockers, SSRIs). The Chemical Structure Classification groups drugs by their molecular composition (e.g., benzodiazepines, statins, penicillins). The Regulatory Classification includes prescription vs. over-the-counter status and controlled substance schedules. The Combination of these systems provides a comprehensive framework for understanding drug actions, interactions, and therapeutic applications.

\medterm{Drug Dependence}-- A state of physiological or psychological adaptation to a drug, characterized by a withdrawal syndrome upon discontinuation or dose reduction. Physical dependence: the body has adapted to the presence of the drug, and abrupt cessation produces a predictable withdrawal syndrome specific to the drug class. Opioid withdrawal: anxiety, mydriasis, lacrimation, rhinorrhea, piloerection, yawning, diarrhea, and muscle cramps. Alcohol withdrawal: anxiety, tremor, autonomic hyperactivity, seizures, and delirium tremens. Benzodiazepine withdrawal: anxiety, insomnia, tremor, and seizures. Psychological (behavioral) dependence: a compulsive pattern of drug use, often associated with craving and difficulty controlling use. Dependence is distinct from addiction (substance use disorder), which involves loss of control over use despite negative consequences.

\medterm{Drug-Drug Interactions}-- When one drug alters the pharmacokinetics or pharmacodynamics of another. Pharmacokinetic interactions affect absorption, distribution, metabolism, or excretion (e.g., CYP inhibitors increasing warfarin levels). Pharmacodynamic interactions involve additive, synergistic, or antagonistic effects at the site of action (e.g., two drugs prolonging QT interval). Clinically significant interactions can be identified through databases, clinical experience, and therapeutic drug monitoring.

\medterm{Drug-Induced Liver Injury}-- Liver damage caused by medications, herbal products, or dietary supplements, representing the most common cause of acute liver failure in the United States. DILI can be dose-dependent and predictable, as with acetaminophen hepatotoxicity from overdose, or idiosyncratic and unpredictable, as with isoniazid, nitrofurantoin, and amoxicillin-clavulanate. The pathogenesis involves formation of reactive metabolites that bind to cellular proteins, triggering immune-mediatedimmune-mediated hepatocyte injury. DILI ranges from asymptomatic transaminase elevation to fulminant hepatic failure and death; early recognition and discontinuation of the offending agent are critical, as no specific therapy exists for most forms of DILI.

\medterm{Drug Interactions}-- Drug interactions occur when one drug affects the activity of another drug when both are administered together, potentially leading to therapeutic failure or toxicity. Pharmacokinetic interactions alter the absorption, distribution, metabolism, or excretion of a drug. Absorption interactions include chelation of tetracyclines with calcium and antacids reducing absorption, and altered GI motility affecting drug absorption. Distribution interactions include displacement from protein binding sites,sites, and drug interactions affecting protein binding. Pharmacodynamic interactions, where drugs have additive, synergistic, or antagonistic effects at the same receptor or physiological system, are equally important and can lead to enhanced therapeutic effects or unexpected toxicity.

\medterm{Drug Metabolism Pharmacogenomics}-- Pharmacogenomics of drug metabolism examines how genetic variations affect individual responses to drugs, enabling personalized medicine approaches to prescribing. The most clinically relevant pharmacogenes include CYP2D6, affecting metabolism of approximately 25 percent of drugs including codeine, tamoxifen, and many antidepressants; CYP2C19, affecting clopidogrel activation and proton pump inhibitor metabolism; CYP2C9, affecting warfarin metabolism; TPMT, affecting thiopurine metabolism; and Dand DPYD, affecting fluorouracil metabolism and risk of severe toxicity. Pre-emptive pharmacogenomic testing is increasingly being integrated into routine clinical practice to guide drug selection and dosing, with significant benefits in terms of efficacy and safety.

\medterm{Drug Resistance}-- The ability of microorganisms, cancer cells, or parasites to survive and proliferate despite exposure to drugs that were previously effective. In antimicrobial resistance, mechanisms include enzymatic inactivation of the drug such as beta-lactamase hydrolyzing penicillins, target modification such as PBP2a in MRSA conferring resistance to beta-lactams, efflux pumps actively pumping drugs out of the cell, decreased permeability through porin mutations, and metabolic bypapassing of the inhibited pathway, allowing the pathogen to survive drug exposure; understanding these diverse mechanisms is essential for the rational design of new therapeutic strategies and for guiding appropriate antimicrobial stewardship interventions.

\medterm{Drug Tolerance}-- A decreased pharmacological response to a drug after repeated administration, requiring dose escalation to achieve the same effect. Tolerance develops through multiple mechanisms: pharmacokinetic tolerance — increased drug metabolism through enzyme induction (e.g., alcohol, barbiturates); pharmacodynamic tolerance — receptor downregulation (decreased number of receptors, e.g., opioids, beta-agonists) or receptor desensitization (uncoupling of receptors from signal transduction, e.g., beta-agonists); functional tolerance — compensatory physiological adaptations; and behavioral tolerance — learned adaptations to drug effects. Tolerance is distinct from dependence and addiction. Cross-tolerance occurs when tolerance to one drug confers tolerance to another drug in the same class (e.g., morphine and heroin).

\medterm{Duloxetine}-- A serotonin-norepinephrine reuptake inhibitor (SNRI) used for major depressive disorder, generalized anxiety disorder, diabetic peripheral neuropathy, fibromyalgia, chronic musculoskeletal pain (including osteoarthritis), and stress urinary incontinence. Adverse effects: nausea (common), dry mouth, insomnia, fatigue, hypertension, and sexual dysfunction; withdrawal syndrome with abrupt cessation (taper). Contraindicated with MAOIs and in uncontrolled narrow-angle glaucoma; caution in hepatic impairment.

\medterm{Echinocandins}-- A class of antifungal agents that inhibit beta-1,3-glucan synthase, an enzyme essential for synthesis of beta-1,3-glucan, a major component of the fungal cell wall. This unique mechanism targets a structure absent in mammalian cells, resulting in favorable toxicity profiles. The available agents include caspofungin, micafungin, and anidulafungin, all administered intravenously. Echinocandins are fungicidal against Candida species and fungistatic against Aspergillus species. TheThe three available echinocandins have similar spectra of activity, though minor differences in pharmacokinetics and dosing schedules exist, and they are generally well tolerated with few drug interactions.

\medterm{Efficacy}-- The maximum therapeutic effect a drug or intervention can produce under ideal conditions (e.g., in a clinical trial). Efficacy is determined by the drug's mechanism of action and the receptor-effector system; it is measured by the Emax (maximal effect) on the dose-response curve. A drug with high efficacy produces a greater maximal response than one with low efficacy, regardless of dose. Efficacy is distinct from potency (the dose required to produce a given effect). Clinically, efficacy determines whether a drug can achieve the desired therapeutic outcome; a drug with insufficient efficacy will not be effective even at high doses. Not the same as effectiveness, which reflects real-world performance.

\medterm{Enfuvirtide}-- A fusion inhibitor used in combination antiretroviral therapy for HIV-1 infection in treatment-experienced patients with multidrug-resistant virus. Enfuvirtide is a synthetic peptide that mimics a region of the HIV-1 gp41 transmembrane protein, preventing the conformational change required for fusion of the viral envelope with the host cell membrane. The drug must be administered twice daily by subcutaneous injection, which is the primary limitation to its use. Common adverse effects include injinjection site reactions, and the drug is associated with an increased rate of lymph node enlargement and potential immune reconstitution syndromes when used as part of combination antiretroviral therapy.

\medterm{Enoxaparin}-- A low-molecular-weight heparin inhibiting factor Xa (and to a lesser extent thrombin) via antithrombin. Used for prophylaxis and treatment of deep vein thrombosis and pulmonary embolism, acute coronary syndromes, and as a bridge during anticoagulation interruption. Subcutaneous dosing (weight-based, BID or daily), predictable pharmacokinetics, no routine monitoring generally (anti-Xa monitoring in renal impairment/pregnancy). Adverse effects: bleeding, thrombocytopenia (less HIT than UFH), and injection-site hematoma. Reversed by protamine (partial). Caution in renal impairment (reduce dose); contraindicated with epidural/spinal anesthesia unless timing respected (neuraxial hematoma risk).

\medterm{Entacapone}-- A peripheral COMT (catechol-O-methyltransferase) inhibitor used as an adjunct to levodopa/carbidopa in Parkinson disease, extending levodopa's duration of action and reducing 'wearing-off'. Does not cross the blood-brain barrier. Adverse effects: dyskinesias (due to increased levodopa), nausea, diarrhea, urine discoloration (orange-brown), and hepatotoxicity (rare). Must be given with levodopa. Tolcapone (central COMT inhibitor) is rarely used due to hepatotoxicity.

\medterm{Enterobacteriaceae carbapenem resistance}-- Resistance to carbapenem antibiotics in Enterobacteriaceae, primarily through production of carbapenemases. The most important carbapenemases include KPC, NDM, and OXA-48-like enzymes. Carbapenem-resistant Enterobacteriaceae (CRE) are a major public health threat due to limited treatment options. Infections are associated with high mortality rates. Treatment options include ceftazidime-avibactam, meropenem-vaborbactam, and aminoglycosides. Infection prevention and antimicrobial stewardship are eessential for controlling the spread of carbapenem-resistant Enterobacteriaceae in healthcare settings.

\medterm{Enzyme Induction}-- A process by which a drug or xenobiotic increases the synthesis or activity of drug-metabolizing enzymes, primarily cytochrome P450 enzymes in the liver, accelerating the metabolism of itself (autoinduction) or other drugs. Enzyme induction typically occurs over days to weeks as it requires new protein synthesis. Important enzyme inducers: rifampin (potent CYP3A4, CYP2C9, CYP2C19 inducer), carbamazepine, phenytoin, phenobarbital, St. John's wort, and chronic alcohol consumption. Clinical consequences: reduced efficacy of co-administered drugs (e.g., rifampin reduces oral contraceptive efficacy, requiring alternative contraception), increased metabolism of warfarin (requiring higher doses), and reduced cyclosporine levels in transplant patients. discontinuing an inducer can lead to increased drug levels of co-administered drugs, requiring dose reduction.

\medterm{Enzyme Inhibition}-- A process by which a drug inhibits the activity of drug-metabolizing enzymes, leading to increased plasma concentrations of co-administered drugs that are metabolized by the same enzyme. Enzyme inhibition occurs rapidly (hours to days) as it does not require new protein synthesis. Important enzyme inhibitors and their substrates: ketoconazole and ritonavir (potent CYP3A4 inhibitors — increase levels of statins, benzodiazepines, calcium channel blockers); cimetidine (inhibits CYP1A2, CYP2D6, CYP3A4); fluoxetine and paroxetine (CYP2D6 inhibitors — increase levels of atomoxetine, TCAs, and tamoxifen); amiodarone (CYP2C9 inhibitor — increases warfarin levels); grapefruit juice (inhibits intestinal CYP3A4 — increases levels of felodipine, simvastatin, cyclosporine). Clinical consequences: increased drug toxicity (e.g., rhabdomyolysis with simvastatin when given with potent CYP3A4 inhibitors), requiring dose reduction or alternative therapy.

\medterm{Eplerenone}-- A selective aldosterone (mineralocorticoid) receptor antagonist, a potassium-sparing diuretic. Used to reduce mortality in heart failure with reduced ejection fraction (post-MI and chronic HFrEF, NYHA II-IV) and for resistant hypertension. More selective than spironolactone with less gynecomastia/breast tenderness. Adverse effects: hyperkalemia (major), gynecomastia (less than spironolactone), and hypotension. Contraindicated with significant renal impairment, hyperkalemia, and strong CYP3A4 inhibitors (azole antifungals). Monitor potassium and renal function.

\medterm{Eravacycline}-- A fluorocycline antibiotic with activity against gram-positive and gram-negative organisms, including multidrug-resistant strains. Eravacycline inhibits bacterial protein synthesis by binding to the 30S ribosomal subunit. It has activity against MRSA, VRE, ESBL-producing Enterobacteriaceae, and anaerobes. The drug is indicated for complicated intra-abdominal infections. Eravacycline is administered intravenously and achieves high concentrations in abdominal tissues. Common adverse effects includincluding nausea, vomiting, and infusion site reactions; eravacycline has a favourable safety profile with minimal phototoxicity compared to other tetracyclines.

\medterm{ESBL-producing organisms}-- Bacteria that produce extended-spectrum beta-lactamases (ESBLs), enzymes that hydrolyze penicillins, cephalosporins, and monobactams. ESBLs are common in Enterobacteriaceae, particularly Klebsiella pneumoniae and E. coli. ESBL production confers resistance to most beta-lactam antibiotics except carbapenems. The prevalence of ESBL-producing organisms is increasing globally. Infections are associated with increased morbidity, mortality, and healthcare costs. Treatment typically requires carbapenemcarbapenems for serious infections, and antimicrobial stewardship is essential to preserve the utility of remaining effective agents.

\medterm{Escitalopram}-- A selective serotonin reuptake inhibitor (SSRI), the S-enantiomer of citalopram, used for major depressive disorder and generalized anxiety disorder. Increases synaptic serotonin; well tolerated, low drug-interaction profile (less QT prolongation than citalopram). Adverse effects: nausea, insomnia or somnolence, sexual dysfunction, and weight change; serotonin syndrome and discontinuation syndrome (taper); black-box warning for suicidal ideation in young adults. Once-daily dosing; onset of effect 2-6 weeks.

\medterm{Esomeprazole}-- The S-enantiomer of omeprazole, a proton pump inhibitor irreversibly inhibiting gastric H+/K+ ATPase. Used for GERD, erosive esophagitis, peptic ulcer disease, H. pylori eradication, and hypersecretory states. Greater acid suppression than omeprazole at equivalent doses; once daily, with an IV formulation for acute GI bleeding. Adverse effects: headache, diarrhea, hypomagnesemia, B12 deficiency, fracture risk, and C. difficile infection with chronic use. Cytochrome P450 (CYP2C19) metabolism.

\medterm{Essential Medicines}-- Drugs that satisfy the priority healthcare needs of a population, as defined by the WHO Model List of Essential Medicines. Selection criteria include efficacy, safety, comparative cost-effectiveness, and suitability for use in resource-limited settings. The list includes approximately 400 drugs covering major disease areas: analgesics, anesthetics, antibiotics, antiepileptics, antihypertensives, antiretrovirals, antimalarials, vaccines, and others. The concept guides national formulary development, procurement, and rational prescribing worldwide.

\medterm{Etravirine}-- A second-generation non-nucleoside reverse transcriptase inhibitor (NNRTI) used in combination antiretroviral therapy for HIV-1 infection, particularly in treatment-experienced patients with NNRTI resistance. Etravirine binds to the HIV reverse transcriptase enzyme and inhibits its activity through a unique mechanism that maintains activity against strains resistant to first-generation NNRTIs. The drug has a high genetic barrier to resistance due to its flexible structure that allows it to adaptadapt to changes in the reverse transcriptase enzyme structure, maintaining antiviral activity against a broad range of NNRTI-resistant variants.

\medterm{Expectorants}-- Medications that promote the clearance of mucus from the respiratory tract by increasing sputum volume and reducing its viscosity, facilitating productive cough. The most common expectorant is guaifenesin (glyceryl guaiacolate), which increases respiratory tract fluid secretion and reduces sputum adhesiveness. Guaifenesin is used for symptom relief in acute respiratory infections, chronic bronchitis, and COPD. Evidence for efficacy is moderate. Other expectorants: potassium iodide (rarely used due to side effects), ammonium chloride, and sodium citrate. Expectorants are distinguished from mucolytics (acetylcysteine, dornase alfa--break down the chemical structure of mucus) and antitussives (dextromethorphan, codeine--suppress the cough reflex). Adequate hydration is essential for expectorant efficacy.

\medterm{Ezetimibe}-- A cholesterol absorption inhibitor that blocks intestinal absorption of dietary and biliary cholesterol at the Niemann-Pick C1-like 1 (NPC1L1) transporter in the small intestine. Used for hypercholesterolemia, usually as add-on to statins or in statin-intolerant patients; reduces LDL and cardiovascular events in combination with statins after ACS. Adverse effects: diarrhea, myalgia, transaminase elevation (with statins). Rare reports of muscle toxicity. Available alone or in combination with simvastatin or atorvastatin.

\medterm{Fecal microbiota transplantation}-- The transfer of fecal material from a healthy donor to a patient with recurrent Clostridioides difficile infection. The procedure aims to restore normal gastrointestinal flora and eliminate C. difficile colonization. FMT is highly effective, with cure rates exceeding 90\% for recurrent C. difficile infection. The procedure can be administered by colonoscopy, nasojejunal tube, or oral capsules. Donor screening includes testing for bloodborne pathogens and gastrointestinal pathogens. FMT has revolurevolutionised the management of recurrent C. difficile infection and is being investigated for other conditions associated with intestinal dysbiosis, including inflammatory bowel disease and metabolic syndrome.

\medterm{Fenofibrate}-- A fibric acid derivative (fibrate) activating PPAR-alpha, increasing HDL and reducing triglycerides, with modest LDL effects. Used for severe hypertriglyceridemia (pancreatitis risk), mixed dyslipidemia, and as an adjunct in residual hypertriglyceridemia. Adverse effects: myopathy (especially with statins, more with gemfibrozil), transaminase and creatinine elevation, gallstones, and pancreatitis. Dose-adjust in renal impairment. Fibrates are avoided with severe hepatic disease and in patients on other fibrates.

\medterm{Fexofenadine}-- A second-generation (non-sedating) H1 antihistamine used for allergic rhinitis and chronic urticaria. Peripherally selective with minimal sedation and anticholinergic effects; no significant CNS penetration. Adverse effects: headache, dizziness, and dry mouth. Transported by P-glycoprotein—avoid concurrent aluminum/magnesium antacids and fruit juices (reduce absorption). Available OTC. Loratadine and cetirizine are alternatives.

\medterm{First-Pass Effect}-- The metabolism of a drug by the liver and intestinal wall before it reaches the systemic circulation after oral administration. When a drug is absorbed from the gastrointestinal tract, it travels through the portal venous system to the liver, where it may be extensively metabolized by cytochrome P450 enzymes before reaching the general circulation. This process significantly reduces the bioavailability of many drugs. The hephepatic extraction ratio quantifies this effect: drugs with high extraction ratios have clearance that depends primarily on hepatic blood flow, while low-extraction drugs have clearance that depends on intrinsic enzyme activity and protein binding; this distinction has important implications for drug dosing in patients with liver disease and for predicting drug interactions.

\medterm{Flecainide}-- A class IC sodium channel blocker that slows conduction velocity with little effect on refractoriness. Used for suppression of supraventricular arrhythmias (atrial fibrillation in patients without structural heart disease—'pill-in-the-pocket', paroxysmal SVT) and ventricular arrhythmias. Contraindicated in ischemic heart disease and structural heart disease (CAST trial showed increased mortality after MI) and in significant LV dysfunction. Adverse effects: proarrhythmia (new/worsened arrhythmia), dizziness, blurred vision, and heart block. Requires ECG/electrolyte monitoring.

\medterm{Fluoroquinolones}-- Broad-spectrum bactericidal antibiotics that inhibit bacterial DNA gyrase and topoisomerase IV, enzymes essential for DNA replication, transcription, and repair. They are synthetic derivatives of nalidixic acid and include ciprofloxacin, levofloxacin, moxifloxacin, and others. Fluoroquinolones have excellent bioavailability, with oral dosing achieving serum concentrations comparable to intravenous administration. They distribute widely to tissues including the lungs, prostatprostate, and bone; fluoroquinolones are generally well tolerated but carry boxed warnings for tendonitis and tendon rupture, peripheral neuropathy, CNS effects, and exacerbation of myasthenia gravis, and their use is reserved for situations where alternative treatment options are limited.

\medterm{Fluoxetine}-- A selective serotonin reuptake inhibitor (SSRI) used for major depressive disorder, obsessive-compulsive disorder, bulimia, panic disorder, and premenstrual dysphoric disorder. Long half-life (fluoxetine ~1-3 days; active metabolite norfluoxetine ~7-15 days) allows once-daily or weekly dosing and produces a washout period (advantageous in discontinuation, problematic in toxicity/interactions). Adverse effects: nausea, insomnia, sexual dysfunction, agitation, hyponatremia (elderly), and serotonin syndrome. Strong CYP2D6 inhibitor—many interactions. Black-box warning for suicidality in young adults.

\medterm{Fondaparinux}-- A synthetic pentasaccharide that selectively inhibits factor Xa via antithrombin, with no direct thrombin inhibition. Used for prophylaxis and treatment of venous thromboembolism and acute coronary syndromes (as an alternative to LMWH). Subcutaneous, once daily, predictable, no routine monitoring; does not cause HIT (no platelet factor 4 interaction). Adverse effects: bleeding; not reversed by protamine. Contraindicated in severe renal impairment (creatinine clearance <30) and low body weight. Neuraxial anesthesia requires timing precautions.

\medterm{Full Agonist}-- A full agonist is a drug that binds to a receptor and produces the maximum possible biological response that the receptor can generate, achieving 100 percent of the maximal effect. Full agonists have both high affinity for the receptor and high intrinsic activity, meaning they are very effective at activating the receptor once bound. The efficacy of a full agonist is inherent to the drug-receptor interaction and cannot be exceeded by increasing the dose. Examples include morphine at mu-opioid rereceptors, with full agonists producing maximal activation, and their effect is determined by both the drug concentration at the receptor site and the receptor density in the target tissue.

\medterm{Fungal Infection}-- See Antifungal. Infections caused by fungi ranging from superficial (dermatophytes: tinea, candidiasis) to invasive/systemic disease (aspergillosis, candidemia, cryptococcosis, histoplasmosis) in immunocompromised hosts. Risk factors: neutropenia, transplantation, HIV/AIDS, corticosteroids, and diabetes. Diagnosis: culture, histopathology, antigen tests (galactomannan, beta-glucan, cryptococcal antigen), and PCR. Treatment depends on severity and organism: azoles, echinocandins, amphotericin B, and terbinafine.

\medterm{Fungal resistance mechanisms}-- Mechanisms by which fungi develop resistance to antifungal agents. Azole resistance occurs through upregulation of efflux pumps (CDR1, CDR2, MDR1), mutations in the target enzyme ERG11, and overexpression of ERG11. Echinocandin resistance is primarily due to mutations in the FKS1 and FKS2 genes encoding beta-(1,3)-D-glucan synthase. Amphotericin B resistance is rare but can occur through alterations in ergosterol biosynthesis. Multidrug resistance is increasingly common in Candida glabrata, withwith azole resistance being the most clinically challenging, particularly in Candida glabrata and Candida auris, and susceptibility testing is increasingly important for guiding antifungal therapy.

\medterm{Furosemide}-- A potent high-ceiling loop diuretic that reversibly inhibits the $Na^+/K^+/2Cl^-$ cotransporter (NKCC2) in the thick ascending limb of the loop of Henle. Induces rapid diuresis and natriuresis by disrupting the hypertonic medullary osmotic gradient. Indications: acute pulmonary edema, acute decompensated heart failure, edematous states (cirrhosis, nephrotic syndrome), hypercalcemia (increases urinary $Ca^{2+}$ excretion), and hypertension. Adverse effects (mnemonic OH DANG): Ototoxicity (sensorineural hearing loss/tinnitus), Hypokalemia/Hypomagnesemia/Hypocalcemia, Dehydration/Hypotension, Allergy (sulfa drug reaction), Nephrotoxicity (interstitial nephritis), Gout (hyperuricemia due to uric acid reabsorption).

\medterm{Gamma-Globulin}-- A fraction of plasma proteins containing antibodies (immunoglobulins). Historically the term referred to a group of serum proteins with the slowest electrophoretic mobility (gamma region on serum protein electrophoresis). Currently synonymous with immunoglobulin preparations used for passive immunisation and immunotherapy. Indications: (1) Replacement therapy for primary immunodeficiency (X-linked agammaglobulinaemia, common variable immunodeficiency); (2) Immunomodulation for autoimmune and inflammatory diseases (Kawasaki disease, Guillain-Barré syndrome, chronic inflammatory demyelinating polyneuropathy, ITP, myasthenia gravis, Stevens-Johnson syndrome). Administration: intravenous (IVIG) or subcutaneous (SCIG). Side effects: headache, fever, chills, myalgia, aseptic meningitis, thrombosis, anaphylaxis (rare, in IgA deficiency). Prepared by cold ethanol fractionation (Cohn method) from pooled human plasma from thousands of donors.

\medterm{Gel}-- A semisolid pharmaceutical formulation in which a liquid (usually water) is immobilized within a three-dimensional network of a gelling agent (e.g., carbomer, cellulose, alginate). Gels provide controlled drug release, cool sensation, and good skin compatibility; they are used for topical dermatologic products (antibiotics, antifungals, NSAIDs), ocular preparations, and oral gels. They are generally clear, non-greasy, and easily washed off.

\medterm{Gelfoam}-- An absorbable hemostatic agent made from purified porcine gelatin. Used in surgical settings to control bleeding where conventional hemostasis is difficult (liver, bone, neurosurgery). Mechanism: provides a physical matrix that promotes platelet aggregation and clot formation. Absorbed within 4-6 weeks. Prepared as a sterile sponge that can be cut to size and applied dry or saturated with thrombin. Contraindicated: infected wounds, gross contamination, intravascular use (risk of embolism). Brand examples: Gelfoam (Pfizer), Surgifoam (Ethicon).

\medterm{Gemfibrozil}-- A fibrate (PPAR-alpha agonist) used primarily for severe hypertriglyceridemia (reducing pancreatitis risk) and mixed hyperlipidemia. Raises HDL and lowers triglycerides markedly; modest LDL lowering. Adverse effects: myopathy/rhabdomyolysis—significantly increased with statin coadministration (interaction via glucuronidation inhibition; contraindicated with many statins), GI upset, gallstones, and increased INR with warfarin. Contrast with fenofibrate, which has less drug interaction risk. Not first-line for isolated hypercholesterolemia.

\medterm{General Anesthetic}-- An agent producing reversible loss of consciousness, amnesia, analgesia, and muscle relaxation sufficient for surgery. Inhaled agents (sevoflurane, isoflurane, desflurane, nitrous oxide) and intravenous agents (propofol, thiopental, etomidate, ketamine) are used; balanced anesthesia combines them with opioids, neuromuscular blockers, and adjuvants. Monitored by depth-of-anesthesia, hemodynamics, and end-tidal agent concentration. Complications: respiratory depression, hypotension, nausea/vomiting, malignant hyperthermia (succinylcholine/volatile agents), and postoperative cognitive dysfunction.

\medterm{GLP-1 Receptor Agonists}-- Glucagon-like peptide-1 receptor agonists are a class of injectable antidiabetic agents that mimic the effects of the incretin hormone GLP-1, which is secreted by intestinal L-cells in response to food intake. GLP-1 stimulates glucose-dependent insulin secretion, suppresses glucagon secretion, delays gastric emptying, and promotes satiety through central nervous system effects. Available agents include exenatide, liraglutide, dulaglutide, semaglutide, and tirzepatide, which is a dual GIP/GLP-1 aagonist. GLP-1 receptor agonists improve glycaemic control with a low risk of hypoglycaemia and promote weight loss, making them particularly valuable in overweight or obese patients with type 2 diabetes; gastrointestinal side effects including nausea, vomiting, and diarrhea are common but often diminish with gradual dose titration.

\medterm{Glycerin}-- A colorless, odorless, sweet-tasting trihydroxy alcohol (glycerol) with humectant, osmotic, and solvent properties. Used as: an osmotic laxative (glycerin suppositories or enemas promote water retention and reflex defecation), a humectant in skin preparations, a vehicle/solvent in pharmaceuticals and cosmetics, and in otolaryngology (glycerin ear drops soften cerumen). Intravenous glycerol is used to reduce intracranial pressure. Generally safe; rectal use may irritate mucosa.

\medterm{Glycerol (Glycerine)}-- A triol compound (CH$_2$OH-CHOH-CH$_2$OH), a colourless, odourless, viscous liquid with a sweet taste. Used in medicine and pharmaceuticals as a: humectant (retains moisture in topical preparations), laxative (glycerol suppositories for constipation), sweetener (in syrups, sugar-free alternatives), solvent for drugs, and osmotherapy agent (intravenous glycerol reduces intracranial pressure). Glycerol is also a by-product of soap manufacturing and biodiesel production (transesterification).

\medterm{GnRH Agonist}-- Gonadotropin-releasing hormone agonists are a class of medications that initially stimulate (flare effect, 1-3 weeks) and then suppress pituitary gonadotropin (LH and FSH) secretion through receptor downregulation, leading to profound suppression of sex steroid production (testosterone in men, estradiol in women). Used in: prostate cancer — medical castration (leuprolide, goserelin, triptorelin); breast cancer — ovarian suppression in premenopausal women; endometriosis and uterine fibroids — suppression of ovarian estrogen; central precocious puberty — halting pubertal progression; and assisted reproduction — controlled ovarian hyperstimulation protocols. The initial testosterone flare in prostate cancer can cause disease exacerbation (bone pain, spinal cord compression, urinary obstruction), mitigated by adding an antiandrogen (bicalutamide) for 2-4 weeks at initiation.

\medterm{GnRH Antagonist}-- A class of medications that competitively block GnRH receptors on pituitary gonadotrophs, producing immediate, reversible suppression of LH and FSH secretion without the initial flare effect seen with GnRH agonists. Used in: prostate cancer — degarelix (monthly subcutaneous injection) and relugolix (oral daily) provide rapid testosterone suppression without flare; endometriosis and uterine fibroids — elagolix (oral); assisted reproduction — cetrorelix and ganirelix prevent premature LH surge during controlled ovarian stimulation. Advantage over GnRH agonists: no initial flare, faster onset of suppression, and reversibility. Adverse effects: hot flashes, injection site reactions (degarelix), and osteoporosis with long-term use.

\medterm{Guaifenesin}-- An expectorant that thins and loosens bronchial mucus by increasing respiratory tract fluid and reducing mucus viscosity, facilitating productive cough. Used for symptomatic relief of chest congestion and productive cough in upper respiratory infections and COPD. Typical adult dose 200-400 mg every 4 hours. Generally well tolerated; common side effects include nausea, vomiting, and dizziness. Efficacy is debated; hydration is an important adjunct. Available in many combination cold medications.

\medterm{H2 Blocker (Histamine-2 Receptor Antagonist)}-- A drug that blocks histamine H2 receptors on gastric parietal cells, reducing gastric acid secretion. Agents: cimetidine, ranitidine, famotidine, nizatidine. Used for GERD, peptic ulcer disease, and dyspepsia—less potent than PPIs. Adverse effects: headache, diarrhea, confusion (elderly, cimetidine), and rare thrombocytopenia. Cimetidine inhibits CYP450 (warfarin, theophylline, phenytoin) and has antiandrogenic effects (gynecomastia). Famotidine has the fewest interactions. Available OTC.

\medterm{Half-Life}-- Elimination half-life is the time required for the plasma concentration of a drug to decrease by 50 percent after reaching equilibrium. It is calculated from the elimination rate constant using the equation t1/2 equals 0.693 divided by the elimination rate constant. Half-life depends on both the volume of distribution and the clearance of the drug, with the relationship t1/2 equals 0.693 times Vd divided by clearance. Drugs with large volumes of distribution or low clearance have long half-liveslives; understanding half-life is essential for determining appropriate dosing intervals, time to reach steady state, and duration of drug action after discontinuation.

\medterm{Histamine}-- A biogenic amine (imidazolylethylamine) synthesised from L-histidine by histidine decarboxylase, stored in basophils, mast cells, and platelets, and released during allergic and inflammatory responses. Histamine acts via four G-protein-coupled receptors: (1) H1 receptor: bronchoconstriction (asthma), vasodilation (hypotension, flushing), increased vascular permeability (urticaria, angioedema), prpruritus, pain, mucus secretion, and CNS effects (wakefulness). H1-antihistamines (cetirizine, loratadine, fexofenadine, chlorphenamine, promethazine, diphenhydramine) are used for allergic rhinitis, urticaria, prpruritus, insomnia, motion sickness. (2) H2 receptor: gastric acid secretion (parietal cells). H2-antagonists (cimetidine, ranitidine, famotidine) are used for peptic ulcer disease, GORD, Zollinger-Ellison syndrome. (3) H3 receptor: pre-synaptic auto- and heteroreceptor in the CNS, inhibits histamine release and synthesis; inverse agonists/antagonists (pitolisant) are used for narcolepsy (increases wakefulness). (4) H4 receptor: expressed on immune cells (eosinophils, neutrophils, mast cells, dendritic cells), involved in chemotaxis and inflammation; H4 antagonists are under investigation for prpruritus and asthma. Degradation: histamine N-methyltransferase and diamine oxidase. Histamine is also involved in gastric acid secretion, neurotransmission (wakefulness, appetite, learning), and smooth muscle contraction. Histamine poisoning (scombroid poisoning): from spoiled fish (tuna, mackerel) containing high levels of histamine due to bacterial decarboxylation, causes flushing, headache, urticaria, diarrhea within minutes, treated with H1-antihistamines.

\medterm{Hydralazine}-- A direct arteriolar vasodilator used for hypertension, particularly in preeclampsia/eclampsia (IV), and in combination with isosorbide dinitrate for heart failure with reduced ejection fraction in African Americans (improves survival). Mechanism: smooth muscle relaxation (nitric oxide pathway), reducing afterload. Adverse effects: reflex tachycardia (coadminister beta-blocker), headache, flushing, fluid retention (with diuretic), and drug-induced lupus with long-term high doses. Rapid acetylators have reduced effect. Also available orally (hypertension add-on).

\medterm{Hydrochlorothiazide}-- A thiazide diuretic inhibiting the sodium-chloride cotransporter in the distal convoluted tubule, promoting natriuresis and reducing plasma volume, with long-term vasodilation. Used for hypertension and mild edema. Lowers blood pressure effectively; useful in calcium stones and osteoporosis (reduces calciuria). Adverse effects: hypokalemia, hyponatremia, hyperglycemia, hyperuricemia (gout), hypercalcemia, hyperlipidemia, and photosensitivity. Less potent than chlorthalidone. Contraindicated in sulfonamide allergy and anuria.

\medterm{Hypnotic}-- A medication that induces, maintains, or improves sleep, used for insomnia. Major classes: benzodiazepine receptor agonists (Z-drugs: zolpidem, zaleplon, eszopiclone) — selective for GABA-A receptors containing alpha-1 subunits, promoting sleep onset with fewer next-day effects than benzodiazepines; melatonin receptor agonists (ramelteon, prolonged-release melatonin) — activate MT1 and MT2 receptors in the suprachiasmatic nucleus, regulating the sleep-wake cycle, no abuse potential; orexin receptor antagonists (suvorexant, daridorexant) — block orexin signaling that promotes wakefulness, used for sleep maintenance insomnia; sedating antidepressants (trazodone, doxepin, mirtazapine) — used in low doses for insomnia, especially with comorbid depression; and antihistamines (diphenhydramine, doxylamine) — sedating H1 antagonists available OTC, limited efficacy with chronic use due to tolerance.

\medterm{Idiosyncratic Reaction}-- An uncommon, unexpected, and often genetically determined adverse drug reaction that is unrelated to the drug's known pharmacology and usually dose-independent (a Type B reaction). Examples: sulfonamide-induced Stevens-Johnson syndrome, abacavir hypersensitivity (HLA-B*5701), isoniazid-induced hepatitis (slow acetylators), and malignant hyperthermia (ryanodine receptor mutation). Idiosyncratic reactions are unpredictable, potentially severe, and necessitate drug discontinuation; some can be prevented by pharmacogenetic screening.

\medterm{Immunosuppressant}-- A drug that reduces immune system activity, used to prevent organ transplant rejection and to treat autoimmune and inflammatory diseases. Agents: calcineurin inhibitors (cyclosporine, tacrolimus), antiproliferatives (azathioprine, mycophenolate), corticosteroids, mTOR inhibitors (sirolimus), and biologic agents (anti-TNF, anti-IL-6, rituximab). Adverse effects include increased infection risk, malignancy (especially lymphoma and skin cancer), nephrotoxicity (calcineurin inhibitors), and bone marrow suppression. Dose is balanced between efficacy and infection risk.

\medterm{Inotrope}-- A drug that increases myocardial contractility (positive inotropy), used for cardiogenic shock, decompensated heart failure, and perioperative low-output states. Agents: dobutamine (beta-1 agonist), milrinone (PDE3 inhibitor), dopamine, epinephrine, norepinephrine (also vasopressors), and digoxin (chronic). Inotropes increase cardiac output and oxygen delivery but may worsen ischemia and arrhythmias and increase mortality with prolonged use. Vasopressors (norepinephrine, vasopressin, phenylephrine) raise blood pressure via vasoconstriction. Choice depends on hemodynamic profile (wet vs dry, cold vs warm).

\medterm{Insulin Aspart}-- A rapid-acting insulin analog (modified amino acid sequence) with onset 10-20 minutes, peak 1-3 hours, duration 3-5 hours. Used to cover prandial glucose (injected at or just after meals) in type 1 and type 2 diabetes; also used in continuous subcutaneous insulin infusion (pump) and IV in acute settings. Adverse effects: hypoglycemia, weight gain, and injection-site lipodystrophy (rotate sites). Lispro and glulisine are similar rapid-acting analogs.

\medterm{Insulin Glargine}-- A long-acting basal insulin analog with a flat, peakless profile providing near-24-hour coverage, administered once daily (or twice in some patients). Formed by modification of human insulin producing slow, steady absorption from the injection depot. Used for type 1 and type 2 diabetes (basal component). Adverse effects: hypoglycemia, weight gain, lipodystrophy. Contrast with intermediate (NPH) and short (regular) insulins. U-100 and U-300 (longer duration) strengths available.

\medterm{Intramuscular (IM)}-- A route of drug administration in which the drug is injected into skeletal muscle (deltoid, vastus lateralis, gluteus maximus). IM injection provides rapid absorption (muscle is highly vascular), avoids first-pass metabolism, and is used for vaccines, antibiotics (benzathine penicillin), depot preparations (medroxyprogesterone), and medications poorly absorbed orally. Absorption rate depends on blood flow and drug lipophilicity. Risks: pain, hematoma, nerve injury (sciatic), abscess, and injection into vessels.

\medterm{Isavuconazole}-- A triazole antifungal agent used in the treatment of invasive aspergillosis and mucormycosis. Isavuconazole inhibits fungal lanosterol 14-alpha demethylase, disrupting ergosterol synthesis and fungal cell membrane integrity. A key advantage is its broad-spectrum activity covering Aspergillus, Candida, Cryptococcus, Mucorales, and endemic fungi. The drug has excellent oral bioavailability of approximately 98\% and is available as both intravenous and oral formulations. Unlike voriconazole, isavucoisavucoazole has minimal effect on CYP3A4 and does not prolong the QT interval to the same degree as voriconazole, giving it a more favourable drug interaction profile.

\medterm{Isosorbide Dinitrate}-- An organic nitrate (and prodrug) that releases nitric oxide, relaxing vascular smooth muscle—predominantly venous (preload reduction) with coronary dilation. Used for chronic stable angina prophylaxis, acute angina (sublingual/IV forms), and, with hydralazine, heart failure with reduced ejection fraction (improves survival in African Americans). Long-term continuous use produces nitrate tolerance (requires nitrate-free interval). Adverse effects: headache, flushing, hypotension, syncope, reflex tachycardia. Contraindicated with PDE5 inhibitors (sildenafil, tadalafil).

\medterm{Laxative}-- A drug that promotes bowel evacuation, used for constipation and bowel preparation. Classes: bulk-forming (psyllium, methylcellulose), osmotic (polyethylene glycol, lactulose, magnesium salts), stimulant (senna, bisacodyl), stool softener (docusate), and lubricant (mineral oil). Choice depends on constipation type and patient factors. Chronic stimulant laxative misuse can cause electrolyte disturbances and colonic dysmotility; bulk and osmotic agents are generally preferred for chronic use. Laxatives are contraindicated in intestinal obstruction.

\medterm{Levodopa}-- The precursor of dopamine that crosses the blood-brain barrier; the most effective symptomatic treatment for Parkinson disease. Converted to dopamine in the brain by aromatic amino acid decarboxylase; always combined with carbidopa (peripheral decarboxylase inhibitor) to prevent peripheral conversion and side effects. Acute benefits: tremor, rigidity, bradykinesia. Long-term complications: motor fluctuations ('wearing off', 'on-off'), dyskinesias, and psychiatric effects. Withheld in psychosis; caution with MAOIs.

\medterm{Liposomal amphotericin B}-- A liposomal formulation of amphotericin B used in the treatment of severe systemic fungal infections, with reduced nephrotoxicity compared to conventional amphotericin B. The liposomal formulation encapsulates amphotericin B within lipid vesicles, which alters the drug's biodistribution and reduces renal tubular damage. It is the preferred formulation for neutropenic patients with fever and suspected fungal infections. The drug has broad-spectrum activity against Candida, Aspergillus, CryptococcAdditionally, liposomal amphotericin B is generally better tolerated with fewer infusion-related reactions and lower rates of nephrotoxicity, allowing for higher doses to be administered with improved therapeutic outcomes in critically ill patients with invasive fungal infections.

\medterm{Loading Dose}-- A higher initial dose of a drug administered to rapidly achieve therapeutic plasma concentrations, followed by lower maintenance doses to sustain the steady state. The loading dose is calculated as: Loading dose = (Vd x desired concentration) / bioavailability. Indications: drugs with long half-lives where waiting 4-5 half-lives to reach steady state would delay clinical effect (amiodarone, half-life 40-55 days; digoxin, half-life 36-48 hours); drugs requiring immediate therapeutic effect (phenytoin for status epilepticus, antibiotics in sepsis, antiarrhythmics in unstable arrhythmias). The loading dose does not depend on clearance, only on the volume of distribution and the desired concentration. Caution is required for drugs with narrow therapeutic indices to avoid toxicity from the high initial dose.

\medterm{Loop Diuretics}-- A class of potent diuretics that act on the thick ascending limb of the loop of Henle by inhibiting the sodium-potassium-chloride cotransporter, blocking sodium and chloride reabsorption and producing a powerful natriuresis and diuresis. The prototypical loop diuretic is furosemide, with other agents including bumetanide and torsemide. Loop diuretics are the most potent diuretics available, producing greater sodium excretion than thiazides. They also increase calcium and magneLoop diuretics are used in the treatment of hypertension, heart failure, and oedematous states including cirrhosis and nephrotic syndrome, and their dosing must be individualised based on renal function and clinical response, with careful monitoring of electrolytes and volume status to avoid hypokalaemia, hypomagnesaemia, and volume depletion.

\medterm{Loperamide}-- A peripheral mu-opioid receptor agonist used as an antidiarrheal. It slows intestinal motility, reduces stool frequency and volume, and increases intestinal transit time without significant CNS effects (does not cross the blood-brain barrier at therapeutic doses) or analgesic activity. Adult dose: 4 mg initially, then 2 mg after each loose stool (max 16 mg/day). Side effects: constipation, abdominal cramps, nausea. Contraindicated in invasive bacterial diarrhea (dysentery, C. difficile) and intestinal obstruction; high doses are abused for opioid-like effects and can cause cardiac toxicity (QT prolongation).

\medterm{Loratadine}-- A second-generation (non-sedating) H1 antihistamine used for allergic rhinitis and chronic urticaria. Peripherally selective—minimal CNS penetration, minimal anticholinergic effects. Once daily; prodrug (desloratadine is active). Adverse effects: headache, dry mouth, and rare sedation. Available OTC. Few drug interactions; metabolized by CYP3A4/2D6. Cetirizine and fexofenadine are alternatives.

\medterm{Macrolides}-- A class of bacteriostatic antibiotics that bind to the 50S ribosomal subunit and inhibit translocation of peptidyl-tRNA during protein synthesis. The major macrolides include erythromycin, clarithromycin, and azithromycin. Erythromycin has a narrow spectrum primarily against gram-positive organisms and is limited by gastrointestinal side effects from motilin receptor agonism. Clarithromycin has improved oral bioavailability and extended spectrum including Haemophilus influenzae anand azithromycin has become a widely used agent for respiratory infections due to its extended half-life and once-daily dosing. Macrolides also have immunomodulatory and anti-inflammatory properties independent of their antimicrobial activity, which contribute to their efficacy in chronic respiratory diseases including cystic fibrosis, bronchiectasis, and diffuse panbronchiolitis.

\medterm{Maintenance Dose}-- The dose of a drug administered at regular intervals to maintain a desired steady-state concentration, balancing the rate of drug administration with the rate of elimination. The maintenance dose is calculated as: Maintenance dose rate = (Clearance x desired concentration) / Bioavailability. At steady state, the rate of drug administration equals the rate of elimination. Maintenance doses are determined by the drug's clearance and target therapeutic concentration. Factors affecting maintenance dose requirements: changes in clearance (renal impairment, hepatic impairment, drug interactions affecting metabolism), changes in volume of distribution (obesity, edema), and patient factors (age, weight, organ function). Maintenance dosing intervals are based on the drug's half-life, typically administered every half-life for drugs with short half-lives or at longer intervals for drugs with long half-lives.

\medterm{Mannitol}-- An osmotic diuretic that is freely filtered at the glomerulus but not reabsorbed, increasing tubular osmotic pressure and water excretion. Used to reduce intracranial pressure (cerebral edema) and intraocular pressure (acute glaucoma), and to promote diuresis in rhabdomyolysis/contrast nephropathy prophylaxis. IV infusion; side effects: fluid overload, hyponatremia or hypernatremia, pulmonary edema (with renal failure), and headache. Contraindicated in anuria and congestive heart failure. Filter for crystallization; avoid in significant renal impairment.

\medterm{Maraviroc}-- A CCR5 antagonist (entry inhibitor) used in combination antiretroviral therapy for HIV-1 infection in patients with CCR5-tropic virus. Maraviroc binds to the human CCR5 coreceptor on CD4+ T cells, preventing HIV-1 gp120 binding and viral entry. A critical requirement for maraviroc use is tropism testing to confirm the presence of CCR5-tropic virus, as the drug has no activity against CXCR4-tropic or dual/mixed-tropic virus. Common adverse effects include hepatotoxicity with elevated transaminases levels, which is typically reversible upon discontinuation; postural hypotension, diarrhea, nausea, and headache are also reported.

\medterm{Maximum Tolerated Dose}-- The highest dose of a drug that produces acceptable, manageable toxicity without causing unacceptable adverse effects, determined through dose-escalation studies, most commonly in Phase I clinical trials. The MTD identifies the dose range for subsequent efficacy studies. In oncology, the MTD is typically defined as the dose at which no more than 33\% of patients experience dose-limiting toxicity (DLT) during the first treatment cycle. The 3+3 design is the classic approach: 3 patients are treated at each dose level; if 0/3 have DLT, the next cohort is escalated; if 1/3 have DLT, 3 more patients are treated at the same level; if 2/3 have DLT, the previous dose is declared the MTD. Modern adaptive designs (e.g., Bayesian continual reassessment method) are increasingly used.

\medterm{Medical Cannabis and Cannabidiol (CBD)}-- Cannabis (marijuana) and cannabidiol (CBD) are cannabinoid compounds derived from \emph{Cannabis sativa} that interact with the endocannabinoid system (CB1 and CB2 receptors) to modulate pain, appetite, mood, and inflammation. THC (delta-9-tetrahydrocannabinol) is the primary psychoactive component responsible for euphoria, cognitive impairment, and dependence risk; CBD is non-intoxicating and has demonstrated anti-inflammatory, anxiolytic, and anticonvulsant properties. FDA-approved cannabinoid preparations: dronabinol and nabilone for chemotherapy-induced nausea and HIV/AIDS anorexia, and purified CBD (Epidiolex) for Dravet and Lennox-Gastaut syndromes. Evidence supports medical cannabis for chronic neuropathic pain, multiple sclerosis spasticity, and chemotherapy-induced nausea, though data for most other conditions are limited by small trials and regulatory barriers. Risks: impaired driving, psychosis vulnerability (especially high-THC products), cannabinoid hyperemesis syndrome, pulmonary harm from smoking, and addiction (~9\% of users).

\medterm{Medication Management}-- The systematic process of optimizing pharmacotherapy to achieve therapeutic goals while minimizing adverse effects, encompassing prescribing, dispensing, administration, monitoring, and deprescribing. Key components: medication reconciliation at transitions of care (comparing patient's current regimen against new orders to prevent errors), adherence assessment (pill counts, pharmacy refill data, electronic monitoring), dose adjustment for renal/hepatic impairment, therapeutic drug monitoring (e.g., warfarin INR, vancomycin trough, lithium levels), and deprescribing (tapering or discontinuing unnecessary or harmful medications). Polypharmacy—use of $\ge$5 medications—increases risk of adverse drug events, falls, cognitive impairment, and hospitalization; an annual medication review using tools like the Beers Criteria and STOPP/START criteria is recommended for older adults. Patient education on indication, dose, timing, side effects, and drug--drug interactions is essential for safe and effective therapy.

\medterm{Medications and Treatments}-- all pharmacologic and therapeutic interventions used to prevent, manage, or cure disease, classified by mechanism of action, therapeutic class, and route of administration. Major drug classes: analgesics (NSAIDs, opioids, acetaminophen), cardiovascular agents (statins, ACE inhibitors, beta-blockers, diuretics, anticoagulants), antimicrobials (antibiotics, antivirals, antifungals), endocrine therapies (insulin, metformin, thyroid hormone, corticosteroids), psychotropics (SSRIs, antipsychotics, benzodiazepines, mood stabilizers), antineoplastics (chemotherapy, targeted therapy, immunotherapy), and disease-modifying biologic agents. Routes include oral, intravenous, intramuscular, subcutaneous, topical, inhaled, transdermal, intrathecal, and intra-articular. The principles of pharmacotherapy: right drug, right dose, right route, right time, right patient—guided by evidence-based guidelines, contraindications, allergies, renal/hepatic function, and cost-effectiveness.

\medterm{Medication Side Effects}-- Unintended harmful responses to medications occurring at standard doses, ranging from mild (nausea, rash, dry mouth) to severe (anaphylaxis, Stevens-Johnson syndrome, QT prolongation, hepatotoxicity). ADRs are classified as Type A (augmented pharmacologic effect, predictable and dose-dependent, ~80\% of ADRs, e.g., bleeding with warfarin, hypoglycemia with insulin) and Type B (bizarre, unpredictable, immune-mediated or idiosyncratic, e.g., penicillin anaphylaxis, statin-induced myopathy). Risk factors: polypharmacy, advanced age, renal/hepatic impairment, genetic polymorphisms (CYP450 variants, HLA-B*5701 for abacavir hypersensitivity), and drug--drug interactions (e.g., simvastatin--clarithromycin). Management: prompt recognition, discontinuation or dose reduction, switching to alternative agent, and reporting to pharmacovigilance systems (FDA MedWatch). The Naranjo algorithm and causality assessment tools aid in determining likelihood of ADR.

\medterm{Memantine}-- An NMDA receptor antagonist (uncompetitive, low-to-moderate affinity) used for moderate to severe Alzheimer disease, often added to an acetylcholinesterase inhibitor. Blocks excessive glutamate-mediated excitotoxicity while preserving normal synaptic transmission. Once or twice daily oral. Adverse effects: dizziness, headache, constipation, confusion, and hypertension. Minimal cognitive side effects; renally cleared (reduce dose in renal impairment). Also used off-label for other dementias and neuropathic pain.

\medterm{Meropenem-vaborbactam}-- A combination of the carbapenem meropenem with the cyclic boronate beta-lactamase inhibitor vaborbactam. Vaborbactam inhibits class A and class C beta-lactamases, including KPC carbapenemases, but does not inhibit class D or metallo-beta-lactamases. This combination has activity against many carbapenem-resistant Enterobacteriaceae. It is indicated for complicated urinary tract infections and is being studied for other indications. The combination is administered intravenously every 8 hours, withwith dose adjustment required for impaired renal function; the most common adverse effects include diarrhea, headache, nausea, and infusion-site reactions.

\medterm{Methicillin-resistant Staphylococcus aureus}-- Staphylococcus aureus resistant to beta-lactam antibiotics through acquisition of the mecA gene encoding penicillin-binding protein 2a (PBP2a). MRSA is a major cause of healthcare-associated and community-associated infections. Treatment options include vancomycin, daptomycin, linezolid, and trimethoprim-sulfamethoxazole. MRSA infections are associated with increased morbidity and mortality. Infection prevention strategies include hand hygiene, contact precautions, and decolonization protocols. MRSA colonisation and infection surveillance, decolonisation strategies using intranasal mupirocin and chlorhexidine bathing, and antimicrobial stewardship to preserve the effectiveness of remaining active agents.

\medterm{Metoclopramide}-- A dopamine (D2) receptor antagonist and 5-HT3 antagonist used as a prokinetic antiemetic for gastroparesis, chemotherapy-induced and postoperative nausea, and GERD (facilitates gastric emptying). Adverse effects: acute dystonic reactions (young women), tardive dyskinesia (chronic use—black-box warning), parkinsonism, hyperprolactinemia, and QT prolongation. Contraindicated in GI obstruction, pheochromocytoma, and Parkinson disease. Domperidone is an alternative (peripheral).

\medterm{Metronidazole}-- A nitroimidazole antibiotic and antiprotozoal agent that is selectively toxic to anaerobic organisms and certain parasites. The drug is activated by reduction of its nitro group within anaerobic cells, generating toxic free radicals that damage DNA. Metronidazole is highly effective against anaerobic bacteria including Bacteroides fragilis, Clostridium species, and Helicobacter pylori, as well as protozoa including Trichomonas vaginalis, Giardia lamblia, and Entamoeba histolyticlytica. It is the drug of choice for C. difficile infection when used orally, for anaerobic infections above the diaphragm, and for bacterial vaginosis. Common adverse effects include metallic taste, nausea, and a disulfiram-like reaction with alcohol; prolonged use may cause peripheral neuropathy.

\medterm{Micafungin}-- An echinocandin antifungal agent used in the treatment and prophylaxis of invasive candidiasis and esophageal candidiasis. Micafungin inhibits beta-(1,3)-D-glucan synthase, disrupting fungal cell wall synthesis and leading to osmotic instability and cell death. A key advantage is its fungicidal activity against Candida species and fungistatic activity against Aspergillus. The drug is available only intravenously, as it is not absorbed orally. Common adverse effects include headache, nausea, and elevated transaminases; micafungin is also associated with an increased risk of hepatocellular tumors in animal studies, though the clinical significance in humans is uncertain.

\medterm{Milrinone}-- A phosphodiesterase-3 inhibitor with positive inotropic and vasodilatory (inodilator) actions, increasing cAMP in cardiac and vascular smooth muscle. Used short-term for acute decompensated heart failure and low cardiac output (bridging therapy, post-cardiac surgery), especially when beta-agonists are inadequate or patients are beta-blocked. IV infusion; half-life 2-3 hours (longer in renal impairment). Adverse effects: hypotension, arrhythmias (including VT), thrombocytopenia, and hypokalemia. May increase mortality with chronic use in class IV heart failure.

\medterm{Minimum inhibitory concentration}-- The lowest concentration of an antimicrobial agent that inhibits visible growth of a microorganism under standardized in vitro conditions. MIC is the gold standard for measuring antimicrobial susceptibility and is determined by broth microdilution or agar dilution methods. The MIC value is used to categorize isolates as susceptible, intermediate, or resistant based on established breakpoints. MIC values guide dosing decisions and can predict clinical outcomes. For some infections, including endoendocarditis and osteomyelitis, maintaining drug concentrations above the MIC for the entire dosing interval is critical for achieving bacteriological cure, and time above the MIC is the key pharmacodynamic parameter for beta-lactam antibiotics.

\medterm{Minoxidil}-- A direct arteriolar vasodilator (potassium channel opener) used for severe refractory hypertension (usually with a beta-blocker and loop diuretic to counter reflex tachycardia and fluid retention), and topically for androgenetic alopecia (promotes hair regrowth). Adverse effects: hypertrichosis, pericardial effusion, reflex tachycardia, fluid retention, and headache. Dose is titrated; discontinue slowly. Topical use is generally well tolerated with minimal systemic effects.

\medterm{Modafinil}-- A wakefulness-promoting agent (mechanism not fully understood, likely dopaminergic) used for narcolepsy, obstructive sleep apnea residual sleepiness, and shift work sleep disorder; used off-label for fatigue and ADHD. Less abuse potential and fewer sympathomimetic effects than amphetamines. Adverse effects: headache, nausea, nervousness, and rare severe rash (SJS); decreases effectiveness of hormonal contraceptives (CYP3A4 induction). Once daily.

\medterm{Montelukast}-- A leukotriene receptor antagonist (LTRA) blocking cysteinyl leukotrienes (LTC4, LTD4, LTE4) that cause bronchoconstriction, mucus secretion, and airway inflammation. Used as add-on maintenance therapy for asthma (particularly exercise-induced and aspirin-sensitive), allergic rhinitis, and bronchospasm prophylaxis. Oral, once daily. Generally well tolerated; adverse effects include headache and abdominal pain. A boxed warning addresses neuropsychiatric events (mood changes, agitation, depression, suicidal ideation). Not indicated for acute bronchospasm.

\medterm{Mycobacterium avium complex}-- A group of slowly growing mycobacteria that cause pulmonary disease and disseminated infection in immunocompromised patients. See microbiology.

\medterm{Narcotic}-- A term derived from the Greek narkotikos (to numb or stupor), historically referring to opium-derived drugs that induce sleep or stupor and relieve pain. In modern medical usage, the term is largely replaced by the more specific term opioid. Legally, the term narcotic in the Controlled Substances Act refers broadly to any drug, including cocaine and coca leaves, that is considered addictive and abused, though cocaine and coca are pharmacologically stimulants, not narcotics. The imprecise legal definition has led to confusion. In clinical practice, the preferred term for morphine-like analgesics is opioid, which specifies drugs acting at opioid receptors.

\medterm{Narrow Therapeutic Index}-- A drug characteristic where there is a small difference between the therapeutic dose and the toxic dose, requiring careful dosing and therapeutic drug monitoring to ensure efficacy while avoiding toxicity. Narrow therapeutic index drugs have a therapeutic index <2-3 (the ratio of toxic dose to effective dose). Examples: warfarin (TI ~2), digoxin, aminoglycosides (gentamicin, tobramycin, amikacin), vancomycin, cyclosporine, tacrolimus, lithium, phenytoin, carbamazepine, valproic acid, theophylline, and thyroid hormones. Managing NTI drugs requires: regular monitoring of drug concentrations; individualized dosing based on pharmacokinetic parameters; awareness of drug interactions that can alter levels; and patient education about signs of toxicity.

\medterm{Neuraminidase Inhibitors}-- A class of antiviral drugs that inhibit the neuraminidase enzyme on the surface of influenza viruses, preventing the release of newly formed viral particles from infected cells and thereby limiting viral spread. The two available agents are oseltamivir, administered orally, and zanamivir, administered by inhalation. These drugs are effective against both influenza A and influenza B viruses. Neuraminidase cleaves sialic acid residues from glycoproteins on the cell surface, allowing release of newly formed viral particles; inhibition of neuraminidase prevents this release, limiting viral spread within the respiratory tract and reducing the duration and severity of influenza symptoms when started within 48 hours of symptom onset.

\medterm{Nitrates}-- Vasodilator drugs that release nitric oxide, causing relaxation of vascular smooth muscle through activation of guanylyl cyclase and increased cyclic GMP. They primarily cause venodilation at low doses, reducing preload and myocardial oxygen demand, and arterial dilation at higher doses, reducing afterload. Nitroglycerin is the prototype nitrate, available in sublingual tablets for acute angina, intravenous infusion for acute coronary syndromes, topical patches for prophylaxis, and and oral sustained-release formulations for long-term prophylaxis. Nitrate tolerance develops rapidly with continuous exposure and requires a daily nitrate-free interval of 8-12 hours to maintain therapeutic efficacy; common adverse effects include headache, flushing, and postural hypotension.

\medterm{Nitroglycerin}-- An organic nitrate that releases nitric oxide, relaxing vascular smooth muscle—predominantly venous. Used for acute angina (sublingual spray/tablet, IV for ACS/heart failure), angina prophylaxis, and preload reduction in acute pulmonary edema. Rapid onset (1-3 min sublingual). Adverse effects: headache, flushing, hypotension, syncope, reflex tachycardia; tolerance with continuous use (nitrate-free interval required). Contraindicated with PDE5 inhibitors (sildenafil/tadalafil—severe hypotension) and in right ventricular infarction, severe aortic stenosis, and raised intracranial pressure.

\medterm{Nitroprusside}-- A direct-acting vasodilator (arterial and venous) releasing nitric oxide, used for hypertensive emergencies and controlled hypotension during surgery. IV infusion with rapid onset and offset; reduces both preload and afterload. Metabolized to cyanide and thiocyanate—toxicity (with prolonged/high-dose infusion, renal/hepatic impairment) causes lactic acidosis, confusion, and seizures. Protect from light. Adverse effects: hypotension, cyanide/thiocyanate toxicity, and methemoglobinemia. Monitor blood pressure continuously and cyanide levels with prolonged use.

\medterm{Noncompetitive Antagonist}-- A noncompetitive antagonist is a drug that reduces the effect of an agonist by a mechanism that cannot be overcome by increasing the agonist concentration. This occurs either through irreversible binding to the receptor, preventing agonist binding permanently, or through binding to an allosteric site that reduces receptor function regardless of agonist binding. The dose-response curve of the agonist is shifted to the right with a decrease in maximal response, meaning the full effect of the agonionist is reduced or eliminated, distinguishing noncompetitive antagonism from competitive antagonism; this type of antagonism is often irreversible on the timescale of clinical drug therapy and requires synthesis of new receptors for recovery of full agonist responsiveness.

\medterm{Non-tuberculous mycobacteria}-- A group of mycobacterial species other than M. tuberculosis complex that cause disease in humans. NTM include rapidly growing species (M. abscessus, M. chelonae, M. fortuitum) and slowly growing species (M. avium complex, M. kansasii). NTM infections are increasingly recognized, particularly in patients with chronic lung disease and immunocompromise. Treatment requires combination therapy based on species identification and susceptibility testing. Macrolides are backbone agents for many NTM infections; treatment is prolonged and requires multiple agents, with regimens tailored to species identification and susceptibility results, and outcomes vary widely depending on the species, extent of disease, and host immune status.

\medterm{NSAIDs}-- Nonsteroidal anti-inflammatory drugs are a class of drugs that inhibit cyclooxygenase enzymes, COX-1 and COX-2, reducing the synthesis of prostaglandins and thromboxanes from arachidonic acid. They produce analgesic, anti-inflammatory, antipyretic, and antiplatelet effects. COX-1 is constitutively expressed and maintains gastric mucosal protection, renal blood flow, and platelet aggregation. COX-2 is inducible and upregulated during inflammation. Traditional NSAIDs including ibuprofen, naproxen,and diclofenac have analgesic, antipyretic, and anti-inflammatory effects but carry risks including gastrointestinal bleeding and ulceration, renal impairment, and cardiovascular events; the risk-benefit profile must be evaluated for each patient based on their individual risk factors.

\medterm{NSAIDs (Nonsteroidal Anti-Inflammatory Drugs)}-- A class of analgesic, antipyretic, and anti-inflammatory drugs that inhibit cyclooxygenase (COX-1 and/or COX-2), reducing prostaglandin and thromboxane synthesis. Nonselective NSAIDs (ibuprofen, naproxen, indomethacin) block both COX-1 (gastric protection, platelets) and COX-2 (inflammation); COX-2 selective agents (celecoxib) reduce GI ulceration but retain cardiovascular risk. Uses: pain, inflammation, fever, arthritis, gout, dysmenorrhea. Adverse effects: GI bleeding/ulceration, renal impairment, fluid retention, hypertension, cardiovascular events, and aspirin-exacerbated respiratory disease.

\medterm{Nucleoside Analogs}-- A class of antiviral drugs that structurally resemble natural nucleosides and inhibit viral replication by interfering with DNA or RNA synthesis. They require intracellular phosphorylation by viral or cellular kinases to their active triphosphate form, which then competes with natural nucleotides for incorporation into the growing viral nucleic acid chain, causing chain termination or lethal mutagenesis. Examples include acyclovir for herpes simplex virus and varicella-zososter virus, tenofovir for HIV and hepatitis B, and ribavirin for hepatitis C, with resistance developing through mutations in the target viral enzymes that reduce drug binding affinity.

\medterm{Off-Label Use}-- The use of a licensed medication for an indication, dose, route, or patient population not included in its marketing authorisation (product licence). Off-label prescribing is legal, common in many areas of medicine (pediatrics—up to 50\% of hospital prescriptions, oncology, psychiatry, obstetrics), and may represent best clinical practice when evidence supports the use (e.g., $\beta$-blockers in heart failure were off-label before regulatory approval). The prescriber bears responsibility for the patient's safety and should: (1) be satisfied that there is a sufficient evidence base (published studies, guidelines, expert consensus) to support the off-label use; (2) explain the off-label nature to the patient and obtain informed consent; (3) document the rationale in the medical record. Regulatory frameworks: in the UK, the MHRA allows off-label prescribing under the prescriber's professional responsibility; the FDA prohibits pharmaceutical companies from promoting off-label use but does not restrict physician prescribing. Notable examples: ketamine for treatment-resistant depression, metformin for gestational diabetes, and aspirin for primary prevention of cardiovascular disease (controversial).

\medterm{Ointment}-- A semi-solid, greasy topical preparation for application to the skin, mucous membranes, or eyes. Ointments are typically composed of a hydrophobic base (petrolatum, lanolin, paraffin, beeswax, vegetable oils) and may contain dissolved or suspended active pharmaceutical ingredients. Ointments have high occlusivity (form a water-repellent barrier, preventing transepidermal water loss) and are preferred for dry, scaly, or lichenified skin conditions (psoriasis, chronic eczema, xerosis). They have higher drug penetration than creams (oil-in-water emulsions) or lotions due to their lipophilic nature. Disadvantages: greasy feel, staining of clothing, and less cosmetic acceptability than creams. Eye ointments (e.g., chloramphenicol, tetracycline) provide prolonged ocular contact time and are often used at bedtime. Ointment bases can be classified as hydrocarbon (petrolatum), absorption (lanolin—can absorb water), water-removable (emulsifying ointment), and water-soluble (polyethylene glycol).

\medterm{Omadacycline}-- A newer aminomethylcycline antibiotic with activity against gram-positive, gram-negative, and atypical organisms. Omadacycline inhibits bacterial protein synthesis by binding to the 30S ribosomal subunit. It has activity against MRSA, vancomycin-resistant enterococci, community-acquired pneumonia pathogens, and skin infections. The drug is available in both intravenous and oral formulations, allowing for IV-to-oral step-down therapy. Omadacycline is indicated for community-acquired bacterial pneneumonia and acute bacterial skin infections, with common adverse effects including nausea, vomiting, and infusion-site reactions; photosensitivity is less common than with earlier tetracyclines.

\medterm{Omeprazole}-- A proton pump inhibitor (PPI) that irreversibly inhibits the gastric H+/K+ ATPase (proton pump) on parietal cells, profoundly reducing acid secretion. Used for GERD, peptic ulcer disease, H. pylori eradication (triple therapy), Zollinger-Ellison syndrome, and stress ulcer prophylaxis. Prodrug activated in acidic canaliculi; once daily before breakfast. Adverse effects: headache, diarrhea, hypomagnesemia, B12 deficiency with long-term use, osteoporosis/fracture risk, C. difficile infection, and drug interactions (CYP2C19). Pantoprazole, esomeprazole, lansoprazole are alternatives.

\medterm{Opiate}-- A natural alkaloid derived from the opium poppy (Papaver somniferum) or a semi-synthetic derivative that binds to opioid receptors. Natural opiates: morphine (prototypical opioid analgesic, gold standard for severe pain), codeine (prodrug metabolized to morphine by CYP2D6—poor metabolisers derive minimal analgesia), and thebaine (used in synthesis of semi-synthetic opioids—oxycodone, buprenorphine, naloxone). Semi-synthetic opiates: heroin (diamorphine—more lipid-soluble than morphine, rapid CNS penetration), oxycodone, and hydromorphone. See also Opioid for the broader class including synthetic opioids (fentanyl, methadone, tramadol). Historically, opium has been used for pain relief, diarrhea, and cough suppression for millennia. Opiates produce euphoria, miosis, respiratory depression, constipation, and physical dependence. Overuse leads to opioid use disorder (addiction).

\medterm{Opioid Analgesics}-- A class of drugs that bind to opioid receptors in the central and peripheral nervous systems, producing analgesia and other effects. There are three main opioid receptor types: mu receptors mediating analgesia, respiratory depression, euphoria, and physical dependence; delta receptors involved in analgesia and mood; and kappa receptors mediating analgesia and sedation. Morphine is the prototype opioid, derived from opium, with other natural opioids including codeine and semsynthetic opioids including fentanyl and methadone. Common adverse effects include respiratory depression, constipation, nausea and vomiting, sedation, tolerance, physical dependence, and the potential for addiction; careful dosing and monitoring are required to balance effective pain relief with the risk of adverse events.

\medterm{Opioids}-- A class of drugs that bind to opioid receptors in the central and peripheral nervous systems, producing analgesia and other effects. See Opioid Analgesics.

\medterm{Oral Hypoglycemic}-- A class of drugs used to treat type 2 diabetes (with lifestyle). Classes: biguanides (metformin—first-line), sulfonylureas (glipizide, glyburide), meglitinides, thiazolidinediones (pioglitazone), DPP-4 inhibitors (sitagliptin), SGLT2 inhibitors (empagliflozin), GLP-1 receptor agonists (semaglutide, liraglutide—injectable), and alpha-glucosidase inhibitors. Choice based on glycemic control, weight, cardiovascular/renal risk, and hypoglycemia risk. SGLT2 inhibitors and GLP-1 agonists have cardiovascular and renal benefits; sulfonylureas and insulin carry hypoglycemia risk.

\medterm{Orphan Drug}-- A pharmaceutical agent developed specifically to treat a rare medical condition (affecting fewer than 200,000 people in the United States, as defined by the Orphan Drug Act of 1983). Orphan drug designation provides incentives including seven-year market exclusivity after FDA approval, tax credits for clinical trial costs, waived FDA application fees, and research grants. Examples include lumacaftor/ivacaftor for cystic fibrosis, nusinersen for spinal muscular atrophy, and eliglustat for Gaucher disease. The Orphan Drug Act has successfully stimulated development of treatments for previously neglected rare diseases.

\medterm{Orphan Drugs}-- Medicinal products developed for the diagnosis, prevention, or treatment of rare diseases (defined in the EU as affecting <5 per 10,000 people; in the US as affecting <200,000 people). Orphan drug designation provides incentives for pharmaceutical companies: 7 years of market exclusivity in the US (Orphan Drug Act 1983), 10 years in the EU (2000), reduced regulatory fees, tax credits for clinical research costs (25\% in the US), and protocol assistance from regulatory agencies. More than 600 orphan drugs have been approved in the US since 1983. Examples: imiglucerase and miglustat (Gaucher disease), ivacaftor (cystic fibrosis—G551D mutation), eculizumab (paroxysmal nocturnal haemoglobinuria, atypical hemolytic uraemic syndrome), nusinersen (spinal muscular atrophy), and eliglustat (Gaucher disease type 1). Challenges: high drug costs (often >\$200,000 per patient per year), limited natural history data, small trial populations, and ethical concerns about equitable access. The UK Rare Diseases Framework (2021) prioritises timely diagnosis and access to treatment for rare diseases.

\medterm{Over-the-Counter Drug}-- A medication that can be purchased without a prescription, deemed safe and effective for self-medication when used according to label instructions. OTC drugs have established safety profiles, low toxicity, and clear dosing guidelines suitable for consumer use. Examples include acetaminophen, ibuprofen, antihistamines (loratadine, diphenhydramine), antacids, and topical antifungals. The FDA determines OTC status based on safety, labeling adequacy, and the condition being treatable without professional supervision. Some drugs are switched from prescription to OTC status after post-marketing safety data accumulate.

\medterm{Over-the-Counter (OTC)}-- Medications available without a prescription, considered safe and effective for self-care with appropriate labeling and low abuse potential. Examples: acetaminophen, ibuprofen, antacids, antihistamines, decongestants, laxatives. Regulated in the US by the FDA (OTC monograph system) and Rx-to-OTC switches (e.g., omeprazole, loratadine). Despite availability, OTC drugs carry risks (hepatotoxicity from acetaminophen overdose, NSAID GI/renal effects, interactions with warfarin and antihypertensives). Pharmacists counsel on appropriate use.

\medterm{Palliative Care}-- Specialized medical care for serious illness focusing on symptom relief and quality of life. See oncology.

\medterm{Pantoprazole}-- A proton pump inhibitor reducing gastric acid secretion by irreversibly inhibiting the H+/K+ ATPase of parietal cells. Used for GERD, peptic ulcer disease, dyspepsia, H. pylori eradication, and acid hypersecretory states. Available oral and IV (used for stress ulcer prophylaxis in the ICU). Generally well tolerated; long-term use is associated with hypomagnesemia, vitamin B12 deficiency, increased fracture risk, and C. difficile infection. Fewer CYP interactions than omeprazole; esomeprazole is the S-enantiomer of omeprazole.

\medterm{Partial Agonist}-- A partial agonist is a drug that binds to a receptor and produces a biological response that is less than the maximum possible response, even when all receptors are occupied. Partial agonists have affinity for the receptor but lower intrinsic activity compared to full agonists, producing a submaximal response regardless of dose. The maximal effect of a partial agonist is always less than that of a full agonist acting at the same receptor. An important property of partial agonists is that they cacan act as functional antagonists in the presence of a full agonist: by occupying receptors that would otherwise be activated by the full agonist, the partial agonist reduces the net response, which is the basis for their use in maintenance therapy for opioid dependence and smoking cessation.

\medterm{Pharmacodynamics}-- The study of the biochemical and physiological effects of drugs on the body and their mechanisms of action, describing what the drug does to the body. It encompasses drug-receptor interactions (binding affinity, intrinsic activity, efficacy), signal transduction pathways (second messengers, gene expression), dose-response relationships, and therapeutic vs. adverse effects. Key concepts: receptors — proteins or macromolecules that drugs bind to exert their effects; agonists — drugs that activate receptors to produce a response; antagonists — drugs that block receptor activation; partial agonists — drugs that produce a submaximal response even at full receptor occupancy; and allosteric modulators — drugs that bind to sites distinct from the orthosteric binding site to modulate receptor activity.

\medterm{Pharmacogenomics}-- The study of how genetic variations influence individual responses to drugs, including drug efficacy, toxicity, and metabolism. Genetic polymorphisms in drug-metabolizing enzymes, transporters, and receptors can lead to significant inter-individual variability in drug response. Examples: CYP2D6 poor metabolizers have reduced codeine activation to morphine, leading to ineffective analgesia; CYP2C19 poor metabolizers require reduced clopidogrel dosing; TPMT deficiency increases risk of severe myelosuppression with azathioprine; HLA-B*5701 testing is required before abacavir to prevent hypersensitivity reaction; UGT1A1*28 homozygotes have increased risk of irinotecan toxicity. Pharmacogenomic testing is increasingly integrated into clinical practice, guiding drug selection and dosing to optimize efficacy and minimize adverse effects.

\medterm{Pharmacokinetics}-- The study of drug movement through the body, describing what the body does to the drug, defined by four processes: Absorption — the process by which a drug enters the systemic circulation from its site of administration (oral, IV, IM, subcutaneous, topical, inhalation); Distribution — the movement of drug from the circulation to body tissues, determined by blood flow, protein binding, tissue affinity, and barriers (blood-brain barrier, placental barrier); Metabolism — the chemical modification of drugs by enzymes (primarily hepatic CYP450) through Phase I (oxidation, reduction, hydrolysis) and Phase II (conjugation) reactions; and Excretion — the elimination of drugs and their metabolites from the body, primarily through the kidneys (glomerular filtration, tubular secretion, reabsorption) and also through bile, feces, lungs, and sweat. Pharmacokinetic parameters: half-life, clearance, volume of distribution, bioavailability, and steady state concentration.

\medterm{Pharmacovigilance}-- The science and activities relating to the detection, assessment, understanding, and prevention of adverse drug reactions and other drug-related problems. It encompasses post-marketing surveillance of drugs after they have been approved for clinical use, as pre-marketing clinical trials may not detect rare adverse effects, long-term effects, or effects in special populations. Pharmacovigilance systems collect adverse event reports from healthcare professionals, patients, andand pharmaceutical companies, and uses data from multiple sources including clinical trials, spontaneous reporting, electronic health records, and published literature to generate safety signals that may lead to regulatory actions including label changes, warnings, and market withdrawals.

\medterm{Phase II Metabolism}-- Phase II metabolism, also known as conjugation reactions, involves the attachment of endogenous hydrophilic molecules to drugs or their Phase I metabolites, dramatically increasing water solubility and facilitating renal excretion. The major Phase II reactions include glucuronidation catalyzed by UDP-glucuronosyltransferases, which is the most common conjugation reaction and accounts for approximately 35 percent of Phase II metabolism; sulfation catalyzed by sulfotransferases; acetylation catalycatalysed by N-acetyltransferases; methylation; and conjugation with glutathione. Phase II metabolites are generally pharmacologically inactive and are rapidly excreted in urine or bile; the efficiency of these conjugation pathways depends on the availability of the endogenous substrate and the activity of the conjugating enzymes, which can be influenced by genetic polymorphisms, drug interactions, and disease states.

\medterm{Phase I Metabolism}-- Phase I metabolism involves chemical modification of drug molecules through oxidation, reduction, and hydrolysis reactions, primarily catalyzed by the cytochrome P450 enzyme system in the liver. These reactions introduce or expose functional groups such as hydroxyl, amino, sulfhydryl, and carboxyl groups, typically increasing the polarity and water solubility of the parent compound. The most common Phase I reaction is oxidation, catalyzed by CYP enzymes requiring molecular oxygen and NADPH. OxidOxidation reactions introduce hydroxyl, carboxyl, or epoxide groups, while reduction reactions add hydrogen or remove oxygen, and hydrolysis reactions cleave ester or amide bonds; the resulting metabolites may be pharmacologically inactive, less active, equally active, or occasionally more active than the parent compound, and some undergo further Phase II conjugation for elimination.

\medterm{Phenylephrine}-- A selective alpha-1 adrenergic agonist causing vasoconstriction and increased blood pressure (systemic vascular resistance) with reflex bradycardia. Used for hypotension/shock (IV bolus/infusion, common in anesthesia and sepsis), nasal congestion (topical decongestant), and mydriasis (ophthalmic). Adverse effects: hypertension, reflex bradycardia, headache, and tissue ischemia with extravasation. Caution in hypertensive and cardiac patients; contraindicated in severe hypertension, hyperthyroidism, and narrow-angle glaucoma (ophthalmic use).

\medterm{Phosphodiesterase Inhibitor}-- A class of drugs that inhibit phosphodiesterase enzymes, which break down cyclic nucleotides (cAMP and cGMP), thereby increasing intracellular levels of these second messengers and potentiating their downstream effects. Nonselective PDE inhibitors: theophylline (PDE3 and PDE4 inhibition) — bronchodilator for asthma and COPD, also has anti-inflammatory effects. PDE3 inhibitors: milrinone — positive inotrope and vasodilator for acute decompensated heart failure. PDE4 inhibitors: roflumilast — used for severe COPD with chronic bronchitis and frequent exacerbations; apremilast — used for psoriasis and psoriatic arthritis. PDE5 inhibitors: sildenafil, tadalafil, vardenafil — inhibit cGMP-specific PDE5 in corpus cavernosum, used for erectile dysfunction; also treat pulmonary arterial hypertension (sildenafil, tadalafil). Each PDE isozyme has distinct tissue distribution, allowing targeted therapeutic effects.

\medterm{Polyethylene Glycol (PEG)}-- An osmotic laxative used for constipation and bowel preparation for colonoscopy. PEG (e.g., Miralax, GoLYTELY) is a non-absorbable polymer that draws water into the colonic lumen, softening stool and stimulating defecation without significant electrolyte shifts. For chronic constipation: 17 g daily in water. For bowel prep: large-volume balanced electrolyte solutions (PEG-ELS) or low-volume with ascorbate. Generally well tolerated; bloating and cramping may occur; contraindicated in bowel obstruction.

\medterm{Potassium-Sparing Diuretics}-- Potassium-sparing diuretics act on the distal nephron to promote sodium excretion while minimizing potassium loss. They include epithelial sodium channel blockers such as amiloride and triamterene, and mineralocorticoid receptor antagonists such as spironolactone and eplerenone. Amiloride and triamterene directly block the epithelial sodium channel in the collecting duct, reducing sodium reabsorption and potassium secretion. Spironolactone and eplerenone competitively inhibit aldosterone bindingto the mineralocorticoid receptor, blocking aldosterone-mediated sodium reabsorption and potassium excretion. Spironolactone is used as a diuretic in conditions associated with secondary hyperaldosteronism, including heart failure and hepatic cirrhosis, and as an antiandrogen in the treatment of acne, hirsutism, and female pattern hair loss. Hyperkalaemia is the most important adverse effect and requires regular monitoring of serum potassium levels, particularly in patients with renal impairment or those receiving concurrent ACE inhibitors or ARBs.

\medterm{Potency vs Efficacy}-- Potency is the amount of drug required to produce a given effect, measured by the EC50 (the concentration producing 50\% of maximal effect). Drugs with higher potency achieve effects at lower doses (e.g., fentanyl is more potent than morphine). Potency is clinically relevant for dosing but does not indicate which drug is therapeutically superior. Efficacy is the maximum effect a drug can produce at saturating concentrations, measured by the Emax (plateau of the dose-response curve). Efficacy determines the clinical utility of a drug — a drug with high efficacy can produce an effect that a low-efficacy drug cannot. For example, morphine (high efficacy mu-opioid agonist) can relieve severe pain that codeine (low efficacy partial agonist) cannot. A drug can have high potency but low efficacy and vice versa.

\medterm{Pramipexole}-- A non-ergot dopamine agonist (D2/D3) used for Parkinson disease (as monotherapy in early disease or adjunct to levodopa) and restless legs syndrome. Activates striatal dopamine receptors. Adverse effects: nausea, orthostatic hypotension, somnolence (sudden sleep attacks—caution driving), impulse control disorders (pathological gambling, hypersexuality—important to warn), hallucinations, and leg edema. Renally cleared—reduce dose in renal impairment. Ropinirole and rotigotine (patch) are alternatives.

\medterm{Prazosin}-- A selective alpha-1 adrenergic antagonist causing vasodilation and reduced peripheral resistance. Used for hypertension, benign prostatic hyperplasia, and post-traumatic stress disorder-related nightmares (off-label). Marked first-dose orthostatic hypotension/syncope—start low and take at bedtime. Adverse effects: dizziness, orthostatic hypotension, headache, palpitations, and priapism (rare). Not first-line for hypertension but useful in combination. Shorter-acting than doxazosin/terazosin.

\medterm{Prescription}-- A written or electronic order by a licensed prescriber (physician, NP, PA, dentist) authorizing a pharmacy to dispense a specific drug to a specific patient. Components: patient identifiers, drug name, dose, route, frequency, quantity, refills, and prescriber signature. Controlled substances require DEA numbers and are scheduled by abuse potential (Schedule II-V). Prescription medications are regulated to ensure safety, require professional supervision, and are dispensed with labeling and counseling. Errors in prescribing are a leading cause of adverse drug events.

\medterm{Prescription Drug}-- A medication that can only be dispensed with a valid prescription from a licensed healthcare provider, requiring professional oversight due to the drug's complexity, toxicity, or potential for harm. Prescription drugs undergo rigorous FDA approval through clinical trials (Phase I-IV) to establish safety and efficacy before marketing. Examples include antibiotics, antihypertensives, antidepressants, and most psychotropic medications. Prescribing requires consideration of indication, dosing, drug interactions, and monitoring parameters.

\medterm{Prescription Medications}-- Pharmaceutical drugs that require authorization from a licensed healthcare provider before they can be dispensed to a patient. These medications are regulated because they carry risks that necessitate professional oversight, including potential side effects, drug interactions, and the need for precise dosing. Physicians prescribe them based on a specific diagnosis, and patients must follow the prescribed regimen closely to achieve therapeutic benefits while minimizing adverse outcomes.

\medterm{Procalcitonin-guided antibiotic therapy}-- The use of procalcitonin levels to guide initiation and discontinuation of antibiotics. Procalcitonin is a biomarker that rises in bacterial infections and falls with effective treatment. Elevated levels support the need for antibiotic therapy, while declining levels support discontinuation. Procalcitonin-guided protocols have been shown to reduce antibiotic duration and exposure in respiratory infections and sepsis. The approach is most validated for lower respiratory tract infections and sepsisis; however, false-positive results can occur in non-infectious inflammatory conditions and false-negative results in localised infections, so procalcitonin should be used as an adjunct to, rather than a replacement for, clinical judgement.

\medterm{Prodrug}-- A prodrug is a pharmacologically inactive compound that is converted to an active drug through metabolic processes within the body. Prodrugs are designed to overcome limitations of the parent drug including poor bioavailability, rapid metabolism, lack of tissue selectivity, or unacceptable side effects. Examples include enalapril, which is hydrolyzed by hepatic esterases to the active form enalaprilat, an ACE inhibitor that is poorly absorbed orally; clopidogrel, which requires CYP2C19-mediated mediated activation to its active metabolite; and valacyclovir, which undergoes hepatic hydrolysis to acyclovir, significantly improving oral bioavailability from 20\% to 55\%.

\medterm{Propafenone}-- A class IC antiarrhythmic with sodium channel blockade and weak beta-blocking activity. Used for paroxysmal atrial fibrillation (in patients without structural heart disease, including 'pill-in-the-pocket'), supraventricular tachycardias, and ventricular arrhythmias. Contraindicated in ischemic heart disease, structural heart disease, and significant LV dysfunction (proarrhythmia). Metabolized by CYP2D6 (poor metabolizers need lower doses). Adverse effects: proarrhythmia, dizziness, metallic taste, bronchospasm (weak beta-blockade), and GI upset.

\medterm{Propylthiouracil (PTU)}-- A thioamide antithyroid drug that inhibits thyroid peroxidase (hormone synthesis) and blocks peripheral T4-to-T3 conversion. Used for hyperthyroidism/Graves disease—preferred in the first trimester of pregnancy (lower teratogenicity than methimazole) and in thyroid storm. Adverse effects: agranulocytosis, hepatotoxicity (fulminant hepatitis—black-box warning, more than methimazole), ANCA vasculitis, rash, and arthralgia. Requires monitoring of CBC and LFTs. Longer action allows less frequent dosing.

\medterm{Protease Inhibitors}-- A class of antiviral drugs that block viral proteases, enzymes essential for processing viral polyproteins into functional components needed for viral assembly and maturation. In HIV treatment, protease inhibitors including ritonavir, lopinavir, atazanavir, and darunavir inhibit HIV-1 protease, preventing the cleavage of Gag and Gag-Pol polyproteins into structural proteins and enzymes. This results in the production of immature, non-infectious viral particles. Protease iinhibitors are also used in the treatment of hepatitis C, where NS3/4A protease inhibitors including simeprevir, grazoprevir, and glecaprevir block viral polyprotein processing and are used in combination with other direct-acting antiviral agents to achieve high rates of sustained virological response.

\medterm{Protein Binding}-- The reversible association of drug molecules with plasma proteins, primarily albumin for acidic drugs and alpha-1 acid glycoprotein for basic drugs. Only unbound or free drug is pharmacologically active and available for distribution to tissues, metabolism by the liver, and excretion by the kidneys. The fraction of drug bound to protein is described by the binding ratio, which depends on the total drug concentration, total protein concentration, and the associationassociation constant; changes in protein binding can significantly affect drug disposition, particularly for highly bound drugs (greater than 90\% bound), where small changes in binding can lead to large changes in free drug concentration, potentially affecting both efficacy and toxicity.

\medterm{Proton Pump Inhibitor}-- A class of medications that irreversibly inhibit the gastric H+/K+ ATPase (proton pump) in parietal cells, the final common pathway of gastric acid secretion, producing profound and sustained suppression of both basal and stimulated acid production. PPIs are prodrugs that require activation in the acidic environment of the parietal cell canaliculus, where they form a sulfenamide that binds covalently to cysteine residues of the proton pump. Available agents: omeprazole, lansoprazole, pantoprazole, rabeprazole, esomeprazole, and dexlansoprazole. Indications: gastroesophageal reflux disease (GERD), erosive esophagitis, peptic ulcer disease, Helicobacter pylori eradication (with antibiotics), Zollinger-Ellison syndrome, and prevention of NSAID-induced ulcers. Adverse effects: increased risk of Clostridioides difficile infection, pneumonia, hypomagnesemia, vitamin B12 deficiency, and osteoporosis-related fractures with long-term use. Rebound acid hypersecretion occurs upon discontinuation after >8 weeks of use.

\medterm{Proton Pump Inhibitors (PPIs)}-- A class of antisecretory drugs (e.g., omeprazole, lansoprazole, pantoprazole, esomeprazole) that irreversibly inhibit the $H^+/K^+$ ATPase enzyme (proton pump) in gastric parietal cells. Prodrugs activated in the acidic canaliculi of parietal cells. Indications: GERD, peptic ulcer disease, H. pylori eradication (triple therapy), Zollinger-Ellison syndrome, NSAID-induced ulcer prophylaxis. Long-term adverse effects: hypomagnesemia, Vitamin B12 deficiency, increased risk of Clostridioides difficile infection, hip fractures (decreased calcium absorption), and interstitial nephritis.

\medterm{Pseudoephedrine}-- A sympathomimetic (alpha-1 agonist, some beta effects) oral decongestant used for nasal congestion in colds and sinusitis. Increases nasal mucosal vasoconstriction. Adverse effects: insomnia, nervousness, tachycardia, hypertension, and urinary retention (BPH); avoid in uncontrolled hypertension, severe CAD, hyperthyroidism, and with MAOIs. Schedule/POS (behind-the-counter) in the US due to methamphetamine precursor use. Phenylephrine (oral) is an alternative with less efficacy.

\medterm{Pseudomonas aeruginosa resistance}-- Pseudomonas aeruginosa has intrinsic and acquired resistance mechanisms making it inherently resistant to many antibiotics. Intrinsic resistance includes low outer membrane permeability, efflux pumps, and chromosomal beta-lactamase. Acquired resistance can develop through mutations in target genes, acquisition of resistance plasmids, and overexpression of efflux pumps. Treatment typically requires antipseudomonal beta-lactams (piperacillin-tazobactam, cefepime, ceftazidime, carbapenems) or fluororoquinolones (ciprofloxacin, levofloxacin), often in combination therapy due to the high rate of resistance emergence during monotherapy; susceptibility testing is essential to guide treatment decisions.

\medterm{Psittacosis}-- See infectious\_disease.

\medterm{QT Prolongation}-- An electrocardiographic finding characterized by prolongation of the QT interval, representing the time from ventricular depolarization to repolarization, which predisposes to the potentially fatal arrhythmia torsades de pointes. The corrected QT interval should be less than 440 milliseconds in men and 460 milliseconds in women. Numerous medications can prolong the QT interval by blocking the cardiac potassium channel encoded by the KCNH2 gene, prolonging phase 3 of the cardiacardiac action potential, prolonging repolarisation and increasing the risk of triggered activity, torsades de pointes, and sudden cardiac death; risk factors include female sex, advanced age, electrolyte disturbances (hypokalaemia, hypomagnesaemia), bradycardia, structural heart disease, and the use of multiple QT-prolonging medications.

\medterm{Quinine}-- An alkaloid derived from cinchona bark, with antimalarial, antipyretic, and skeletal muscle relaxant properties. Quinine is a blood schizonticide with activity against all Plasmodium species, but resistance has emerged in P. falciparum. Current indications: (1) Treatment of severe/complicated falciparum malaria (WHO—quinine IV infusion plus doxycycline or clindamycin; artesunate IV is now preferred due to lower mortality). (2) Babesiosis (combined with clindamycin). (3) Nocturnal leg cramps (used off-label—risk of serious adverse effects outweighs benefit; only if cramps are disabling and no alternative exists). Mechanism: accumulates in the parasite's food vacuole, inhibits haem polymerase, and prevents detoxification of haem (toxic to the parasite). Side effects: cinchonism (tinnitus, dizziness, headache, nausea, visual disturbances—dose-related, reversible), QT prolongation, hypoglycaemia (stimulates insulin secretion—dangerous in severe malaria), hypersensitivity (urticaria, flushing, hemolytic anemia), blackwater fever (massive intravascular haemolysis—rare, idiosyncratic). Contraindications: G6PD deficiency (hemolytic anemia), QT prolongation, myasthenia gravis, optic neuritis. Therapeutic range: 8--15 \(\mu\)g/mL. Quinine is also used in tonic water (bitter flavouring, subtherapeutic concentrations).

\medterm{Quinupristin-dalfopristin}-- A combination streptogramin antibiotic that inhibits bacterial protein synthesis by binding to the 50S ribosomal subunit. The combination is bactericidal against most gram-positive organisms including MRSA. It is indicated for vancomycin-resistant Enterococcus faecium infections. The drug is administered intravenously every 8-12 hours. Common adverse effects include arthralgia, myalgia, and infusion-related reactions including phlebitis. The drug is a potent inhibitor of CYP3A4 and can interact with numerous medications, particularly those that inhibit CYP3A4 and prolong the QT interval.

\medterm{Receptor Desensitization}-- A rapid decrease in the response of a receptor to continuous or repeated agonist stimulation, occurring within minutes to hours. Desensitization mechanisms: receptor phosphorylation by G protein-coupled receptor kinases which uncouples the receptor from G proteins, leading to rapid functional desensitization; beta-arrestin binding which sterically prevents G protein coupling and targets the receptor for internalization; receptor internalization (endocytosis) into intracellular compartments, reducing cell-surface receptor number; and receptor degradation (downregulation) after prolonged stimulation, reducing total receptor protein. Desensitization contributes to drug tolerance (e.g., beta-agonist desensitization in asthma, opioid receptor desensitization in chronic pain) and can be reversed by removing the agonist, allowing receptor recycling to the cell surface.

\medterm{Renin Inhibitor}-- A class of antihypertensive drugs that directly inhibit renin, the first and rate-limiting enzyme in the renin-angiotensin-aldosterone system. Renin catalyzes the conversion of angiotensinogen to angiotensin I, which is then converted to angiotensin II by ACE. Aliskiren is the only direct renin inhibitor available. It binds to the active site of renin, preventing angiotensinogen cleavage and reducing angiotensin I, angiotensin II, and aldosterone levels. Indication: hypertension. Aliskiren is used alone or in combination with other antihypertensives. Adverse effects: diarrhea (especially at higher doses), cough, hyperkalemia, and angioedema. Contraindicated in pregnancy (fetal renin-angiotensin system toxicity). Combination with ACE inhibitors or ARBs in diabetic patients with renal impairment is contraindicated due to increased risk of hyperkalemia, hypotension, and renal impairment.

\medterm{Rezafungin}-- A long-acting echinocandin antifungal agent used in the treatment of invasive candidiasis and candidemia. Rezafungin inhibits beta-(1,3)-D-glucan synthase with a unique pharmacokinetic profile characterized by an extremely long half-life of approximately 130 hours, allowing for once-weekly dosing after an initial loading regimen. This extended dosing interval is a significant advantage over other echinocandins and may improve adherence. The drug has excellent stability in solution and can be stored at room temperature, maintaining potency during extended infusion periods; common adverse effects include headache, nausea, diarrhea, and hypokalaemia; hepatotoxicity has been observed in clinical trials.

\medterm{Ritalin}-- A central nervous system stimulant, a piperidine derivative that blocks the reuptake of dopamine and noradrenaline (norepinephrine) by inhibiting the dopamine transporter (DAT) and noradrenaline transporter (NAT) in the prefrontal cortex, increasing their synaptic concentrations. See Methylphenidate.

\medterm{Rivaroxaban}-- A direct oral anticoagulant and selective factor Xa inhibitor. Used for stroke prevention in nonvalvular atrial fibrillation, treatment/prophylaxis of venous thromboembolism, and thromboprophylaxis after hip/knee replacement. Fixed once-daily dosing, rapid onset, no routine monitoring. Reduces dose in renal impairment (avoid if CrCl <15-30 depending on indication) and with strong CYP3A4/P-gp inhibitors. Adverse effects: bleeding; reversal agent: andexanet alfa. Avoid in mechanical valves and moderate/severe mitral stenosis.

\medterm{Ropinirole}-- A non-ergot dopamine agonist (D2/D3) used for Parkinson disease and restless legs syndrome. Adverse effects: nausea, dizziness, somnolence (sudden sleep attacks), orthostatic hypotension, impulse control disorders (gambling, hypersexuality), hallucinations, and leg edema. Extended-release for once-daily Parkinson dosing. Renally cleared (caution in renal impairment). Dose slowly; risk of augmentation in restless legs syndrome with long-term use.

\medterm{Route of Administration}-- The path by which a drug enters the body, influencing absorption, onset, bioavailability, and first-pass metabolism. Enteral routes: oral, sublingual, buccal, rectal. Parenteral: intravenous (100 percent bioavailability, immediate onset), intramuscular, subcutaneous, intradermal. Other routes: inhalation, topical/transdermal, ophthalmic, otic, intranasal, vaginal, intrathecal, and intra-articular. Choice depends on drug properties (stability, absorption), therapeutic goal, and patient factors. Oral is most convenient; IV is used for emergencies and drugs with poor oral absorption.

\medterm{Sedative}-- A medication that reduces central nervous system activity, producing calming effects, drowsiness, and sleep at higher doses. Sedatives are used for anxiety, procedural sedation, insomnia, and as adjuncts in anesthesia. Classes include: benzodiazepines (midazolam, lorazepam, diazepam) — enhance GABA-A receptor activity, used for sedation, anxiolysis, and amnesia; barbiturates (phenobarbital, thiopental) — enhance GABA-A activity with broader CNS depression, largely replaced by benzodiazepines due to narrow therapeutic index and high dependence potential; propofol — an IV anesthetic that potentiates GABA-A receptors, rapid onset and offset, used for procedural sedation and general anesthesia; dexmedetomidine — an alpha-2 agonist producing sedation with minimal respiratory depression; and chloral hydrate — older sedative/hypnotic used for short-term sedation in children.

\medterm{Selective Estrogen Receptor Modulator}-- A class of medications that bind to estrogen receptors (ER-alpha and ER-beta) with tissue-specific agonist or antagonist activity, depending on the target tissue and differential co-regulator recruitment. Tamoxifen: antagonist in breast tissue (first-line adjuvant therapy for hormone receptor-positive breast cancer in premenopausal women), partial agonist in bone (preserves bone density) and endometrium (increases risk of endometrial cancer and hyperplasia), and improves lipid profile. Raloxifene: antagonist in breast and agonist in bone without endometrial stimulation (used for osteoporosis prevention and treatment in postmenopausal women, reduces breast cancer risk). Toremifene: similar to tamoxifen for metastatic breast cancer. Bazedoxifene: used for postmenopausal osteoporosis, often combined with conjugated estrogens (duavee). Adverse effects: tamoxifen — hot flashes, vaginal dryness, and thromboembolism (2-3\% risk). SERMs do not cause the systemic side effects of estrogen therapy.

\medterm{Selegiline}-- A selective, irreversible MAO-B inhibitor used for Parkinson disease (mild early disease or adjunct to levodopa) and (as oral disintegrating/selegiline transdermal) for depression. Spares MAO-A in the gut at low oral doses, reducing tyramine risk (still require dietary precautions with higher doses or patch). Adverse effects: nausea, insomnia, dyskinesias, hypotension, and interactions with opioids (meperidine—serotonin syndrome) and SSRIs/SNRIs (avoid overlap).

\medterm{Senna}-- A stimulant laxative containing anthraquinone glycosides (sennosides) that are activated by colonic bacteria to stimulate colonic motility and fluid secretion. Onset: 6-12 hours; used for constipation and bowel preparation. Common side effects: abdominal cramps, nausea; chronic use may cause melanosis coli and electrolyte disturbances. Safe in pregnancy (short-term) and generally first-line among stimulants. Available as tablets, syrup, and suppositories.

\medterm{Sertraline}-- A selective serotonin reuptake inhibitor (SSRI) used for major depressive disorder, panic disorder, PTSD, OCD, social anxiety disorder, and premenstrual dysphoric disorder. Effective and well tolerated; mild CYP2D6 inhibition. Adverse effects: nausea, diarrhea, insomnia, sexual dysfunction, and hyponatremia (elderly); serotonin syndrome and discontinuation syndrome (taper); black-box warning for suicidality in young adults. Once-daily dosing; available in pregnancy-adjacent indications (postpartum depression) with monitoring.

\medterm{SGLT2 Inhibitors}-- Sodium-glucose co-transporter 2 inhibitors are a class of oral antidiabetic agents that block the reabsorption of glucose in the proximal convoluted tubule of the kidney, promoting glycosuria and reducing blood glucose levels. The kidney normally filters approximately 180 grams of glucose daily, with SGLT2 responsible for approximately 90 percent of glucose reabsorption. By blocking this transporter, SGLT2 inhibitors increase urinary glucose excretion by 60 to 80 grams per day. Common agents incinclude dapagliflozin, empagliflozin, and canagliflozin. In addition to glycaemic control, SGLT2 inhibitors have demonstrated significant cardiovascular and renal benefits in major outcome trials, reducing the risk of heart failure hospitalisation and progression of chronic kidney disease, making them an important treatment option in patients with type 2 diabetes and established cardiovascular or renal disease.

\medterm{Side Effect}-- An unintended pharmacological effect of a drug occurring at therapeutic doses, usually predictable from the drug's mechanism of action (a Type A reaction). Side effects are dose-dependent and may be bothersome but are often tolerable or transient (e.g., dry mouth with anticholinergics, sedation with antihistamines, cough with ACE inhibitors, GI upset with antibiotics). Side effects differ from toxicity (excessive dose), allergic reactions, and idiosyncratic reactions. Management includes dose reduction, alternate dosing, adjunctive treatment, or switching drugs.

\medterm{Sleep Medications and Supplements}-- Pharmacological agents used to treat insomnia and other sleep disorders when behavioral interventions alone are insufficient. Prescription sleep aids include benzodiazepine receptor agonists such as zolpidem, melatonin receptor agonists, and low-dose sedating antidepressants, while common over-the-counter options include antihistamines and melatonin supplements. These agents should be used with caution due to risks of dependence, tolerance, daytime drowsiness, and potential interactions with other medications, and are typically recommended for short-term use under medical supervision.

\medterm{Solution}-- A homogeneous mixture in which a drug (solute) is dissolved in a liquid vehicle (solvent, usually water, saline, alcohol, or polyethylene glycol). Solutions provide rapid absorption (oral solutions, parenteral IV solutions), ease of swallowing, and flexible dosing; they require preservatives and have limited stability. Used for oral (syrups, drops), injectable (IV/IM), topical, ophthalmic, otic, and nasal administration. Contrast with suspensions (particles) and emulsions. Isotonicity, pH, and sterility are important for parenteral and ophthalmic solutions.

\medterm{Somatostatin Analog}-- A class of synthetic peptides that mimic the actions of native somatostatin, a hormone that inhibits the secretion of growth hormone, thyroid-stimulating hormone, insulin, glucagon, gastrin, and other GI hormones. Octreotide (short-acting, subcutaneous TID) and lanreotide (long-acting, IM every 4 weeks) are somatostatin analogs with longer half-lives (1-2 hours for octreotide vs. 2-3 minutes for native somatostatin). Indications: acromegaly — suppress GH and IGF-1 levels, reduce tumor size; neuroendocrine tumors (carcinoid syndrome, VIPoma, glucagonoma) — control hormone-related symptoms and may slow tumor growth; variceal bleeding — reduce splanchnic blood flow (octreotide infusion); and refractory diarrhea from chemotherapy or short bowel syndrome. Adverse effects: gallstones (25-50\% due to inhibition of gallbladder contraction), bradycardia, hyperglycemia/hypoglycemia, and GI symptoms.

\medterm{Spironolactone}-- A nonselective aldosterone (mineralocorticoid) receptor antagonist and potassium-sparing diuretic. Used for heart failure with reduced ejection fraction (NYHA II-IV—reduces mortality, often with ACE/ARB + beta-blocker), resistant hypertension, primary hyperaldosteronism, ascites in cirrhosis, and acne/hirsutism in women (off-label). Adverse effects: hyperkalemia, gynecomastia (common, limiting), breast tenderness, menstrual irregularities, and hypotension. Monitor potassium and renal function; avoid with potassium supplements and strong potassium-raising drugs.

\medterm{SSRIs}-- Selective serotonin reuptake inhibitors are a class of antidepressant drugs that selectively inhibit the reuptake of serotonin from the synaptic cleft into the presynaptic neuron by blocking the serotonin transporter, increasing serotonergic neurotransmission. Common SSRIs include fluoxetine, sertraline, paroxetine, citalopram, escitalopram, and fluvoxamine. They are first-line agents for major depressive disorder, generalized anxiety disorder, panic disorder, obsessive-compulsive disorder, posttraumatic stress disorder, premenstrual dysphoric disorder, and bulimia nervosa. Common adverse effects include nausea, insomnia, sexual dysfunction, and weight changes; SSRIs have a favourable safety profile in overdose compared to TCAs and MAOIs, making them the preferred first-line pharmacotherapy for most depressive and anxiety disorders.

\medterm{Steady State}-- The condition where the rate of drug administration equals the rate of elimination, resulting in constant average plasma concentrations with repeated dosing. During regular dosing at fixed intervals, steady state is achieved after approximately 4 to 5 half-lives, regardless of the dosing frequency or dose size. At steady state, the peak and trough concentrations fluctuate within a therapeutic range, with the degree of fluctuation determined by the dosing interval relative to the the drug's half-life; the concept of steady state is fundamental to therapeutic drug monitoring and ensures that drug concentrations remain within the therapeutic window for optimal efficacy and safety.

\medterm{Stenotrophomonas maltophilia resistance}-- Stenotrophomonas maltophilia is a gram-negative bacterium with intrinsic resistance to many antibiotics including carbapenems. The organism produces a chromosomal beta-lactamase (L1) that confers resistance to beta-lactams. Treatment options include trimethoprim-sulfamethoxazole, fluoroquinolones, and minocycline. S. maltophilia is an emerging opportunistic pathogen, particularly in immunocompromised patients and those with cystic fibrosis. The organism forms biofilms and is resistant to many disinfectants, contributing to its persistence in healthcare environments; treatment requires combination therapy based on susceptibility testing, and antimicrobial stewardship is critical for preserving the limited treatment options available for this organism.

\medterm{Stevens-Johnson Syndrome}-- A severe, life-threatening mucocutaneous reaction characterized by widespread skin erythema, necrosis, and bullae formation with mucosal involvement (oral, ocular, genital). See cardiology.

\medterm{Streptokinase}-- A thrombolytic agent that forms a plasminogen activator complex, converting plasminogen to plasmin and lysing fibrin clots. Used for STEMI, massive pulmonary embolism, and deep vein thrombosis (largely replaced by tissue plasminogen activator, which is more fibrin-specific and less antigenic). Given IV; associated with allergic reactions and antibodies after prior exposure (reduced efficacy with repeat use). Bleeding is the major adverse effect; contraindicated in active hemorrhage, recent surgery, and intracranial pathology.

\medterm{Sulfonamide}-- A class of antibiotics (sulfamethoxazole, sulfadiazine, sulfasalazine) that inhibit bacterial dihydropteroate synthase, blocking folate synthesis—bacteriostatic. Often combined with trimethoprim (co-trimoxazole: TMP-SMX). Used for UTIs, PCP (pneumocystis pneumonia), toxoplasmosis, Nocardia, and MRSA. Adverse effects: hypersensitivity reactions (rash, Stevens-Johnson syndrome), crystalluria (hydrate, alkalinize), photosensitivity, hemolysis in G6PD deficiency, and hyperkalemia (TMP). Avoid in sulfa allergy and in pregnancy near term.

\medterm{Sulfonamides}-- A class of bacteriostatic antibiotics that competitively inhibit dihydropteroate synthase, an enzyme essential for folic acid synthesis in bacteria. They are structural analogs of para-aminobenzoic acid and compete with PABA for the active site of the enzyme. Trimethoprim is combined with sulfamethoxazole to produce synergistic bactericidal activity by blocking sequential steps in folate synthesis. Co-trimoxazole, the combination of trimethoprim and sulfamethoxazole, is used forurinary tract infections, Pneumocystis jirovecii pneumonia, and other infections; hypersensitivity reactions to sulfonamides are common, particularly in patients with HIV infection, and range from mild rash to severe cutaneous adverse drug reactions including Stevens-Johnson syndrome.

\medterm{Sulfonylureas}-- Oral antidiabetic agents that stimulate insulin secretion from pancreatic beta cells by blocking ATP-sensitive potassium channels on the beta cell membrane. This causes membrane depolarization, opening of voltage-gated calcium channels, calcium influx, and exocytosis of insulin-containing granules. They are effective at lowering blood glucose but carry a risk of hypoglycemia because they stimulate insulin secretion regardless of glucose levels. First-generation sulfonylureas ininclude tolbutamide and chlorpropamide; second-generation sulfonylureas including glipizide, glyburide, and glimepiride are more potent, have a more favourable side effect profile, and are dosed once or twice daily. Sulfonylureas remain an important treatment option for type 2 diabetes, particularly in resource-limited settings where cost is a significant consideration.

\medterm{Sulphonamides}-- A class of antimicrobial agents containing the sulfonamide group (SO$_2$NH$_2$). Sulphonamides were the first effective systemic antibacterial agents (Prontosil, 1935—Gerhard Domagk). Mechanism: competitive inhibition of dihydropteroate synthase, blocking folate synthesis (para-aminobenzoic acid analogue)—bacteriostatic. Spectrum: broad—Gram-positive (Streptococcus, Staphylococcus), Gram-negative (E. coli, Klebsiella, Haemophilus, Nocardia), and some protozoa (Toxoplasma, Plasmodium). Clinical uses: (1) Sulfamethoxazole combined with trimethoprim (co-trimoxazole)—first-line for Pneumocystis jirovecii pneumonia (PCP), Nocardia, Toxoplasma encephalitis, and UTIs. (2) Sulfadiazine—with pyrimethamine for toxoplasmosis. (3) Sulfasalazine—used in inflammatory bowel disease and rheumatoid arthritis (5-ASA prodrug). (4) Mafenide and silver sulfadiazine—topical for burn wounds. (5) Sulfacetamide—ophthalmic drops for conjunctivitis. Adverse effects: hypersensitivity reactions (rash, Stevens–Johnson syndrome/toxic epidermal necrolysis—rare but life-threatening, fever, serum sickness), crystalluria (prevented by adequate hydration and urine alkalinisation—risk with older, less soluble sulphonamides), hemolytic anemia in G6PD deficiency, bone marrow suppression, photosensitivity, and kernicterus in neonates (displaces bilirubin from albumin—contraindicated in pregnancy, breastfeeding, and infants <2 months). Drug interactions: potentiates warfarin, methotrexate, and sulfonylureas.

\medterm{Suppository}-- A solid, single-dose dosage form designed to be inserted into a body orifice (rectum, vagina, or urethra) where it melts, softens, or dissolves, releasing the active pharmaceutical ingredient for local or systemic effects. Rectal suppositories (the most common): base materials (cocoa butter/theobroma oil, hard fat/Witepsol, glycero-gelatin, polyethylene glycol). Indications: local—hemorrhoids (anusol, hydrocortisone), constipation (glycerol, bisacodyl); systemic—antiemetics (prochlorperazine, promethazine), antipyretics/analgesics (paracetamol, diclofenac), bronchodilators (theophylline), and antiemetics when oral route is unavailable (nausea/vomiting, post-operative). Advantages: bypasses first-pass hepatic metabolism (some drugs achieve higher bioavailability), useful when oral route is unavailable (NPO, vomiting, unconscious patients), and avoids GI irritation. Insertion: supine or lateral position, push past the internal anal sphincter, retain for 15--30 minutes. Vaginal suppositories (pessaries): for local treatment of vaginal infections (clotrimazole, metronidazole, oestrogen) and contraception.

\medterm{Suspension}-- A pharmaceutical formulation in which fine solid drug particles are dispersed (not dissolved) in a liquid vehicle, often with suspending agents and stabilizers. Suspensions are used when a drug is poorly soluble, for oral (antibiotic suspensions, antacids), topical, ophthalmic, and parenteral (depot) administration. They require shaking before use to ensure uniform dosing and may settle over time. Bioavailability depends on particle size. Contrast with solutions (dissolved) and emulsions (liquid-liquid).

\medterm{Sympatholytic}-- An agent that blocks the sympathetic nervous system, used to treat hypertension, sympathetic overactivity, and catecholamine excess. Classes: beta-blockers (propranolol, metoprolol), alpha-blockers (prazosin, doxazosin, phenoxybenzamine), centrally acting agents (clonidine, methyldopa), and neuron blockers (reserpine, guanethidine). Effects: reduced heart rate and contractility, vasodilation, lowered blood pressure, and reduced sympathetic reflexes. Adverse effects: bradycardia, hypotension, fatigue, bronchospasm (beta-blockers in asthma), and rebound hypertension on abrupt withdrawal.

\medterm{Tablet}-- A solid dosage form of compressed drug plus excipients (binders, fillers, disintegrants, lubricants, coatings). Types: immediate-release, enteric-coated (resistant to gastric acid), sustained/controlled-release, sublingual, buccal, effervescent, chewable, and orally disintegrating. Tablets provide dose accuracy, stability, ease of administration, and mass production. They should generally not be crushed (sustained-release/enteric forms) unless directed. Some tablets are scored for splitting.

\medterm{Tedizolid pharmacokinetics}-- Tedizolid has a longer half-life than linezolid, allowing for once-daily dosing. The drug has approximately 90\% oral bioavailability. Tedizolid is metabolized to an active metabolite (tzd-M1) that contributes to efficacy. The drug achieves high concentrations in skin and soft tissues. Tedizolid does not inhibit monoamine oxidase to the same degree as linezolid, reducing drug interaction potential. The drug is excreted primarily in feces.

\medterm{Teicoplanin}-- A glycopeptide antibiotic related to vancomycin with a longer half-life allowing for less frequent dosing. Teicoplanin inhibits cell wall synthesis by binding to D-Ala-D-Ala precursors. It has activity against gram-positive organisms including MRSA. The drug is administered intravenously or intramuscularly, with once-daily or less frequent dosing after loading. Common adverse effects include injection site reactions, rash, and nephrotoxicity. Teicoplanin achieves higher tissue concentrations thathan vancomycin, making it useful for deep-seated infections including bone and joint infections, endocarditis, and prosthetic device infections; teicoplanin is not available in the United States but is widely used in Europe and other regions.

\medterm{Tenecteplase}-- A modified tissue plasminogen activator (TNK-tPA) thrombolytic with longer half-life, higher fibrin specificity, and resistance to PAI-1. Used for acute ischemic stroke and STEMI. Given as a single IV bolus (weight-based), simplifying administration. Adverse effects: hemorrhage (intracranial—major concern), allergic reactions. Contraindications: active bleeding, recent surgery/trauma, uncontrolled hypertension, and intracranial pathology. Time to treatment is critical.

\medterm{Teratogenicity}-- The ability of a drug, chemical, or infectious agent to cause birth defects or developmental abnormalities when exposure occurs during pregnancy. The risk of teratogenic effects depends on the timing of exposure, with the embryonic period from 3 to 8 weeks after fertilization being the most critical for organogenesis. Certain drugs have well-established teratogenic effects: thalidomide causes limb defects, isotretinoin causes craniofacial, cardiac, and CNS abnormalities,and sodium valproate is associated with neural tube defects; the management of drug therapy during pregnancy requires careful risk-benefit assessment, minimising exposure during organogenesis whenever possible, and using alternative agents with established safety profiles for the treatment of chronic conditions in women of childbearing potential.

\medterm{Terazosin}-- A selective alpha-1 adrenergic blocker causing vascular and prostatic smooth muscle relaxation. Used for benign prostatic hyperplasia (improves urinary flow) and hypertension. Adverse effects: first-dose orthostatic hypotension/syncope (start low, at bedtime), dizziness, nasal congestion, and fatigue. Not first-line for hypertension but useful in men with BPH. Doxazosin and prazosin are related alpha-blockers.

\medterm{Terbinafine}-- An allylamine antifungal agent that inhibits squalene epoxidase, an early enzyme in the ergosterol biosynthesis pathway, causing accumulation of squalene and depletion of ergosterol in the fungal cell membrane. The accumulation of squalene is fungicidal, while ergosterol depletion disrupts membrane function. Terbinafine has excellent activity against dermatophytes including Trichophyton, Microsporum, and Epidermophyton species, making it the treatment of choice for dermatophyte ininfections of the skin and nails. Unlike azole antifungals, terbinafine does not significantly inhibit CYP450 enzymes and has a lower potential for drug interactions; adverse effects include mild gastrointestinal disturbances, headache, and taste disturbance, with rare cases of hepatotoxicity requiring periodic liver function monitoring during prolonged therapy.

\medterm{Tetracyclines}-- A class of bacteriostatic antibiotics that bind reversibly to the 30S ribosomal subunit and block the attachment of aminoacyl-tRNA to the acceptor site, inhibiting protein synthesis. They include tetracycline, doxycycline, and minocycline. Tetracyclines have a broad spectrum of activity against gram-positive and gram-negative bacteria, atypical organisms including Chlamydia, Mycoplasma, and Rickettsia, and the malaria parasite. Doxycycline is the most commonly used tetracyclinefor clinical use, with excellent oral bioavailability, favourable pharmacokinetics allowing once- or twice-daily dosing, and a well-established safety profile; tetracyclines also have anti-inflammatory and immunomodulatory properties that are independent of their antimicrobial activity and contribute to their efficacy in acne vulgaris and rosacea.

\medterm{Therapeutic Drug Monitoring}-- The clinical practice of measuring specific drug concentrations in blood or plasma to individualize drug dosing and ensure therapeutic effectiveness while minimizing toxicity. TDM is particularly important for drugs with narrow therapeutic indices, where small changes in concentration can lead to therapeutic failure or toxicity. Examples include aminoglycoside antibiotics where both subtherapeutic and supratherapeutic concentrations can occur, vancomycin with target trough concentrations for vancomycin requiring trough monitoring in the range of 10-20 mcg/mL depending on the infection type and severity; other drugs requiring TDM include digoxin, lithium, theophylline, cyclosporine, tacrolimus, and certain anticonvulsants including phenytoin and valproic acid.

\medterm{Therapeutic Index}-- The therapeutic index is a quantitative measure of the relative safety of a drug, defined as the ratio of the dose that produces toxicity in 50 percent of the population to the dose that produces the desired therapeutic effect in 50 percent of the population. A wider therapeutic index indicates a larger margin of safety, meaning there is a greater difference between effective and toxic doses. Drugs with narrow therapeutic indices require careful dosing and monitoring because small changes in doses in drug concentration produce significant clinical consequences; drugs with narrow therapeutic indices require individualised dosing based on patient factors, therapeutic drug monitoring, and careful management of drug interactions to maintain concentrations within the therapeutic window and avoid toxicity.

\medterm{Therapeutics}-- The branch of medicine concerned with the treatment of disease, encompassing all interventions used to prevent, manage, or cure illness. Therapeutics includes: pharmacotherapy (drug treatment—most common), surgery (therapeutic procedures), radiotherapy (ionizing radiation for cancer), physical therapy (physiotherapy, rehabilitation), psychological therapy (CBT, psychotherapy), gene and cell therapy, immunotherapy, and complementary/alternative medicine. Principles: (1) Evidence-based medicine: treatment decisions based on the best available evidence (randomised controlled trials, systematic reviews, meta-analyses, clinical guidelines). (2) Individualised therapy: patient-specific factors (age, genetics, comorbidities, concurrent medications, preferences) guide treatment choice. (3) Risk–benefit assessment: consider efficacy, safety, tolerability, cost, and quality of life. (4) Therapeutic index: the ratio of toxic to therapeutic dose; drugs with a narrow therapeutic index (warfarin, lithium, digoxin, phenytoin) require therapeutic drug monitoring. (5) Adherence/compliance: the extent to which a patient follows prescribed treatment; poor adherence (50\% of chronic disease patients) is a major cause of treatment failure. (6) Polypharmacy: use of multiple medications (\(\ge\)5) increases risk of adverse drug reactions, drug interactions, and non-adherence. (7) Deprescribing: systematic withdrawal of inappropriate or unnecessary medications.

\medterm{Therapeutic Window}-- The range of plasma drug concentrations between the minimum effective concentration (producing a therapeutic effect) and the minimum toxic concentration (producing toxicity), in which efficacy is achieved with acceptable toxicity. Drugs with a narrow therapeutic window (digoxin, warfarin, lithium, aminoglycosides, cyclosporine) require therapeutic drug monitoring and careful dosing because small changes in concentration cause therapeutic failure or toxicity. The width of the window relates to the therapeutic index (TD50/ED50).

\medterm{Thiazide Diuretics}-- Thiazide diuretics act on the distal convoluted tubule by inhibiting the sodium-chloride cotransporter, reducing sodium and chloride reabsorption and producing natriuresis and diuresis. They are less potent than loop diuretics but have a longer duration of action, typically 6 to 12 hours. Common thiazides include hydrochlorothiazide, chlorthalidone, and metolazone. Thiazides are first-line antihypertensive agents, particularly effective in elderly patients with isolated systolic hypertension andand in patients of African descent; adverse effects include hypokalaemia, hyperuricaemia, hyperglycaemia, hypercalcaemia, and metabolic alkalosis, and they should be used with caution in patients with diabetes mellitus, gout, and electrolyte imbalances.

\medterm{Thiopental Sodium}-- An ultra-short-acting barbiturate, used for induction of general anesthesia (IV, onset 30--60 seconds, duration 5--10 minutes) and as a third-line agent for raised intracranial pressure (cerebroprotective effect: reduces cerebral metabolism, cerebral blood flow, and intracranial pressure). Mechanism: enhances GABA$_A$ receptor activity (prolongs chloride channel opening). Side effects: dose-dependent respiratory depression, hypotension (decreased cardiac output, peripheral vasodilation), laryngospasm, and tissue necrosis/extravasation injury (highly alkaline, pH 10.5—administer through a secure IV line). Contraindications: porphyria (acute intermittent porphyria—induces haem oxygenase, precipitates attack), severe shock, status asthmaticus, and hypersensitivity to barbiturates. Thiopental was also used in lethal injection protocols (US) and for barbiturate coma (refractory intracranial hypertension). Now largely replaced by propofol (faster recovery, fewer side effects) for routine anesthetic induction. Historical: 'truth serum' myth—thiopental reduces inhibition but does not produce truthful statements reliably.

\medterm{Thrombolytic}-- A drug that dissolves blood clots by activating plasminogen to plasmin, which degrades fibrin. Agents: alteplase (tPA), tenecteplase (TNK-tPA), reteplase, streptokinase, urokinase. Indications: acute ischemic stroke (within 4.5 hours), STEMI, massive pulmonary embolism, and acute limb ischemia. Administration is IV; the major risk is hemorrhage (especially intracranial). Contraindications include active bleeding, recent surgery/trauma, and uncontrolled hypertension. Time to treatment is critical for benefit.

\medterm{Tigecycline}-- A glycylcycline antibiotic derived from minocycline with expanded spectrum of activity. Tigecycline inhibits bacterial protein synthesis by binding to the 30S ribosomal subunit. It has activity against gram-positive organisms including MRSA and VRE, gram-negative organisms including ESBL-producing Enterobacteriaceae, and anaerobes. Tigecycline is indicated for complicated intra-abdominal infections, complicated skin infections, and community-acquired pneumonia. The drug is administered intravenously twice daily, with common adverse effects including nausea, vomiting, abdominal pain, and hepatotoxicity; tigecycline carries a FDA black box warning for increased all-cause mortality compared to comparator agents, particularly in patients treated for hospital-acquired pneumonia, and should be reserved for situations where alternative treatment options are limited or inadequate.

\medterm{Tolbutamide}-- A first-generation sulfonylurea oral hypoglycaemic agent used in the management of type 2 diabetes. Mechanism: stimulates insulin secretion from pancreatic $\beta$-cells by binding to the sulfonylurea receptor (SUR1) on ATP-sensitive potassium ($\mathrm{K}_{\mathrm{ATP}}$) channels, causing membrane depolarisation, calcium influx, and exocytosis of insulin granules. Tolbutamide has a short half-life (4--7 hours) and is metabolized by the liver to inactive metabolites (CYP2C9). It is less potent and shorter-acting than second-generation sulfonylureas (glibenclamide/glyburide, gliclazide, glipizide, glimepiride). Clinical use: now rarely used (largely replaced by newer sulfonylureas and other oral agents). Indication: adjunct to diet and exercise in type 2 diabetes when metformin is contraindicated or not tolerated. Side effects: hypoglycaemia (less frequent than with longer-acting sulfonylureas), weight gain, and hypersensitivity reactions (rash, photosensitivity). The University Group Diabetes Program (UGDP) trial (1970s) suggested an increased cardiovascular mortality with tolbutamide, though this finding has been disputed. Drug interactions: potentiates effect of alcohol (disulfiram-like reaction), warfarin, and NSAIDs; diminished effect with thiazide diuretics, corticosteroids, and rifampicin.

\medterm{Tolerance}-- A decreased response to a drug after repeated administration, requiring dose escalation to achieve the same effect. Unlike tachyphylaxis, tolerance develops over weeks to months of chronic use. The mechanisms include pharmacokinetic tolerance from increased drug metabolism through enzyme induction, pharmacodynamic tolerance from receptor downregulation and desensitization, and learned or behavioral tolerance from compensatory physiological responses. Opioid tolerance is a major clinical challenge in chronic pain management, leading to the need for gradual dose escalation, careful assessment of risk versus benefit, and the use of multidisciplinary approaches including non-opioid analgesics, interventional procedures, and behavioural therapies.

\medterm{Torsemide}-- A potent loop diuretic inhibiting the Na-K-2Cl cotransporter in the loop of Henle. Used for edema in heart failure, cirrhosis, and renal disease, and hypertension. More consistently absorbed than furosemide; long duration; oral-to-IV conversion is nearly 1:1. May have less mortality effect than furosemide in heart failure trials. Adverse effects: hypokalemia, hyponatremia, hypomagnesemia, hypocalcemia, dehydration, ototoxicity, and hyperuricemia. Half-life ~3-4 hours (longer in renal/hepatic failure).

\medterm{TPMT Polymorphism}-- Thiopurine S-methyltransferase is an enzyme that metabolizes thiopurine drugs including azathioprine, 6-mercaptopurine, and thioguanine, which are used in the treatment of leukemia, inflammatory bowel disease, and autoimmune conditions. TPMT catalyzes the S-methylation of thiopurines, diverting them away from the formation of active thioguanine nucleotides that cause cytotoxicity. Genetic polymorphisms in the TPMT gene result in variable enzyme activity, with approximately 10 percent of the popuopulation having intermediate activity and approximately 0.3\% having low or absent activity; patients with intermediate or low TPMT activity are at increased risk of severe myelosuppression when treated with standard doses of thiopurines, and pre-treatment TPMT genotyping or phenotyping is recommended to guide initial dosing and minimise the risk of haematological toxicity.

\medterm{Tramadol}-- An atypical opioid analgesic (weak mu-agonist plus serotonin and norepinephrine reuptake inhibition) used for moderate pain. Prodrug (M1, via CYP2D6) contributes analgesia; efficacy varies with CYP2D6 genotype. Adverse effects: nausea, dizziness, sedation, constipation, seizures (at high doses, in epilepsy, with SSRIs/SNRIs/TCAs), serotonin syndrome, and dependence (Schedule IV). Caution in renal/hepatic impairment; avoid in opioid-naive children post-tonsillectomy.

\medterm{Tranquillisers}-- A class of drugs that have a calming or sedative effect, reducing anxiety, tension, and agitation. Historically classified as major tranquillisers (antipsychotics/neuroleptics—for psychosis) and minor tranquillisers (anxiolytics—for anxiety). The term is now less common in clinical practice. Major tranquillisers (antipsychotics): first-generation/typical (chlorpromazine, haloperidol, fluphenazine—block D$_2$ receptors), second-generation/atypical (risperidone, olanzapine, quetiapine, aripiprazole, clozapine—block D$_2$ and 5-HT\(_{2A}\) receptors). Used for schizophrenia, bipolar mania, agitation, and treatment-resistant depression. Minor tranquillisers (anxiolytics): benzodiazepines (diazepam, lorazepam, alprazolam—enhance GABA$_A$ activity), buspirone (5-HT\(_{1A}\) partial agonist—non-sedating, no dependence), and $\beta$-blockers (propranolol—for performance anxiety, somatic symptoms). Adverse effects: (1) Benzodiazepines—dependence, tolerance, withdrawal (rebound anxiety, insomnia, seizures), sedation, cognitive impairment, falls. (2) Antipsychotics—extrapyramidal side effects (dystonia, parkinsonism, akathisia, tardive dyskinesia), weight gain, metabolic syndrome, hyperprolactinaemia, QT prolongation. Non-pharmacological alternatives: cognitive behavioural therapy, relaxation techniques, and mindfulness.

\medterm{Transdermal}-- A route of drug administration delivering drug through intact skin into the systemic circulation via a patch or gel, bypassing first-pass metabolism and providing steady plasma concentrations over hours to days. Examples: nitroglycerin, fentanyl, nicotine, estrogen, clonidine, rotigotine. Absorption depends on drug lipophilicity, molecular size, skin site, and occlusion. Advantages: sustained release, improved adherence, avoidance of GI irritation. Risks: skin irritation, dosing limitations (potent, lipophilic drugs only), and accidental exposure (fentanyl patch).

\medterm{Trimethoprim-sulfamethoxazole}-- A fixed-dose combination of trimethoprim and sulfamethoxazole that synergistically inhibits folate synthesis. TMP-SMX is indicated for urinary tract infections, Pneumocystis jirovecii pneumonia, and community-acquired pneumonia. The combination is bactericidal against many gram-positive and gram-negative organisms. Common adverse effects include hypersensitivity reactions, bone marrow suppression, and hyperkalemia. The drug is contraindicated in patients with folate deficiency and in late pregnanancy; resistance is increasingly common among many bacterial species, particularly E. coli, and susceptibility testing is essential before prescribing TMP-SMX for urinary tract infections.

\medterm{Triptan}-- A class of medications used for acute migraine treatment, acting as 5-HT1B/1D receptor agonists. 5-HT1B receptors are located on cranial blood vessels; activation causes vasoconstriction of dilated meningeal arteries. 5-HT1D receptors are located on presynaptic trigeminal nerve terminals; activation inhibits the release of vasoactive neuropeptides (calcitonin gene-related peptide, substance P), reducing neurogenic inflammation and pain transmission. Available triptans: sumatriptan (first, available as oral, subcutaneous, intranasal, and rectal), rizatriptan, zolmitriptan, naratriptan, eletriptan, almotriptan, and frovatriptan. Contraindications: ischemic heart disease, Prinzmetal angina, uncontrolled hypertension, hemiplegic migraine, and basilar migraine. Adverse effects: chest tightness/pressure (may mimic cardiac ischemia but typically benign), paresthesias, dizziness, and injection site reactions. Triptans should be taken early in the migraine attack for optimal efficacy.

\medterm{Tyrosine Kinase Inhibitor}-- A class of targeted cancer therapy drugs that inhibit tyrosine kinases, enzymes that phosphorylate tyrosine residues on proteins and activate signaling pathways involved in cell growth, proliferation, differentiation, and survival. Tyrosine kinase mutations or overexpression drive many cancers. TKIs are small molecules (oral) or monoclonal antibodies (IV). Examples: imatinib — BCR-ABL inhibitor for chronic myeloid leukemia (first TKI approved, revolutionized CML treatment); erlotinib and gefitinib — EGFR inhibitors for non-small cell lung cancer with EGFR mutations; sorafenib and sunitinib — multi-kinase inhibitors (VEGFR, PDGFR, RAF) for renal cell carcinoma and hepatocellular carcinoma; osimertinib — third-generation EGFR TKI for T790M mutant NSCLC. Adverse effects vary by target: hypertension (VEGF inhibitors), rash and diarrhea (EGFR inhibitors), fatigue, and hand-foot skin reaction.

\medterm{Unguentum}-- A semi-solid, greasy topical preparation for application to the skin, mucous membranes, or eyes. See Ointment.

\medterm{Valproate (Sodium Valproate)}-- Valproate (valproic acid, sodium valproate, divalproex sodium) is a broad-spectrum antiepileptic drug effective against all seizure types (generalised tonic-clonic, absence, myoclonic, and partial seizures). It is also used as a mood stabiliser in bipolar disorder (particularly for acute mania and maintenance) and for prophylaxis of migraine headaches. Mechanism: multiple actions including enhancement of GABAergic transmission (by increasing GABA synthesis via glutamic acid decarboxylase activation and inhibiting GABA transaminase), blockade of voltage-gated sodium channels, and inhibition of T-type calcium channels. Pharmacokinetics: well absorbed orally, highly protein-bound (90\%), extensively metabolized in the liver (glucuronidation, beta-oxidation, and CYP-mediated oxidation), half-life 9-16 hours. Adverse effects: dose-related (tremor, weight gain, alopecia, sedation, thrombocytopenia, elevated liver enzymes), idiosyncratic (hepatotoxicity — highest risk in children under 2 years on polytherapy, pancreatitis), teratogenic (neural tube defects, spina bifida — absolute contraindication in pregnancy unless no alternative, requiring high-dose folic acid supplementation), and polycystic ovary syndrome with long-term use. Monitoring: serum levels (therapeutic range 50-100 mcg/mL), LFTs, CBC, and pregnancy testing. Drug interactions: inhibits CYP2C9, displaces protein-bound drugs (phenytoin).

\medterm{Vancomycin-resistant Enterococcus}-- Enterococci resistant to vancomycin through acquisition of vanA or vanB gene clusters encoding altered cell wall precursors. VRE is a significant cause of healthcare-associated infections, particularly in immunocompromised patients. Treatment options include linezolid, daptomycin, and quinupristin-dalfopristin (for E. faecium only). VRE infections are associated with increased mortality and healthcare costs. Infection prevention strategies include hand hygiene, contact precautions, and environmeonmental decontamination, and antimicrobial stewardship to preserve the effectiveness of remaining active agents.

\medterm{Vaughan Williams Classification}-- The Vaughan Williams classification categorizes antiarrhythmic drugs into four classes based on their primary mechanism of action. Class I agents block sodium channels and slow conduction velocity, further subdivided into IA with moderate sodium channel blockade and action potential prolongation including quinidine and procainamide, IB with fast sodium channel binding and action potential shortening including lidocaine and mexiletine, and IC with slow sodium channel binding and minimal effect onon action potential duration; class II agents are beta-blockers that reduce sympathetic activity and slow AV conduction; class III agents block potassium channels, prolonging repolarisation; and class IV agents are calcium channel blockers that slow AV conduction; while clinically useful, the classification has limitations as many drugs have multiple mechanisms of action and it does not account for all clinically important antiarrhythmic drugs.

\medterm{Viagra (Sildenafil)}-- Sildenafil is a phosphodiesterase type 5 (PDE5) inhibitor used primarily for erectile dysfunction (ED) and pulmonary arterial hypertension (PAH, as Revatio). Mechanism: PDE5 degrades cGMP in penile corpus cavernosum smooth muscle; sildenafil inhibits PDE5, increasing cGMP levels and enhancing nitric oxide (NO)-mediated vasodilation and smooth muscle relaxation, improving penile blood flow during sexual stimulation. For PAH: sildenafil dilates pulmonary vasculature, reducing pulmonary vascular resistance. Pharmacokinetics: oral, onset within 30-60 minutes, duration 4-6 hours, metabolized by CYP3A4. Adverse effects: headache, flushing, dyspepsia, nasal congestion, visual disturbances (blue tinged vision/cyanopsia due to PDE6 inhibition in retinal photoreceptors), and myalgia. Serious: priapism (erection >4 hours — medical emergency), hypotension (especially when combined with nitrates — absolute contraindication due to risk of severe hypotension and syncope), and non-arteritic anterior ischemic optic neuropathy (NAION, rare). Contraindications: concurrent nitrate therapy, severe hepatic impairment, severe hypotension, and recent stroke or MI.

\medterm{Vibrio vulnificus infections}-- Infections caused by Vibrio vulnificus, a gram-negative bacterium found in warm seawater and raw seafood. V. vulnificus can cause wound infections, septicemia, and gastroenteritis. The most severe form is necrotizing fasciitis and primary septicemia, which can be fatal. Risk factors include liver disease, immunocompromise, and iron overload states. Treatment requires doxycycline combined with a third-generation cephalosporin. Early wound care and antibiotic therapy are essential. The organism hahas a high case-fatality rate, and prevention involves avoiding consumption of raw or undercooked seafood, particularly in patients with liver disease or immunocompromising conditions.

\medterm{Volume of Distribution}-- A theoretical pharmacokinetic parameter that relates the total amount of drug in the body to the plasma concentration, calculated as Vd equals amount of drug in the body divided by plasma concentration. It represents the apparent volume into which a drug distributes and is not a true anatomical volume. A small volume of distribution, such as 3 to 5 liters, indicates that the drug remains primarily in the plasma and does not extensively distribute into tissues, as seen with drugs that are highly protein-bound and extensively distributed to tissues; understanding volume of distribution is critical for calculating loading doses, interpreting drug concentrations, and predicting the extent of drug distribution to the site of action.

\medterm{Voriconazole}-- A triazole antifungal agent used in the treatment of invasive aspergillosis, candidemia, and other serious fungal infections. Voriconazole inhibits the fungal cytochrome P450 enzyme lanosterol 14-alpha demethylase, preventing ergosterol synthesis and disrupting fungal cell membrane integrity. The drug has broad-spectrum activity against Aspergillus, Candida, Cryptococcus, and other fungi. A key advantage of voriconazole over itraconazole is its excellent oral bioavailability of approximately 96\%and has a long half-life allowing for twice-daily oral dosing after an initial intravenous loading regimen; voriconazole is associated with visual disturbances (photopsia, altered color perception), hepatotoxicity, and QT prolongation, and requires therapeutic drug monitoring due to non-linear pharmacokinetics and significant inter-individual variability in plasma concentrations; drug interactions are numerous due to CYP2C19 and CYP3A4 inhibition.

\medterm{Weight Loss Medications}-- Pharmacologic agents approved for the treatment of obesity when used in conjunction with lifestyle modifications. These medications work through various mechanisms, including appetite suppression through GLP-1 receptor agonists such as semaglutide, inhibition of fat absorption with orlistat, or combination therapies targeting multiple neuropeptide pathways. They are indicated for individuals with a body mass index of 30 or higher, or 27 or higher with obesity-related comorbidities, and require ongoing medical supervision due to potential side effects and the need for monitoring.

\medterm{Withdrawal Symptoms}-- See Withdrawal Syndrome. used primarily for erectile dysfunction (ED) and pulmonary arterial hypertension (PAH, as Revatio). Mechanism: PDE5 degrades cGMP in penile corpus cavernosum smooth muscle; sildenafil inhibits PDE5, increasing cGMP levels and enhancing nitric oxide (NO)-mediated vasodilation and smooth muscle relaxation, improving penile blood flow during sexual stimulation. For PAH: sildenafil dilates pulmonary vasculature, reducing pulmonary vascular resistance. Pharmacokinetics: oral, onset within 30-60 minutes, duration 4-6 hours, metabolized by CYP3A4. Adverse effects: headache, flushing, dyspepsia, nasal congestion, visual disturbances (blue tinged vision/cyanopsia due to PDE6 inhibition in retinal photoreceptors), and myalgia. Serious: priapism (erection >4 hours — medical emergency), hypotension (especially when combined with nitrates — absolute contraindication due to risk of severe hypotension and syncope), and non-arteritic anterior ischemic optic neuropathy (NAION, rare). Contraindications: concurrent nitrate therapy, severe hepatic impairment, severe hypotension, and recent stroke or MI.

\medterm{Withdrawal Syndrome}-- A predictable set of signs and symptoms that occur when a drug is discontinued or the dose is reduced after chronic use, reflecting the physiological rebound from drug adaptation. Specific withdrawal syndromes: Opioid withdrawal (not life-threatening but intensely uncomfortable): yawning, lacrimation, rhinorrhea, piloerection (cold turkey), mydriasis, abdominal cramps, diarrhea, nausea, vomiting, muscle aches, and drug craving. Alcohol withdrawal (potentially life-threatening): tremor, anxiety, autonomic hyperactivity (tachycardia, hypertension, diaphoresis), insomnia, seizures (6-48 hours after last drink), and delirium tremens (48-96 hours later: agitation, hallucinations, autonomic instability, and altered mental status). Benzodiazepine withdrawal: anxiety, insomnia, tremor, palpitations, and seizures (may be delayed with long-acting agents). SSRI discontinuation syndrome: dizziness, paresthesias, nausea, headache, and irritability (not life-threatening). Management involves gradual tapering (detoxification), supportive care, and pharmacotherapy (benzodiazepines for alcohol withdrawal, methadone/buprenorphine for opioid withdrawal).

\medterm{Xanthine Oxidase Inhibitor}-- A class of medications that inhibit xanthine oxidase, the enzyme that catalyzes the final steps of purine metabolism: hypoxanthine to xanthine and xanthine to uric acid. By inhibiting this enzyme, XOIs reduce uric acid production and lower serum uric acid levels, used for management of chronic gout and tumor lysis syndrome. Allopurinol (first-line, 40+ years of clinical use): a purine analog that acts as a competitive substrate and is metabolized to oxypurinol (active metabolite with a long half-life of 18-30 hours). Dosing starts at 100 mg daily, titrated to target serum urate <6 mg/dL. Febuxostat: a non-purine selective XOI, metabolized in the liver, dose 40-80 mg daily, preferred in patients with allopurinol hypersensitivity or renal impairment. Cardiotoxicity warning: febuxostat has a higher risk of cardiovascular death compared to allopurinol in patients with preexisting CVD. Adverse effects: allopurinol hypersensitivity syndrome (DRESS, Stevens-Johnson syndrome, most severe, risk increased in HLA-B*5801 carriers, especially Han Chinese and African Americans), rash, and GI upset.

\medterm{Zalcitabine}-- A nucleoside reverse transcriptase inhibitor (NRTI) that was previously used in combination antiretroviral therapy for HIV-1 infection but is now rarely used due to the availability of more effective and better-tolerated agents. Zalcitabine is a cytidine analogue that inhibits viral reverse transcriptase by competing with deoxycytidine triphosphate and causing DNA chain termination. The drug has a relatively short half-life of approximately 1-2 hours, requiring dosing three times daily. The primprimary adverse effect was peripheral neuropathy, which occurred in a significant proportion of patients and was the dose-limiting toxicity; pancreatitis and lactic acidosis were also reported. Zalcitabine was withdrawn from the market in 2006 due to limited clinical use and the availability of safer, more effective NRTIs.

\medterm{Zidovudine (AZT)}-- Zidovudine (azidothymidine, AZT, ZDV) is a nucleoside analogue reverse transcriptase inhibitor (NRTI), the first antiretroviral drug approved for the treatment of HIV infection (FDA approval 1987). Mechanism: zidovudine is a thymidine analogue that is phosphorylated intracellularly to zidovudine triphosphate by cellular kinases. Zidovudine triphosphate competes with endogenous thymidine triphosphate for incorporation into the growing viral DNA chain by HIV-1 reverse transcriptase. Because zidovudine lacks the 3'-hydroxyl group needed for further chain elongation, its incorporation results in premature termination of the viral DNA chain. Zidovudine has a higher affinity for HIV-1 RT than for human DNA polymerases (particularly DNA polymerase gamma, the mitochondrial enzyme), which accounts for its relative selectivity for viral inhibition. The historical landmark ACTG 019 trial (1989) showed that zidovudine monotherapy delayed disease progression in asymptomatic HIV-infected patients. Zidovudine was the first drug shown to reduce maternal-fetal HIV transmission (ACTG 076, 1994: zidovudine given antepartum and intrapartum to the mother and then for 6 weeks to the neonate reduced transmission by 67\%). Combination therapy (HAART, highly active antiretroviral therapy) replaced monotherapy in 1996 due to rapid development of resistance. Current use: zidovudine is still used in fixed-dose combinations (Combivir = zidovudine + lamivudine) and in resource-limited settings, and for prevention of mother-to-child transmission. Resistance: mutations at codons 41, 67, 70, 210, 215, and 219 of HIV-1 RT (thymidine analogue mutations, TAMs) confer zidovudine resistance. TAMs are selected by zidovudine and stavudine and confer cross-resistance to all NRTIs. Adverse effects: bone marrow suppression (macrocytic anemia, neutropenia — the dose-limiting toxicity, more common with advanced HIV, monitored by serial CBC), myopathy (mitochondrial toxicity due to inhibition of DNA polymerase gamma, causing skeletal muscle weakness and wasting, elevated CK, myalgia, and histologic findings of ragged red fibers on muscle biopsy), lactic acidosis and hepatic steatosis (mitonuchondrial toxicity, hepatic steatosis and severe hepatomegaly with lactic acidosis, a rare but potentially fatal complication), headache, nausea, vomiting, and skin rash. Zidovudine should be avoided in patients with low baseline hemoglobin (<9 g/dL) or neutrophil count (<750 cells/microlitre)., the enzyme that catalyzes the final steps of purine metabolism: hypoxanthine to xanthine and xanthine to uric acid. By inhibiting this enzyme, XOIs reduce uric acid production and lower serum uric acid levels, used for management of chronic gout and tumor lysis syndrome. Allopurinol (first-line, 40+ years of clinical use): a purine analog that acts as a competitive substrate and is metabolized to oxypurinol (active metabolite with a long half-life of 18-30 hours). Dosing starts at 100 mg daily, titrated to target serum urate <6 mg/dL. Febuxostat: a non-purine selective XOI, metabolized in the liver, dose 40-80 mg daily, preferred in patients with allopurinol hypersensitivity or renal impairment. Cardiotoxicity warning: febuxostat has a higher risk of cardiovascular death compared to allopurinol in patients with preexisting CVD. Adverse effects: allopurinol hypersensitivity syndrome (DRESS, Stevens-Johnson syndrome, most severe, risk increased in HLA-B*5801 carriers, especially Han Chinese and African Americans), rash, and GI upset.

